\chapter{Solutions to the Exercises}\label{app:solutions}



\section*{Chapter~\ref{ch:predictive}: Causal Questions}

\paragraph{Exercise~\ref{ex:ch01:1}.} Under complete randomization with $n_1$ treated and $n_0$ controls, each unit is treated with probability $n_1/n$. The treated sample mean $\bar Y_1 = n_1^{-1}\sum_i T_i Y_i(1)$ has expectation $n_1^{-1}\sum_i (n_1/n)Y_i(1) = \bar Y(1)$, the finite population mean of treated potential outcomes, and similarly for controls, so $\bar Y_1 - \bar Y_0$ is unbiased for $\tau_{\mathrm{fp}}$. For the variance write $\bar Y_1 - \bar Y_0 = \sum_i T_ic_i - \sum_i Y_i(0)/n_0$ with $c_i = Y_i(1)/n_1 + Y_i(0)/n_0$. Under complete randomization $\mathrm{Var}(T_i) = n_1n_0/n^2$ and $\mathrm{Cov}(T_i, T_j) = -n_1n_0/\{n^2(n-1)\}$ for $i \neq j$, so $\mathrm{Var}(\sum_i T_ic_i) = (n_1n_0/n)S_c^2$ with $S_c^2 = \sum_i(c_i - \bar c)^2/(n-1)$. Expanding $S_c^2$ and using $S_\tau^2 = S_1^2 + S_0^2 - 2S_{10}$ gives
\[
\mathrm{Var}(\bar Y_1 - \bar Y_0) = \frac{S_1^2}{n_1} + \frac{S_0^2}{n_0} - \frac{S_{\tau}^2}{n},
\]
where $S_t^2$ is the finite population variance of $Y_i(t)$, $S_{10}$ the finite population covariance of $Y_i(1)$ and $Y_i(0)$, and $S_\tau^2$ the variance of the unit effects, all with divisor $n - 1$. The last term is not estimable, because it involves both potential outcomes of the same units. The treated units are a simple random sample of the $n$ units, so the sample variance $s_1^2$ of the observed treated outcomes, with divisor $n_1 - 1$, has expectation $S_1^2$, and likewise $s_0^2$ has expectation $S_0^2$. Neyman's estimator $s_1^2/n_1 + s_0^2/n_0$ therefore has expectation $\mathrm{Var}(\bar Y_1 - \bar Y_0) + S_\tau^2/n$: it is conservative, and exact when every unit has the same effect.

\paragraph{Exercise~\ref{ex:ch01:2}.} Write the observed outcome as $Y = \alpha + \tau T + \beta^\top X + \gamma^\top X T + \varepsilon$, with coefficient vector $\vartheta = (\alpha, \tau, \beta, \gamma)$ of dimension $d = 2p + 2$, a flat prior on $\vartheta$ and Jeffreys' prior on $\sigma^2$. Let $Z$ be the design matrix, $\hat\vartheta$ the least squares estimate and $s^2$ the residual variance on $\nu = n - d$ degrees of freedom. The posterior of $\vartheta$ is multivariate Student $t$ with location $\hat\vartheta$ and scale matrix $s^2(Z^\top Z)^{-1}$. Since $\tau(x) = \tau + \gamma^\top x$ is linear in $x$, the average effect for the covariate law $Q$ is
\[
\tau_Q = \int \tau(x)\,Q(dx) = \tau + \gamma^\top m_Q = a^\top\vartheta,
\]
with $m_Q = \int x\,Q(dx)$ and $a = (0, 1, 0_p, m_Q)$, a fixed linear combination of the coefficients. Its posterior is a scaled Student $t$ with $\nu$ degrees of freedom, location $\hat\tau + \hat\gamma^\top m_Q$ and squared scale $s^2 a^\top (Z^\top Z)^{-1} a$, and for $\nu > 2$ its variance is
\[
\frac{\nu s^2}{\nu - 2}\bigl(B_{\tau\tau} + 2\,m_Q^\top B_{\gamma\tau} + m_Q^\top B_{\gamma\gamma}\,m_Q\bigr),
\]
where $B$ is the block of $(Z^\top Z)^{-1}$ belonging to $(\tau, \gamma)$. Of $Q$ only $m_Q$ appears. Two populations with the same mean of $X$ and very different spreads have the same posterior for the average effect, and without the interaction even $m_Q$ drops out, as Section~\ref{sec:firstexample} says.

\paragraph{Exercise~\ref{ex:ch01:3}.} Given $\theta$, the missing outcome of unit $i$ is normal with variance $\sigma_{1-T_i}^2(1-\rho^2)$, so $\mathrm{Var}(\tau_{\mathrm{fp}}\mid\theta,\text{data})=n^{-2}\sum_i\sigma_{1-T_i}^2(1-\rho^2)$ decreases with $|\rho|$ and vanishes at $\rho=\pm1$. In the two arm model without covariates, with known variances and flat priors, $\mu_t\mid\text{data}\sim N(\bar y_t,\sigma_t^2/n_t)$ independently, giving $\mathrm{Var}(\tau\mid\text{data})=\sigma_1^2/n_1+\sigma_0^2/n_0$. The imputed outcome of a treated unit with observed value $y_i$ has mean $\mu_0+\rho(\sigma_0/\sigma_1)(y_i-\mu_1)$, and that of a control has mean $\mu_1+\rho(\sigma_1/\sigma_0)(y_i-\mu_0)$. The conditional mean of $n\tau_{\mathrm{fp}}$ is therefore $(n_0+\rho n_1\sigma_0/\sigma_1)\mu_1-(n_1+\rho n_0\sigma_1/\sigma_0)\mu_0$ plus data terms. Total variance gives
\[
\begin{split}
n^2\mathrm{Var}(\tau_{\mathrm{fp}}\mid\text{data})={}&(n_0+\rho n_1\sigma_0/\sigma_1)^2\frac{\sigma_1^2}{n_1}+(n_1+\rho n_0\sigma_1/\sigma_0)^2\frac{\sigma_0^2}{n_0}\\
&+(1-\rho^2)(n_1\sigma_0^2+n_0\sigma_1^2).
\end{split}
\]
The $\rho^2$ terms cancel, leaving
\[
\mathrm{Var}(\tau_{\mathrm{fp}}\mid\text{data})=\mathrm{Var}(\tau\mid\text{data})-\frac{\sigma_1^2+\sigma_0^2-2\rho\sigma_1\sigma_0}{n}.
\]
The subtracted term is the variance of the unit effects divided by $n$, the Bayesian twin of the $S_\tau^2/n$ term of Exercise~\ref{ex:ch01:1}. Thus $\tau_{\mathrm{fp}}$ has no greater posterior variance than $\tau$, with equality only for constant effects: $\rho=1$ and $\sigma_1=\sigma_0$, when Neyman's estimator is exact. For $n_1=n_0=1$, $\sigma_1=\sigma_0=1$ and $\rho=0$, the variances are 1 and 2.

\paragraph{Exercise~\ref{ex:ch01:4}.} Draw both balls without replacement from an urn containing $0$ and $1$. The pair $(Z_1,Z_2)$ is $(0,1)$ or $(1,0)$ with probability $1/2$ each, so it is exchangeable, with covariance $-1/4$. A mixture of independent and identically distributed sequences has $\mathrm{Cov}(Z_1,Z_2)=\mathrm{Var}\{\E(Z_1\mid\theta)\}\geq0$, ruling out such a representation. For the first $k$ of $n$ draws the total variation distance from sampling with replacement is at most $k(k-1)/(2n)$: a repeated index has at most this probability, and conditional on no repeat the laws agree. Finite exchangeable sequences are mixtures over urns, giving the approximation of Diaconis and Freedman (1980). The approximation suffices in Section~\ref{sec:exchangeability}, where $k^2$ is small relative to $n$. For a target averaging all $n$ units, $k=n$ makes the bound uninformative, and the target is the mean of the urn itself and not a parameter of an infinite mixture.

\paragraph{Exercise~\ref{ex:ch01:6}.} In a two person household let $T_i=1$ mean vaccinated and $Y_i=1$ mean infected. Suppose infection enters only when neither is vaccinated:
\[
Y_i(t_1,t_2)=\one\{t_1=t_2=0\},\qquad Y_i=Y_i(T_1,T_2),\quad i=1,2.
\]
Each unit has four potential outcomes. The stable unit treatment value assumption requires $Y_1(t_1,0)=Y_1(t_1,1)$, which at $t_1=0$ says $1=0$. The effect must therefore specify the other unit's treatment: $\tau_1(t_2)=Y_1(1,t_2)-Y_1(0,t_2)$ is $-1$ at $t_2=0$ and $0$ at $t_2=1$, and likewise for unit 2. If unit 2 is vaccinated with probability $q$, averaging gives $-(1-q)$, which depends on the assignment law. Nor do these effects add to a campaign's effect: vaccinating both reduces infections from 2 to 0, while the two effects measured with the other unit vaccinated sum to 0.

\paragraph{Exercise~\ref{ex:ch01:7}.} Write $V(\delta)$ for the value net of cost. Under ignorability and given $\theta$, $\mu_1(x) - \mu_0(x) = \tau + \gamma^\top x$, and $\mu_0(x) = \alpha + \beta^\top x$ does not depend on the policy, so
\[
V(\delta) = \int \mu_0(x)\,Q(dx) + \int \delta(x)\,(\tau + \gamma^\top x - c)\,Q(dx).
\]
Averaging over the posterior, the first term does not involve $\delta$ and the second becomes $\int \delta(x)\,g(x)\,Q(dx)$ with $g(x) = \E(\tau + \gamma^\top x \mid \text{data}) - c$. Since $\delta(x) \in \{0, 1\}$, this is largest when $\delta(x) = 1$ wherever $g(x) > 0$ and $\delta(x) = 0$ wherever $g(x) < 0$, and ties do not matter. That is $\delta^\ast$. For any $\delta$,
\begin{align*}
V(\delta^\ast) - V(\delta) &= \int \{\delta^\ast(x) - \delta(x)\}(\tau + \gamma^\top x - c)\,Q(dx)\\
&= w(\tau - c) + u^\top\gamma,
\end{align*}
with $w = \int (\delta^\ast - \delta)\,dQ$ and $u = \int (\delta^\ast - \delta)\,x\,Q(dx)$ fixed once the data are given, so the difference is affine in $(\tau, \gamma)$, with constant term $-cw$. Take $\delta \equiv 0$. Then $w = p_\ast$ is the share of $Q$ that $\delta^\ast$ treats, $u = p_\ast\,m_\ast$ with $m_\ast$ the mean of $X$ under $Q$ among the treated, and $D = V(\delta^\ast) - V(0) = p_\ast(\tau + \gamma^\top m_\ast - c)$. If $(\tau, \gamma)$ has a normal posterior with mean $(\hat\tau, \hat\gamma)$ and covariance $\Sigma$, and $p_\ast > 0$, then with $b = (1, m_\ast)$ and $\Phi$ the standard normal distribution function
\[
\Prob(D > 0 \mid \text{data}) = \Phi\!\left(\frac{\hat\tau + \hat\gamma^\top m_\ast - c}{\sqrt{b^\top \Sigma\, b}}\right).
\]
The numerator is the average of $g$ over the treated, so it is positive and the probability exceeds $1/2$. With $\sigma^2$ unknown the posterior is Student $t$, $\Sigma$ is its scale matrix, and $\Phi$ becomes the $t$ distribution function with $\nu$ degrees of freedom. Given posterior draws $(\tau^{(s)}, \gamma^{(s)})$, $s = 1, \ldots, S$, the probability is the share of draws with a positive gain,
\[
\frac1S\sum_{s=1}^S \one\Bigl\{\int \delta^\ast(x)\,(\tau^{(s)} + \gamma^{(s)\top}x - c)\,Q(dx) > 0\Bigr\},
\]
which needs no normality.

\paragraph{Exercise~\ref{ex:ch01:8}.} Write $r_{ab}=\Prob\{Y(1)=a,Y(0)=b\}$. The margins give $r_{10}=p_1-r_{11}$ and $\max(0,p_1+p_0-1)\le r_{11}\le\min(p_1,p_0)$, hence
\[
\max(0,p_1-p_0)\le r_{10}\le\min(p_1,1-p_0).
\]
For $U\sim\mathrm{Unif}(0,1)$, take $Y(1)=\one\{U\le p_1\}$. Setting $Y(0)=\one\{U\le p_0\}$ gives $r_{11}=\min(p_1,p_0)$ and attains the lower bound; setting $Y(0)=\one\{U>1-p_0\}$ gives $r_{11}=\max(0,p_1+p_0-1)$ and attains the upper bound. For $p_1=0.6$, $p_0=0.3$, the two laws have $(r_{11},r_{10},r_{01},r_{00})=(0.3,0.3,0,0.4)$ and $(0,0.6,0.3,0.1)$. Finally the margins give $r_{10}=p_1-p_0+r_{01}$ under every joint law, so when $p_1\ge p_0$ the lower bound is attained exactly when $r_{01}=0$, that is, when no unit has $Y(1)=0$ and $Y(0)=1$.

\paragraph{Exercise~\ref{ex:ch01:9}.} The observed table can be anything: the imputed outcomes are linear in the observed ones and enter $\tau_{\mathrm{fp}}$ with fixed coefficients, so the variance across draws does not depend on them. We take $y_i \sim N(1, \sigma_1^2)$ for the 50 treated units and $y_i \sim N(0, \sigma_0^2)$ for the 50 controls, draw $S = 200{,}000$ pairs $(\mu_0, \mu_1)$ with $\mu_t \sim N(\bar y_t, \sigma_t^2/50)$, impute each missing outcome from its normal conditional law (mean $\mu_0 + \rho(\sigma_0/\sigma_1)(y_i - \mu_1)$ and variance $\sigma_0^2(1 - \rho^2)$ for a treated unit, the same with the arms exchanged for a control), and compute $\tau_{\mathrm{fp}}$ from the completed table. With $\sigma_1 = \sigma_0 = 1$ the variance of $\tau$ is $0.040$ at every $\rho$, and that of $\tau_{\mathrm{fp}}$ at $\rho = -1, -0.5, 0, 0.5, 1$ is $0.000$, $0.010$, $0.020$, $0.030$, $0.040$, the values of $0.040 - 0.02(1 - \rho)$ up to a Monte Carlo error of about $0.0003$. At $\rho = -1$ every missing outcome is $\mu_0 + \mu_1 - y_i$, the means cancel between the arms, and $\tau_{\mathrm{fp}} = \bar y_1 - \bar y_0$ in every draw. With $\sigma_0 = 2$ the variance of $\tau$ is $1/50 + 4/50 = 0.100$, and the formula $0.100 - (5 - 4\rho)/100$ gives $0.010$, $0.030$, $0.050$, $0.070$, $0.090$, which the simulation reproduces. At $\rho = 1$ the two variances are $0.090$ and $0.100$. There the unit effects are $\tau_i = \mu_1 - \mu_0 + (\sigma_1/\sigma_0 - 1)\{Y_i(0) - \mu_0\}$, with variance $(\sigma_1 - \sigma_0)^2 = 1$, so the subtracted term $1/n = 0.010$ does not vanish. The variances are equal only when the unit effects are constant, which takes $\rho = 1$ and $\sigma_1 = \sigma_0$.

\paragraph{Exercise~\ref{ex:ch01:moscow}.} Each group's mean has a Student $t$ posterior with location the group mean, scale $s_t/\sqrt{n_t}$ and $n_t - 1$ degrees of freedom, 184 and 2,978. With 184 degrees of freedom the 97.5 percent point is 1.973 against the normal 1.960, so both laws are close to normal and so is their difference. A Student $t$ variable with $\nu > 2$ degrees of freedom has variance $\nu/(\nu - 2)$, so the posterior variance of the effect is
\[
\frac{s_1^2}{n_1}\,\frac{n_1 - 1}{n_1 - 3} + \frac{s_0^2}{n_0}\,\frac{n_0 - 1}{n_0 - 3} = \frac{15.04^2}{185}\,\frac{184}{182} + \frac{16.48^2}{2979}\,\frac{2978}{2976} = 1.236 + 0.091 = 1.327
\]
squared points, a standard deviation of 1.15 points. Dropping the factors $(n_t - 1)/(n_t - 3)$ gives the large sample variance $s_1^2/n_1 + s_0^2/n_0 = 1.314$, one percent smaller. The interval $-8.53 \pm 1.96 \times 1.15$ is $(-10.8, -6.3)$, the interval of the 200,000 draws. The assigned group gives 93 percent of the variance: its spread is a little smaller, but it has a sixteenth as many stations. Halving the width means cutting the variance to a quarter, $1.314/4 = 0.328$. With $n_1$ assigned and $3164 - n_1$ unassigned stations, $15.04^2/n_1 + 16.48^2/(3164 - n_1) \le 0.328$ first holds at $n_1 = 1{,}182$, 37 percent of the city. The variance is smallest at $n_1 = 1{,}510$, where the group sizes are in proportion to the standard deviations (Neyman allocation), and there it is 0.314.

\paragraph{Exercise~\ref{ex:ch01:gujarat}.} Write $G_i(t) = A_i(t) - E_i(t)$ for the reporting gap, so $A_i(t) = E_i(t) + G_i(t)$. Because plants were assigned by lot, the mean audit reading in arm $t$ estimates $\E A(t)$, and the difference between the arms estimates
\[
\E A(1) - \E A(0) = \{\E E(1) - \E E(0)\} + \{\E G(1) - \E G(0)\},
\]
the effect on emissions plus the effect on reporting. Audit readings alone cannot split the sum: a reform that made auditors honest and plants cleaner could leave the mean reading unchanged. The backchecked plants are a random sample within each arm and their readings measure $E_i(t)$ without bias, so the backcheck means estimate $\E E(1)$ and $\E E(0)$, their difference is the effect on emissions, and the mean of audit minus backcheck over the backchecked plants of arm $t$ estimates $\E G(t)$. The paper reports the reporting term in just this way: on audit minus backcheck, the treatment effect is 0.15 to 0.21 standard deviations, against 0.30 standard deviations of underreporting in the control group. Its fall in emissions, 0.21 standard deviations, comes from a survey six months after the last audits, not from the backchecks, so the two numbers are not the two terms of one decomposition at one date. At control plants alone, audits against backchecks identify $\E G(0)$, the misreporting of the old system. That is a fact about plants and auditors that were observed, and it needs no lottery. It says nothing about $\E G(1)$ or $\E E(1)$, what auditors would report and plants would emit under a reform nobody had tried, and those need the treatment arm.

\section*{Chapter~\ref{ch:frameworks}: Frameworks for Causal Inference}

\paragraph{Exercise~\ref{ex:ch01:5}.} Under consistency, $p(Y \mid F_T = t, X) = p(Y(t) \mid X)$ and $p(Y \mid F_T = \emptyset, T = t, X) = p(Y(t) \mid T = t, X)$. Dawid's condition that these are equal for each $t$ is $p(Y(t) \mid X) = p(Y(t) \mid T = t, X)$. For binary $T$ with $0 < e(x) < 1$, the law $p(Y(t) \mid X)$ is a mixture of $p(Y(t) \mid T = t, X)$ and $p(Y(t) \mid T = 1 - t, X)$ with positive weights, so it equals the first only if it equals the second, and the condition is $Y(t) \indep T \mid X$ for each $t$. The joint condition $\{Y(0), Y(1)\} \indep T \mid X$ of the definition of ignorability (Definition~\ref{def:ignorability}) implies this and is stronger, because the regime condition involves the laws of $Y(0)$ and $Y(1)$ only one at a time. To see that it is strictly stronger, take $X$ trivial and let $Y(0)$ and $Y(1)$ be independent fair coin flips given $T = 0$ and equal to each other given $T = 1$. Each $Y(t)$ is a fair coin flip in both arms, so $Y(t) \indep T$ for $t = 0, 1$, but $\Prob(Y(0) = Y(1) = 1 \mid T = 0) = 1/4$ and $\Prob(Y(0) = Y(1) = 1 \mid T = 1) = 1/2$, so the pair is not independent of $T$. No identification of a total effect in this book needs more than the separate version; mediation does, since the natural direct and indirect effects of Chapter~\ref{ch:frameworks} need a cross world independence.

\paragraph{Exercise~\ref{ex:chframeworks:1}.} By consistency $Y = Y(1)$ when $T = 1$ and $Y = Y(0)$ when $T = 0$, so $\E[Y \mid T=1] - \E[Y \mid T=0] = \E[Y(1) \mid T=1] - \E[Y(0) \mid T=0]$. Adding and subtracting $\E[Y(0) \mid T=1]$ splits this into $\E[Y(1) \mid T=1] - \E[Y(0) \mid T=1]$, the effect on the treated, and $\E[Y(0) \mid T=1] - \E[Y(0) \mid T=0]$, the selection bias, which is the decomposition \eqref{fw:eq:decomp}. For the population, let $Y$ be a health score and let every unit gain 2 from treatment, $Y(1) = Y(0) + 2$. A unit is sick, with $Y(0) = 0$, or healthy, with $Y(0) = 4$. Suppose treatment goes mostly to the sick: 75\% of the treated and 25\% of the controls are sick. Then $\E[Y \mid T=1] = 0.75 \cdot 2 + 0.25 \cdot 6 = 3$ and $\E[Y \mid T=0] = 0.25 \cdot 0 + 0.75 \cdot 4 = 3$, so the naive contrast is zero. The effect on the treated is 2 and the selection bias is $\E[Y(0) \mid T=1] - \E[Y(0) \mid T=0] = 1 - 3 = -2$, which cancels it exactly.

\paragraph{Exercise~\ref{ex:chframeworks:2}.} The pooled data have $87 + 270 = 357$ small and $263 + 80 = 343$ large stones among 700 patients, weights $0.51$ and $0.49$. Weighting each procedure's stratum rates by these gives
\begin{align*}
\text{open:} \quad & 0.51\cdot\frac{81}{87} + 0.49\cdot\frac{192}{263} = 0.833,\\
\text{percutaneous:} \quad & 0.51\cdot\frac{234}{270} + 0.49\cdot\frac{55}{80} = 0.779.
\end{align*}
Open surgery leads by 5.4 points, in agreement with both strata and against the pooled comparison of Table~\ref{fw:tab:simpson}. The raw pooled rates, 78\% and 83\%, weight the strata differently, because 75\% of the open operations but only 23\% of the percutaneous ones were on large stones. The standardized rate describes a hypothetical population: the 700 patients of the study, with their own stone sizes, all given one procedure. It is the back door formula \eqref{fw:eq:backdoor} with the pooled stone sizes as $\Prob(X)$, and the difference between the two rates is the effect of the procedure only if stone size is the sole confounder. It is an accounting identity in the sense of Chapter~\ref{ch:mccullagh}: a known functional of the stratum rates and of a law $Q$ chosen by the analyst, here the pooled stone sizes, which does not enter the likelihood. Another $Q$ gives another number from the same table. For a clinic whose stones are 90\% large the standardized rates are $0.750$ and $0.705$, and the gap is 4.5 points.

\paragraph{Exercise~\ref{ex:chframeworks:3}.} Split $T$ into its natural value, the random node $T$, and the fixed value $t$. The natural value keeps the arrow from $X$ and loses every outgoing arrow; the fixed value has no parents and takes over the arrow into the outcome, which becomes the potential outcome $Y(t)$. The covariate is not a descendant of $T$, so it stays $X$.
\begin{center}
\begin{tikzpicture}[node distance=18mm]
\node[v] (X) {$X$};
\node[v,below left=of X] (T) {$T$};
\node[v,right=2mm of T] (t) {$t$};
\node[v,below right=of X] (Y) {$Y(t)$};
\draw[e] (X) -- (T); \draw[e] (X) -- (Y); \draw[e] (t) -- (Y);
\end{tikzpicture}
\end{center}
The only path from $T$ to $Y(t)$ is the fork $T\leftarrow X\rightarrow Y(t)$. It is open marginally and blocked by $X$, so $d$ separation gives ignorability, $Y(t)\indep T\mid X$ for each $t$. With consistency and overlap,
\[
\E[Y(t)]=\E_X\{\E[Y(t)\mid X]\}=\E_X\{\E[Y(t)\mid T=t,X]\}=\E_X\{\E(Y\mid T=t,X)\}.
\]
Each graph contains one potential outcome, so it does not imply the stronger joint condition $\{Y(0),Y(1)\}\indep T\mid X$. That condition holds in the structural model $Y(t)=f(t,X,\varepsilon)$, $T=g(X,\eta)$ with $\varepsilon\indep\eta\mid X$, since both potential outcomes are functions of $(X,\varepsilon)$.

\paragraph{Exercise~\ref{ex:chframeworks:4}.} For $\bar a=(1,1)$ write $e_0=\Prob(A_0=1)$, $e_1(\ell)=\Prob(A_1=1\mid A_0=1,L_1=\ell)$ and $m(\ell)=\E[Y\mid A_0=1,L_1=\ell,A_1=1]$. Assume $e_0>0$ and $e_1(\ell)>0$ wherever $\Prob(L_1=\ell\mid A_0=1)>0$. With $W=\one\{A_0=1,A_1=1\}/\{e_0e_1(L_1)\}$,
\[
\E[WY]=\sum_\ell\frac{e_0\Prob(L_1=\ell\mid A_0=1)e_1(\ell)m(\ell)}{e_0e_1(\ell)}=\sum_\ell m(\ell)\Prob(L_1=\ell\mid A_0=1),
\]
the g formula. For continuous $L_1$, replace sums by integrals. In Figure~\ref{fw:fig:gmethods}, $A_0$ has no parent and $L_1$ blocks the path from $U$ to $A_1$, giving sequential ignorability and identifying the common value as $\E[Y(1,1)]$. Correct models for $\Prob(L_1\mid A_0)$ and $\E[Y\mid A_0,L_1,A_1]$, or for $e_0$ and $e_1$, make the respective estimators consistent. They agree in large samples, though generally differ in finite samples.

\paragraph{Exercise~\ref{ex:chframeworks:5}.} Every unit has $\tau_i = \pm 1$, so no unit has $Y_i(1) = Y_i(0)$ and Fisher's sharp null is false, while $\tau_{\mathrm{fp}} = \frac{1}{2}(1) + \frac{1}{2}(-1) = 0$, so Neyman's weak null is true. A randomized experiment identifies the laws of $Y(1)$ and of $Y(0)$, and it never shows both outcomes of one unit, so it reaches the proportion harmed only through bounds. A binary outcome shows the limit most plainly. Let higher outcomes be better and use the $r_{ab}$ of Exercise~\ref{ex:ch01:8}, so the proportion harmed is $r_{01}$. Exchanging the roles of the two outcomes in that exercise gives $\max(0, p_0 - p_1) \le r_{01} \le \min(p_0, 1 - p_1)$. Here half are helped ($r_{10} = \frac{1}{2}$) and half harmed ($r_{01} = \frac{1}{2}$), so $p_1 = p_0 = \frac{1}{2}$ and the bounds are $[0, \frac{1}{2}]$. The value 0 belongs to the law with $r_{11} = r_{00} = \frac{1}{2}$, under which $Y(1) = Y(0)$ for every unit. That world, where the sharp null holds, and the world of this exercise give the same law of $(T, Y)$, so no sample size separates them. With other outcomes the law of $Y(1)$ can move, and then the experiment can exclude zero. For any $y$, if $Y(1) \le y$ and $Y(0) > y$ then the unit is harmed, and if $Y(0) \le y$ and $Y(1) > y$ then it is not, so with $F_t$ the distribution function of $Y(t)$,
\[
\sup_y\,\{F_1(y) - F_0(y)\} \;\le\; \Prob\{Y(1) < Y(0)\} \;\le\; 1 - \sup_y\,\{F_0(y) - F_1(y)\}.
\]
Let $Y(0) \sim N(0, 1)$ with the sign of the effect independent of $Y(0)$, so $F_1(y) = \frac{1}{2}\Phi(y - 1) + \frac{1}{2}\Phi(y + 1)$. Both suprema equal 0.103, attained at $y = -1.085$ and $y = 1.085$, so the proportion harmed lies in $[0.103, 0.897]$ against a true value of $\frac{1}{2}$. The lower end excludes the sharp null, through the spread of $Y(1)$ (variance 2 against 1), although the difference in means has expectation 0. The exact answer needs the premise itself. If every effect is known to be $+1$ or $-1$, then $\tau_{\mathrm{fp}} = 1 - 2\Prob\{Y(1) < Y(0)\}$, and $\tau_{\mathrm{fp}} = 0$ gives $\frac{1}{2}$, but this comes from the assumption and not from the randomization. For a Bayesian the same limit appears as a posterior that the data move only through the margins, with the rest set by the prior on the association of $Y(0)$ and $Y(1)$, the parameter $\rho$ of Chapter~\ref{ch:predictive}.

\paragraph{Exercise~\ref{ex:chframeworks:6}.} The four cells $(T, Y) = (0,0), (0,1), (1,0), (1,1)$ have probability $1/4 = \tfrac12\cdot\tfrac12$ each, so $T$ and $Y$ are independent. Only the cell $(0,0)$ has $C = 0$. Restricting to $C = 1$ leaves $(0,1)$, $(1,0)$ and $(1,1)$, each with probability $1/3$, so
\begin{align*}
\Prob(Y=1 \mid T=1, C=1) &= \frac{1/3}{2/3} = \frac12,\\
\Prob(Y=1 \mid T=0, C=1) &= \frac{1/3}{1/3} = 1.
\end{align*}
Given $C = 1$ the coins are dependent, with a risk difference of $-1/2$, although neither affects the other and each is a fair coin. Both coins cause $C$, a collider, and conditioning on a common effect makes its causes informative about each other: once $C = 1$ is known, $T = 0$ forces $Y = 1$. A regression of $Y$ on $T$ and $C$ would report a negative effect of $T$ on $Y$ where there is none. Selecting units on $C$ produces the same dependence without any regression: in 3000 units with $C = 1$, about 1000 fall in each of the three cells, and the share with $Y = 1$ is about 0.5 among those with $T = 1$ and exactly 1 among those with $T = 0$.

\paragraph{Exercise~\ref{ex:chframeworks:7}.} The target has the same stratum effects, $\E^\ast[Y(1) - Y(0) \mid X = x] = \tau(x)$, so by the law of total expectation in the target population
\begin{align*}
\E^\ast[Y(1) - Y(0)] &= \sum_x \Prob^\ast(X = x)\,\E^\ast[Y(1) - Y(0) \mid X = x]\\
&= \sum_x \tau(x)\,\Prob^\ast(X = x).
\end{align*}
This needs $\tau(x)$ at every $x$ the target has, so $\Prob(X = x) > 0$ whenever $\Prob^\ast(X = x) > 0$. In Chapter~\ref{ch:mccullagh} an average effect is $\tau_Q = \int \tau(x)\,Q(dx)$, the stratum effects averaged over a law $Q$ chosen by the analyst and not part of the likelihood. The display is this with $Q = \Prob^\ast$: the source data supply $\tau(x)$ and the target supplies the weights. With $w(x) = \Prob^\ast(X = x)/\Prob(X = x)$, which has mean 1 under the source, the target and source averages differ by
\begin{align*}
\sum_x \tau(x)\{\Prob^\ast(X = x) - \Prob(X = x)\} &= \E[\tau(X)\{w(X) - 1\}]\\
&= \mathrm{Cov}\{\tau(X), w(X)\},
\end{align*}
with expectation and covariance taken under the source. They coincide exactly when the effect is uncorrelated with the weight, in particular when $\tau(x)$ is constant or the two covariate laws are equal. If $\tau = 1$ for young and $\tau = 3$ for old patients, and the share of old patients is $0.2$ in the source and $0.6$ in the target, the averages are $1.4$ and $2.2$, and $\mathrm{Cov}\{\tau(X), w(X)\} = 0.8$.

\paragraph{Exercise~\ref{ex:chframeworks:8}.} Let $\mathcal T$ be the finite set of allocations, $Y_t$ the outcomes under $t$, $\tau$ the average effect and $L(a,\tau)$ the reporting loss. A fixed allocation has expected posterior loss $R(t)=\E[\min_a\E\{L(a,\tau)\mid x,t,Y_t\}\mid x,t]$, averaging over the predictive law of $Y_t$. For the first Bayesian, $p(y(0),y(1),\theta\mid x)$ does not depend on $t$. A randomized design contributes only $\pi(t\mid x)$ to her likelihood, which cancels, leaving the same posterior and predictive law as a fixed allocation $t$. Thus $R(\pi)=\sum_{t\in\mathcal T}\pi(t\mid x)R(t)\ge\min_tR(t)$, with equality exactly when $\pi$ is supported on minimizers. For an example let $m$ be even, $m$ patients have $x = 1$ and $m$ have $x = 0$, and $Y = \mu + \tau t + \beta x + \varepsilon$ with $\varepsilon \sim N(0, \sigma^2)$, $\sigma$ known, and flat priors. If a balanced allocation treats $a$ of those with $x = 1$, the residual sum of squares of $t$ on $(1, x)$ is $m/2 - 2(a - m/2)^2/m$, and under squared error $R(t) = \sigma^2/\{m/2 - 2(a - m/2)^2/m\}$. The minimum, $2\sigma^2/m$, is at $a = m/2$, and complete randomization mixes in every other $a$, including $a = 0$ and $a = m$, where $\tau$ is not identified.

The second Bayesian adds an unrecorded prognosis $U$ and a mechanism $p(t \mid x, u, y(0), y(1), \psi)$ with a prior on $\psi$. With $y_{\mathrm{mis}}$ the unobserved potential outcomes, her posterior is
\[
p(\theta, \psi, u, y_{\mathrm{mis}} \mid x, t, y_t) \propto p(\theta)\,p(\psi)\,p(u, y(0), y(1) \mid x, \theta)\,p(t \mid x, u, y(0), y(1), \psi).
\]
Under random allocation the last factor is $\pi(t\mid x)$ and cancels, so complete randomization, blocks and an $x$ dependent coin give the same posterior for given $(t,y_t)$. Hand allocation leaves a factor depending on $u$. For an exact example take $m$ recorded pairs with $Y_{ji}(t)=\mu_j+\tau t+U_{ji}+\varepsilon_{ji}$, independent $U_{ji}\sim N(0,\kappa^2)$ and $\varepsilon_{ji}\sim N(0,\sigma^2)$, and flat priors on $\mu_j$. Each pair has one treated patient, chosen independently to have the larger $U$ with probability $\lambda$. The pair sum is uninformative under the flat prior and independent of the choice, since normal sums and differences are independent. Only the treated minus control difference $D_j$ matters: $D_j-\tau$ is $|V|+e$ with probability $\lambda$ and $-|V|+e$ otherwise, for independent $V\sim N(0,2\kappa^2)$ and $e\sim N(0,2\sigma^2)$. At $\lambda=1/2$ this is $N(0,s^2)$, $s^2=2\kappa^2+2\sigma^2$, so a flat prior gives $\tau\mid D\sim N(\bar D,s^2/m)$. Otherwise $\E D_j=\tau+(2\lambda-1)2\kappa/\sqrt\pi$. With $\kappa=\sigma=1$, $m=50$ and $\lambda=0.9$, the shift is 0.90 against a fair coin posterior standard deviation of 0.28. The same allocation and outcomes therefore give different posteriors under the two mechanisms. For unknown $\lambda$, $p(\tau\mid D)=\int p(\tau\mid D,\lambda)p(\lambda\mid D)\,d\lambda$ carries her suspicion into the answer.

This illustrates the public part of Kadane and Seidenfeld's answer. By her own lights the first Bayesian needs no randomization, and the first part of the exercise is the theorem behind that. Randomization protects against selection on unrecorded variables and gives observers with different priors a common likelihood for the allocation, so that their disagreement is confined to the science. The two can settle on a fair coin within the recorded pairs: exchangeability within pairs gives every one of the $2^m$ allocations the same $R(t)$ for the first Bayesian, so the mixture costs her nothing, and the second has nothing left to model.

\section*{Chapter~\ref{ch:mccullagh}: Models, Accounting and Randomization}

\paragraph{Exercise~\ref{ex:ch02:1}.} In stratum A the odds ratio is $(0.80/0.20)/(0.50/0.50) = 4$ and in stratum B it is $(0.50/0.50)/(0.20/0.80) = 4$. With two strata of equal size the pooled risks are $(0.50 + 0.20)/2 = 0.35$ and $(0.80 + 0.50)/2 = 0.65$, and the pooled odds ratio is $(0.65/0.35)/(0.35/0.65) = 3.45$, the entries of Table~\ref{tab:noncollapse}. For a marginal odds ratio below 2 with conditional odds ratio 4, keep two strata of equal size and spread the baseline risks: take stratum A with control risk 0.95 and treated risk $4 \times 19/(1 + 4 \times 19) \approx 0.987$, and stratum B with control risk 0.05 and treated risk $4/23 \approx 0.174$. The pooled risks are 0.50 and 0.580, with odds ratio about 1.38.

\paragraph{Exercise~\ref{ex:ch02:2}.} Take two strata $k = A, B$ with probabilities $q_k$, baseline hazards $h_A(t) > h_B(t)$ and cumulative hazards $H_k$, and randomize treatment with the same probability in each stratum, so that there is no confounding. Let the conditional hazard ratio be $c > 1$ in both strata. Arm $a \in \{0, 1\}$ has survival function $S_a(t) = \sum_k q_k e^{-c^a H_k(t)}$ and marginal hazard
\[
h_a(t) = -\frac{d}{dt}\log S_a(t) = c^a \sum_k w_{ak}(t)\,h_k(t), \qquad w_{ak}(t) = \frac{q_k e^{-c^a H_k(t)}}{S_a(t)},
\]
the stratum hazard averaged over the survivors of the arm. The weight of the high hazard stratum, $w_{aA} = \{1 + (q_B/q_A)e^{c^a(H_A - H_B)}\}^{-1}$, is smaller for $a = 1$ than for $a = 0$ at every $t > 0$, because treatment removes members of that stratum faster. A weighted average of $h_A$ and $h_B$ rises with the weight on $A$, so
\[
\frac{h_1(t)}{h_0(t)} = c\,\frac{w_{1A}h_A + (1 - w_{1A})h_B}{w_{0A}h_A + (1 - w_{0A})h_B} < c \qquad (t > 0).
\]
The marginal hazard ratio equals $c$ only at $t = 0$, and the marginal hazards are not proportional (for $c < 1$ the inequality reverses). With $q_A = q_B = 1/2$, $h_A = 1$, $h_B = 0.1$ and $c = 2$ the ratio is 1.26 at $t = 1$ and 1.09 at $t = 2$. It returns to 2 only as nearly all survivors come from stratum $B$ (1.998 at $t = 10$).

\paragraph{Exercise~\ref{ex:ch02:3}.} Write $\tau_Q(\theta) = \int \{\expit(\eta_1(x)) - \expit(\eta_0(x))\}Q(dx)$ with $\eta_t(x) = \alpha + \beta t + \gamma^\top x$. Its gradient with respect to $\theta = (\alpha, \beta, \gamma)$ is $g = \int \{w_1(x)(1, 1, x) - w_0(x)(1, 0, x)\}Q(dx)$, where $w_t = \expit(\eta_t)\{1 - \expit(\eta_t)\}$. To first order the posterior mean is $\tau_Q(\hat\theta)$ and the posterior variance is $g^\top V g$ with $g$ evaluated at $\hat\theta$, where $\hat\theta$ and $V$ are the posterior mean and covariance.

\paragraph{Exercise~\ref{ex:ch02:4}.} With errors $e_{ij}\sim N(0,\sigma_e^2)$ independent across units and of $b_i$, the cluster covariance matrix is $\sigma_e^2I+(\sigma_b^2/m_i)\mathbf1\mathbf1^\top$. Coherence asks that the model for $m_i$ units be the marginal of the model for $m_i+1$ units. Marginalizing the larger model leaves pairwise covariance $\sigma_b^2/(m_i+1)$, whereas the model for $m_i$ units requires $\sigma_b^2/m_i$. The same children's joint law changes with the sample size. Instead take $b_i\sim N(0,\sigma_b^2)$ with $\sigma_b^2$ free of $m_i$. Conditional on $b_i$ the units are independent $N(\mu+b_i,\sigma_e^2)$, so marginalization is coherent. The cluster mean variance $\sigma_b^2+\sigma_e^2/m_i$ still falls with sample size. If larger clusters should have smaller effects, let $\mathrm{Var}(b_i)$ depend on the cluster population size $M_i$, which does not change with sampling.

\paragraph{Exercise~\ref{ex:ch02:5}.} The interaction makes the integrand of \eqref{eq:tauQ} $\expit(\alpha+\beta+(\gamma+\delta)^\top x)-\expit(\alpha+\gamma^\top x)$, which has the sign of $\beta+\delta^\top x$. Without it the sign is that of $\beta$ at every $x$. Choose $Q$ supported where $\beta+\delta^\top x>0$ and $Q'$ where it is negative. For a posterior concentrated near $\hat\theta$, the posterior means are close to $\tau_Q(\hat\theta)>0$ and $\tau_{Q'}(\hat\theta)<0$ by Exercise~\ref{ex:ch02:3}. In the worked example $\alpha=-0.5$, $\beta=0.8$, $\gamma=0.6$ and $\delta=-0.7$, so the logit effect $0.8-0.7a$ changes sign at $a=8/7$. Standardized age laws $N(0,1)$ and $N(2,1)$ give $\tau_Q=0.19$ (0.158 in the chapter's data) and $\tau_Q=-0.13$: one analyst reports a benefit and the other a harm. Opposite signs require different covariate laws, as eligibility criteria can produce; a random sample nearly shares the population law and mean.

\paragraph{Exercise~\ref{ex:ch02:6}.} ICH E9(R1) names five attributes. Treatment: three quarters of a pint of milk a day for four months ($T = 1$) against no supplement ($T = 0$). Population: the schoolchildren of Lanarkshire in 1930, with $Q$ the law of their covariates $X$ (age, sex, school, initial weight and height). Variable: the four month gain $G$ in weight, and separately in height. Intercurrent events: illness, absence and a change of school, handled by a treatment policy strategy, so that $G$ is used whatever happened. Summary: the difference in mean gain. With $\mu(t, x; \theta) = \E(G \mid T = t, X = x, \theta)$ the estimand is
\[
\tau_Q(\theta) = \int \{\mu(1, x; \theta) - \mu(0, x; \theta)\}\,Q(dx).
\]
A different population changes $Q$ (for poor children, the covariate law of that subgroup), a different summary changes the functional, and a different strategy for intercurrent events changes what $G$ and the model describe: a hypothetical strategy for children who stopped drinking the milk needs an adherence model. Each draw of $\theta$ gives a draw of $\tau_Q$, so the prespecified $Q$ and functional define the estimand and the posterior supplies its estimate.

\paragraph{Exercise~\ref{ex:ch03:1}.} In the simulation, gain is $1 - 0.3S + 0.2T + \varepsilon$ and $T$ depends on $S$. Write the growth slope as $b$ ($b = -0.3$ here) and the selection as $\Prob(T = 1 \mid S) = \expit(\kappa S)$ with $S$ standard normal ($\kappa = -1.2$ here). The naive difference has expectation $0.2 + b\{\E(S \mid T = 1) - \E(S \mid T = 0)\}$. By symmetry $\Prob(T = 1) = 1/2$, so $\E(S \mid T = 1) = 2\E\{S\expit(\kappa S)\}$ and $\E(S \mid T = 0) = -2\E\{S\expit(\kappa S)\}$, and Stein's lemma gives $\E\{S\expit(\kappa S)\} = \kappa\E\{\expit(\kappa S)(1 - \expit(\kappa S))\}$. The bias is therefore
\[
4b\kappa\,\E\{\expit(\kappa S)(1 - \expit(\kappa S))\},
\]
linear in the growth slope and close to $b\kappa$ for a small selection slope, where the expectation is near $1/4$. At $\kappa = -1.2$ the expectation is 0.194, the treated have lower $S$ by 0.93 on average, and the bias is $0.28$, so the expected naive estimate is $0.48$, close to the simulated value of $0.49$.

\paragraph{Exercise~\ref{ex:ch03:2}.} Let $(S, U)$ be standard bivariate normal with correlation $r$, let the gain be $G = a + bS + cU + \tau T + \varepsilon$ with independent noise $\varepsilon$, and let assignment depend on $U$ alone. Then $S = rU + \sqrt{1 - r^2}\,W$ with $W$ independent of $(U, T)$. The population regression of $G$ on $T$ and $S$ has coefficient $\tau + c\,\delta_U$ on $T$, where $\delta_U$ is the coefficient of $T$ in the linear regression of $U$ on $T$ and $S$. With $d = \mathrm{Cov}(U, T)$ and $v = \mathrm{Var}(T)$ we have $\mathrm{Cov}(S, T) = rd$, and
\[
\delta_U = \frac{d - r\cdot rd}{v - (rd)^2} = \frac{(1 - r^2)\,d}{v - r^2d^2}.
\]
Write $\Delta_U = d/v = \E(U \mid T = 1) - \E(U \mid T = 0)$ and $\rho^2 = d^2/v$, the squared correlation of $U$ with assignment. The bias is
\[
c\,\delta_U = c\,\Delta_U\,\frac{1 - r^2}{1 - r^2\rho^2}.
\]
At $r = 0$ it is the full imbalance bias $c\Delta_U$ of the naive comparison, and at $|r| = 1$ it vanishes, because $U = \pm S$ is adjusted for. In between, adjusting for $S$ removes slightly less than the fraction $r^2$ of the bias. Take the selection rule of the chapter with $U$ in place of $S$, $\Prob(T = 1 \mid U) = \expit(-1.2\,U)$, and $c = -0.3$. Then $\Delta_U = -0.93$ and $\rho^2 = 0.22$, and the bias is $0.28$ at $r = 0$ and $0.22$, $0.12$ and $0.03$ at $r = 0.5$, $0.8$ and $0.95$.

\paragraph{Exercise~\ref{ex:ch03:3}.} Within a twin pair the shared effect cancels, so with $m$ pairs the variance of the mean difference is $2\sigma_\varepsilon^2/m$. With $N$ unpaired children split equally between arms it is $4(\sigma_\alpha^2 + \sigma_\varepsilon^2)/N$, since each child carries the full variance. Equality at $m = 50$, $N = 20{,}000$ requires $\sigma_\varepsilon^2/25 = (\sigma_\alpha^2 + \sigma_\varepsilon^2)/5000$, that is $\sigma_\alpha^2 + \sigma_\varepsilon^2 = 200\sigma_\varepsilon^2$, so $\sigma_\alpha^2/\sigma_\varepsilon^2 = 199$. Student's claim holds against an ideal randomized trial only if family effects dominate individual variation by a factor of about two hundred; against the actual, biased trial it holds much more easily, because bias does not shrink with $N$.

\paragraph{Exercise~\ref{ex:ch03:4}.} Let $W_{is}$ be the true weight of child $i$ in February ($s = 0$) and June ($s = 1$), $C_{is}$ the weight of clothing, $R_{is} = W_{is} + C_{is}$ the recorded weight, $P_i \in \{0, 1\}$ poverty as the teachers judged it, and $X_i$ the recorded covariates. The model is
\[
W_{i1} - W_{i0} = \mu + \beta^\top X_i + \tau T_i + \varepsilon_i, \qquad C_{is} = c_s - \zeta_s P_i + \eta_{is},
\]
with $\varepsilon_i \sim N(0, \sigma^2)$, $\eta_{is} \sim N(0, \sigma_C^2)$ and $T_i$ depending on $(X_i, P_i)$. The poor wear $\zeta_s$ less than the non-poor in season $s$. The seasonal drop $c_0 - c_1$ of the non-poor has a normal prior centred on Student's one pound, and $\zeta = \zeta_0 - \zeta_1$, the amount by which the poor shed less, has a prior on $[0, c_0 - c_1]$. Weak priors on $\mu$, $\beta$, $\tau$, $\sigma^2$ and $\sigma_C^2$ complete the model. When poverty is recorded only as a share of poor children per school, $P_i$ is a latent Bernoulli variable whose probability is that share. Since $R_{i1} - R_{i0} = (W_{i1} - W_{i0}) - (C_{i0} - C_{i1})$,
\[
R_{i1} - R_{i0} = \mu - (c_0 - c_1) + \beta^\top X_i + \tau T_i + \zeta P_i + \text{noise}.
\]
The naive milk minus control comparison of recorded gains has expectation $\tau+\beta^\top\{\E(X\mid T=1)-\E(X\mid T=0)\}+\zeta\Delta_P$, where $\Delta_P=\Prob(P=1\mid T=1)-\Prob(P=1\mid T=0)$. Teachers favoured poor children, so $\Delta_P>0$: with $\zeta=0.5$ pounds, 60 percent poor children among those given milk and 30 percent among controls, clothing adds $0.5\times0.3=0.15$ pounds to the apparent gain whatever the true effect. This needs a season by poverty interaction; if $\zeta_0=\zeta_1$, the poverty term cancels. With latent $P_i$ the likelihood sees only $\tau T_i+\zeta P_i$. The priors carry the information, and to first order the posterior of $\tau$ is the comparison adjusted for $X$ less $\zeta\Delta_P$, propagating both factors' uncertainty.

\paragraph{Exercise~\ref{ex:ch03:5}.} Complete randomization is uniform on $\mathcal T$, the $N=\binom n{n_1}$ assignments with $n_1$ treated. Under the sharp null $Y_i(1)=Y_i(0)=y_i$, outcomes are fixed, and a statistic $s(T,y)$ has the $p$ value
\[
p(T)=N^{-1}\#\{t\in\mathcal T:s(t,y)\ge s(T,y)\}.
\]
For nonempty $A=\{t:p(t)\le\alpha\}$, choose $t^*$ minimizing $s$ on $A$. Then $A\subseteq\{t:s(t,y)\ge s(t^*,y)\}$, so $|A|\le Np(t^*)\le N\alpha$. Thus $\Prob\{p(T)\le\alpha\}\le\alpha$, trivially also for empty $A$, with equality if $N\alpha$ is an integer and there are no ties. Exactness holds for any outcomes and statistic because the assignment law is known. Lanarkshire's substitutions depended on size and nourishment, making $T$ nonuniform with unknown law. The mass of $A$ can then far exceed $\alpha$ when children likely to grow most receive milk. In 200 simulations of the chapter's mechanism with zero milk effect, 2,000 children each and 999 permutations per test, the mean difference averages 0.28 against a permutation standard deviation of 0.026; the 5 percent test rejects every time.

\paragraph{Exercise~\ref{ex:ch03:6}.} One formalization is a game of incomplete information. The effect is $\tau \in \{0, \delta\}$ with common prior $\Prob(\tau = \delta) = \pi$, and the experimenter knows no more about $\tau$ than the reader does, so her choices carry no news about it. Nature makes the experimenter honest with probability $q$ and partisan with probability $1 - q$, and only she knows which. She picks a design: r, a verifiable random allocation at cost $k \ge 0$, or h, allocation by hand. The trial then reports favourable ($G$) or not ($N$). Under r, and under h by an honest experimenter, $\Prob(G \mid \tau = \delta) = a_1$ and $\Prob(G \mid \tau = 0) = a_0$, with $a_1 > a_0$. Under h by a partisan, who gives the treatment the patients with good prognosis, these become $b_1 \ge a_1$ and $b_0 > a_0$. The reader sees the design and the report and accepts or rejects the claim. He gets $1 - c$ for accepting a true claim, $-c$ for accepting a false one and 0 for rejecting. She gets 1 if he accepts, less $k$ if she randomized. He accepts when the posterior probability of $\tau = \delta$ is at least $c$, that is, when the likelihood ratio of what he saw is at least $\ell = c(1 - \pi)/\{(1 - c)\pi\}$. If he believes that a share $q_{\mathrm{h}}$ of those using h are honest, the ratio after h and $G$ is
\[
\Lambda(q_{\mathrm{h}}) = \frac{q_{\mathrm{h}} a_1 + (1 - q_{\mathrm{h}}) b_1}{q_{\mathrm{h}} a_0 + (1 - q_{\mathrm{h}}) b_0},
\]
which increases from $b_1/b_0$ at $q_{\mathrm{h}} = 0$ to $a_1/a_0$ at $q_{\mathrm{h}} = 1$, since the assumptions that follow imply $a_1/a_0 > b_1/b_0$. Assume a clean favourable result persuades, $a_1/a_0 > \ell$, an unfavourable result never does, $(1 - a_1)/(1 - a_0) < \ell$ and $(1 - b_1)/(1 - b_0) < \ell$, and the reader is sceptical, $\Lambda(q) < \ell$, so that $q$ lies below the root $q^*$ of $\Lambda(q^*) = \ell$. Let $g = \pi a_1 + (1 - \pi) a_0$ and $g_{\mathrm{P}} = \pi b_1 + (1 - \pi) b_0$, so $g_{\mathrm{P}} > g$, and take $k < g$.

Then in every perfect Bayesian equilibrium both types randomize with probability one, and the reader accepts exactly on $G$. Under r the signal is clean whatever the type, so he accepts on $G$ only and r pays $g - k > 0$. Let $\rho$ be the probability that he accepts after h and $G$; he rejects after h and $N$. Then h pays $\rho g$ to the honest type and $\rho g_{\mathrm{P}}$ to the partisan. Suppose h is used with positive probability. If $\rho = 0$, h pays 0, which is below $g - k$, a contradiction. If $\rho > 0$, then $\Lambda(q_{\mathrm{h}}) \ge \ell$, so $q_{\mathrm{h}} \ge q^* > q$, which by Bayes' rule requires the honest type to use h more often than the partisan. So the honest type uses h, $\rho g \ge g - k$, and then $\rho g_{\mathrm{P}} > \rho g \ge g - k$, so the partisan uses h for certain, and $q_{\mathrm{h}} \le q$, a contradiction. An equilibrium exists: with off path belief $q_{\mathrm{h}} = q$ the reader rejects after h, and h pays 0.

This is the protection against suspicion. An honest experimenter using h is rejected and loses $g-k$, because anything she could say in her favour a partisan could say too. The random allocation is the one move whose meaning does not depend on who makes it. The reader's doubt matters: if $q\ge q^*$ and $k>0$, both types using h with acceptance on $G$ is an equilibrium, and randomization is not bought.

\section*{Chapter~\ref{ch:groups}: Groups, Urns and Information}

\paragraph{Exercise~\ref{ex:gr:1}.} (a) Since $1-g_\psi(p)=(1-p)/\{1+(e^{\psi}-1)p\}$, the odds of $g_\psi(p)$ are $e^{\psi}$ times the odds of $p$, so $\logit g_\psi(p)=\logit p+\psi$ and $g_\psi=\expit\circ(\,\cdot+\psi)\circ\logit$. Translations of the line compose by adding, and conjugation by $\logit$ carries this over to $g_{\psi_1}\circ g_{\psi_2}=g_{\psi_1+\psi_2}$. Also $g_0$ is the identity, $g_{-\psi}=g_\psi^{-1}$, and $g_\psi(\tfrac12)=\expit(\psi)$ determines $\psi$, so $\psi\mapsto g_\psi$ is one to one and an isomorphism of $(\R,+)$ with a group of bijections of $(0,1)$. (b) Take $\psi>0$ and write $p=\expit\gamma$. With $\expit a=\{1+\tanh(a/2)\}/2$, $\tanh x-\tanh y=\sinh(x-y)/(\cosh x\cosh y)$ and $2\cosh x\cosh y=\cosh(x+y)+\cosh(x-y)$,
\[
g_\psi(p)-p=\tfrac12\Bigl\{\tanh\frac{\gamma+\psi}{2}-\tanh\frac{\gamma}{2}\Bigr\}=\frac{\sinh(\psi/2)}{\cosh(\gamma+\psi/2)+\cosh(\psi/2)}.
\]
The numerator is positive and does not depend on $\gamma$, and the denominator is smallest, with $\cosh(\gamma+\psi/2)=1$, at $\gamma=-\psi/2$ and nowhere else. The maximum is $\sinh(\psi/2)/\{1+\cosh(\psi/2)\}=\tanh(\psi/4)$ by the half angle formulas. The display is symmetric about $-\psi/2$, the baseline at which the control and treated risks lie symmetrically about one half. (c) The average $\tau_Q(\psi)=\E_Q\{\expit(\gamma+\psi)-\expit(\gamma)\}$ is an average of numbers each at most $\tanh(\psi/4)$, so $\tau_Q(\psi)\le\tanh(\psi/4)$ for every $Q$ and every $\psi>0$, the bound of Proposition~\ref{gr:prop:range}. An average effect $\tau$ therefore needs $\psi\ge4\operatorname{artanh}\tau$, and
\[
e^{\psi}\ge e^{4\operatorname{artanh}\tau}=\Bigl(\frac{1+\tau}{1-\tau}\Bigr)^2,\qquad\Bigl(\frac{1.2}{0.8}\Bigr)^2=2.25\quad(\tau=0.2).
\]
The same formula gives the $2.57$ and $1.36$ of Section~\ref{gr:sec:many} for the averages $0.232$ and $0.076$.

\paragraph{Exercise~\ref{ex:gr:2}.} Write $y_i=Y_i(1)$, so that $\bar Y_1=n_1^{-1}\sum_iI_iy_i$ and only the $I_i$ are random. With $c=\mathrm{Var}(I_i)=n_1(n-n_1)/n^2$ the covariance is $-c/(n-1)$, and using $\sum_{i\ne j}y_iy_j=(\sum_iy_i)^2-\sum_iy_i^2$,
\[
\mathrm{Var}\Bigl(\sum_iI_iy_i\Bigr)=c\Bigl\{\sum_iy_i^2-\frac1{n-1}\sum_{i\ne j}y_iy_j\Bigr\}=\frac{cn}{n-1}\Bigl\{\sum_iy_i^2-\frac1n\Bigl(\sum_iy_i\Bigr)^2\Bigr\}.
\]
For a binary urn $y_i^2=y_i$ and $\sum_iy_i=np_1$, so the braces equal $np_1(1-p_1)$ and
\[
\mathrm{Var}(\bar Y_1)=\frac1{n_1^2}\cdot\frac{n_1(n-n_1)}{n^2}\cdot\frac{n}{n-1}\cdot np_1(1-p_1)=\frac{p_1(1-p_1)}{n_1}\cdot\frac{n-n_1}{n-1}.
\]
For $n=10$, $n_1=5$ and four ones, $p_1=0.4$ and the variance is $(0.24/5)(5/9)=0.0267$, against $0.048$ for five independent draws. The urn variance is $44$ percent smaller, the finite population correction $5/9$ at work, and the negative covariances of the indicators produce it, since a unit that has been drawn makes each other unit slightly less likely to be drawn. For the check, let $X$ be the number of ones among the five units drawn. It is hypergeometric, $\Prob(X=k)=\binom4k\binom6{5-k}/\binom{10}5$, with numerators $6,60,120,60,6$ over $252$ for $k=0,\dots,4$. Then $\E X=2$, $\E X^2=1176/252=14/3$, $\mathrm{Var}(X)=2/3$ and $\mathrm{Var}(\bar Y_1)=\mathrm{Var}(X)/25=0.0267$, the same value.

\paragraph{Exercise~\ref{ex:gr:3}.} Both arms of this urn have $p_1=p_0=\tfrac12$, so they look alike, yet the unit effects are $+1$ for five units, $-1$ for five and $0$ for ten. The finite population variances with divisor $n-1=19$ are $S_1^2=S_0^2=20\cdot\tfrac14/19=5/19$ and $S_\tau^2=10/19$, so Proposition~\ref{gr:prop:fpc} gives, with $n_1=n_0=10$,
\[
\mathrm{Var}(\hat\tau)=\frac{5/19}{10}+\frac{5/19}{10}-\frac{10/19}{20}=\frac1{19}-\frac1{38}=\frac1{38}=0.02632,
\]
and the conservative value that drops $S_\tau^2/n$ is $1/19=0.05263$. The simulation draws $100{,}000$ complete randomizations of the 20 units, $10$ treated each time, with seed $20261019$. The mean of $\hat\tau$ is $0.0001$ and its empirical variance is $0.02633$, with Monte Carlo standard error $0.00011$. The simulation matches Proposition~\ref{gr:prop:fpc}, and the conservative variance overstates it by the factor two. An interval built on it is $\sqrt2$ times too wide.

\paragraph{Exercise~\ref{ex:gr:4}.} (a) Here $\tau_{\mathrm{fp}}=p_1-p_0=0.3$ and Proposition~\ref{gr:prop:frechet} gives $\max(0,0.3)\le p_{10}\le\min(0.8,0.5)$, that is $0.3\le p_{10}\le0.5$. The other cells follow from the margins, $p_{11}=p_1-p_{10}$, $p_{01}=p_0-p_{11}$ and $p_{00}=1-p_{11}-p_{10}-p_{01}$. In the order $(p_{11},p_{10},p_{01},p_{00})$ the two ends are
\[
(0.5,\ 0.3,\ 0,\ 0.2)\qquad\text{and}\qquad(0.3,\ 0.5,\ 0.2,\ 0).
\]
At the first end nobody is harmed, and at the second nobody fails under both arms. (b) Proposition~\ref{gr:prop:sharpvar} gives $0.3-0.09=0.21\le V_\tau\le2\times0.5-0.3-0.09=0.61$. Directly, $V_\tau=\E\tau_i^2-\tau_{\mathrm{fp}}^2=p_{10}+p_{01}-\tau_{\mathrm{fp}}^2$, which is $0.3+0-0.09=0.21$ at the first end and $0.5+0.2-0.09=0.61$ at the second. Because $p_{01}=p_{10}-\tau_{\mathrm{fp}}$ along the segment, $V_\tau=2p_{10}-\tau_{\mathrm{fp}}-\tau_{\mathrm{fp}}^2$ is linear and increasing in $p_{10}$, so the two ends give the extremes and both bounds are attained. (c) The overstatement is $S_\tau^2/n=V_\tau/(n-1)$, which for $n=100$ runs from $0.21/99=0.0021$ to $0.61/99=0.0062$, whatever the arm sizes. With $50$ units in each arm the conservative variance is $(S_1^2+S_0^2)/50=0.00828$, where $S_1^2=100\times0.16/99$ and $S_0^2=100\times0.25/99$. The true variance of $\hat\tau$ then lies between $0.0021$ and $0.0062$, and the conservative value is between $1.3$ and $3.9$ times too large.

\paragraph{Exercise~\ref{ex:gr:5}.} (a) Fix margins $p_1$ and $p_0$ in the open unit interval and index the joint laws by $s=p_{10}$, so that $p_{11}=p_1-s$, $p_{01}=p_0-p_1+s$ and $p_{00}=1-p_0-s$, with $s$ in the Fr\'echet segment of Proposition~\ref{gr:prop:frechet}, which has positive length. Then
\[
\begin{aligned}
\int f\,dF&=f(1,1)p_1+f(0,1)(p_0-p_1)+f(0,0)(1-p_0)\\
&\quad+s\,\{f(1,0)+f(0,1)-f(1,1)-f(0,0)\}.
\end{aligned}
\]
The integral is constant on the segment exactly when the bracket vanishes, which is the stated condition, and by Theorem~\ref{gr:thm:sep} this is separability. Directly, $f=\alpha(a)+\beta(b)$ makes both sums equal to $\alpha(0)+\alpha(1)+\beta(0)+\beta(1)$. Conversely $\alpha(a)=f(a,0)$ and $\beta(b)=f(0,b)-f(0,0)$ reproduce $f$ at $(0,0)$, $(1,0)$ and $(0,1)$, and by the condition at $(1,1)$. (b) With margins $p_1=p_0=\tfrac12$ take the comonotone law, $p_{11}=p_{00}=\tfrac12$, and the countermonotone law, $p_{10}=p_{01}=\tfrac12$. For $f=(a-b)^2$ the integrals are $0$ and $1$, and for $f=\one\{a=1,b=0\}$ they are $0$ and $\tfrac12$, so neither integral is constant. (c) For $f(a,b)=\kappa(a-b)$ the two sums of (a) are $2\kappa(0)$ and $\kappa(1)+\kappa(-1)$, so, with the margins in the open unit interval as in (a), the mean is the same for every joint law with those margins only if $\kappa(1)+\kappa(-1)=2\kappa(0)$. Then $\kappa(-1)-\kappa(0)=-\{\kappa(1)-\kappa(0)\}$, and with $p_{10}-p_{01}=p_1-p_0=\tau$, the ATE,
\[
\begin{aligned}
\E\,\kappa(a-b)&=\kappa(0)+\{\kappa(1)-\kappa(0)\}p_{10}+\{\kappa(-1)-\kappa(0)\}p_{01}\\
&=\kappa(0)+\{\kappa(1)-\kappa(0)\}\,\tau,
\end{aligned}
\]
an affine function of the ATE. The choice $\kappa(u)=u$ returns the ATE itself, and $\kappa(u)=u^2$ fails the condition.

\paragraph{Exercise~\ref{ex:gr:6}.} (a) Both stratum effects are $0.95-0.90=0.55-0.50=0.05$ and the strata have equal size, so $\tau_{\mathrm{fp}}=0.05$. Within a stratum the treated units are a uniform subset, so $\E\bar Y_{x,t}$ is the stratum mean of $Y(t)$. The treated arm has $n_1=200$ units, with weights $50/200=0.25$ and $150/200=0.75$ on the two strata, and the control arm has $n_0=200$ units with weights $0.75$ and $0.25$. Hence
\[
\begin{aligned}
\E(\text{collapsed})&=0.25\times0.95+0.75\times0.55-(0.75\times0.90+0.25\times0.50)\\
&=0.65-0.80=-0.15,
\end{aligned}
\]
which has the opposite sign to $\tau_{\mathrm{fp}}$: the bias is $-0.20$ although assignment is random within each stratum. (b) Write $\mu_{x,t}=n_x^{-1}\sum_{i\in x}Y_i(t)$, so that $\tau_{\mathrm{fp}}=\sum_x(n_x/n)(\mu_{x,1}-\mu_{x,0})$, and let $a_x=n_{x,1}/n_1-n_x/n$ and $b_x=n_{x,0}/n_0-n_x/n$, each of which sums to zero over $x$. Subtracting $\tau_{\mathrm{fp}}$ from $\E(\text{collapsed})=\sum_x(n_{x,1}/n_1)\mu_{x,1}-\sum_x(n_{x,0}/n_0)\mu_{x,0}$ gives
\[
\E(\text{collapsed})-\tau_{\mathrm{fp}}=\sum_xa_x\mu_{x,1}-\sum_xb_x\mu_{x,0}.
\]
If every $a_x$ and $b_x$ is zero the bias vanishes for every urn. Conversely, take the urn with $Y_i(1)=1$ for the units of one stratum $x_0$, $Y_i(1)=0$ elsewhere and $Y_i(0)=0$ throughout: the bias is $a_{x_0}$, and the mirror urn gives $b_{x_0}$. So a bias that vanishes for every urn forces $n_{x,1}/n_1=n_{x,0}/n_0=n_x/n$ for all $x$, which is proportional allocation in both arms. In the example $a=(-0.25,0.25)$ and $b=(0.25,-0.25)$, and the formula gives $-0.10-0.10=-0.20$. Equal case mix in the two arms, $a_x=b_x$, is not enough, since the contrast then estimates a different weighted average of the stratum effects.

\paragraph{Exercise~\ref{ex:gr:7}.} (a) In nats, with $\KL(q\,\|\,p)=q\log(q/p)+(1-q)\log\{(1-q)/(1-p)\}$,
\[
\KL(0.45\,\|\,0.30)=0.45\log1.5+0.55\log\tfrac{0.55}{0.70}=0.1825-0.1326=0.0498,
\]
\[
\KL(0.30\,\|\,0.45)=0.30\log\tfrac{0.30}{0.45}+0.70\log\tfrac{0.70}{0.55}=-0.1216+0.1688=0.0472.
\]
(b) Both rates are computed exactly, from the binomial tail $\Prob\{\mathrm{Bin}(n,0.30)\ge0.45n\}$ and from the tail of the Beta$(0.3n+1,0.7n+1)$ posterior beyond $0.45$, so there is no simulation and no seed.
\[
\begin{array}{lcccc}
n & 100 & 400 & 1600 & \text{limit}\\
\text{Sanov rate} & 0.0683 & 0.0561 & 0.0518 & 0.0498\\
\text{posterior rate} & 0.0674 & 0.0538 & 0.0493 & 0.0472
\end{array}
\]
Both rates approach their limits from above. The excess is $0.018$, $0.006$, $0.002$ for the Sanov rate and $0.020$, $0.007$, $0.002$ for the posterior rate, falling about threefold when $n$ quadruples. These are of order $(\log n)/n$, close to $(\log n)/(2n)=0.023$, $0.0075$, $0.0023$, as a tail prefactor of order $n^{-1/2}$ gives. The Sanov limit belongs to Theorem~\ref{gr:thm:sanov}, $\inf_{q\ge0.45}\KL(q\,\|\,0.30)=\KL(0.45\,\|\,0.30)$; the posterior limit belongs to Theorem~\ref{gr:thm:isanov}, $\inf_{p\ge0.45}\KL(0.30\,\|\,p)=\KL(0.30\,\|\,0.45)$. Each divergence increases beyond $0.45$, placing the infimum at the boundary. The limits differ because KL is asymmetric: the empirical law comes first in Sanov's exponent, the data law in the inverse.

\paragraph{Exercise~\ref{ex:gr:8}.} (a) At $n=700$, $\tfrac12d_G\log n=1.5\times6.551=9.827$ for each group and $2\times6.551=13.102$ for the heterogeneous model, whose gap is zero. With $D_G=1.1\times10^{-5}$, $7.3\times10^{-5}$ and $7.2\times10^{-4}$ the products $nD_G$ are $0.0077$, $0.0511$ and $0.504$. Relative to the risk ratio the weights are therefore $e^{-0.0434}=0.957$ for the risk difference, $e^{-0.496}=0.609$ for the odds ratio and $e^{-3.268}=0.038$ for heterogeneity. They sum to $2.604$ with the risk ratio's $1$, and the normalized weights are
\[
0.384,\quad0.368,\quad0.234,\quad0.015
\]
for the log risk ratio, the risk difference, the log odds ratio and the heterogeneous model. The text quotes $0.38$, $0.37$, $0.23$ and $0.01$, which agree, and the first split rests on only $n(D_{\mathrm{RD}}-D_{\mathrm{RR}})=0.043$ nats. Refitting the three stable models to the counts of Table~\ref{fw:tab:simpson} gives the gaps $1.120\times10^{-5}$, $7.311\times10^{-5}$ and $7.201\times10^{-4}$, and the same weights to four decimals. (b) Both groups have $d_G=3$, so the factor $n^{-d_G/2}$ is common to their weights and cancels from the odds,
\[
\frac{w_{\mathrm{RR}}}{w_{\mathrm{OR}}}=\exp\{n(D_{\mathrm{OR}}-D_{\mathrm{RR}})\}.
\]
Setting the odds equal to $19$ gives
\[
n=\frac{\log19}{7.2\times10^{-4}-1.1\times10^{-5}}=\frac{2.944}{7.09\times10^{-4}}\approx4{,}153,
\]
about six times the trial and the $4{,}150$ of the text. Polynomial factors matter only across models of different dimension. Against the heterogeneous model $d_{\mathrm{sat}}-d_G=1$ leaves the odds $n^{-1/2}e^{nD_{\mathrm{RR}}}$ for heterogeneity over the risk ratio, which reach one only at $n\approx593{,}000$, or $605{,}000$ with the rounded gap $1.1\times10^{-5}$.

\paragraph{Exercise~\ref{ex:gr:9}.} (a) With $g=(\sigma,\mu)$ and $g_0=(\sigma_0,\mu_0)$ the composition rule gives $gg_0=(\sigma_0\sigma,\ \mu+\sigma\mu_0)$ and $g_0g=(\sigma_0\sigma,\ \mu_0+\sigma_0\mu)$. The Jacobian matrices of $(\sigma,\mu)\mapsto gg_0$ and $(\sigma,\mu)\mapsto g_0g$ are
\[
\begin{pmatrix}\sigma_0&0\\\mu_0&1\end{pmatrix}\qquad\text{and}\qquad\begin{pmatrix}\sigma_0&0\\0&\sigma_0\end{pmatrix},
\]
with determinants $\sigma_0$ and $\sigma_0^2$. A density $\pi$ is invariant under a map $T$ when $\pi(Tg)\,|\det DT|=\pi(g)$. For $\pi=1/\sigma$ and $T:g\mapsto gg_0$ the left side is $\sigma_0/(\sigma_0\sigma)=1/\sigma$, and for $\pi=1/\sigma^2$ and $T:g\mapsto g_0g$ it is $\sigma_0^2/(\sigma_0\sigma)^2=1/\sigma^2$. So $d\sigma\,d\mu/\sigma$ is the right Haar measure and $d\sigma\,d\mu/\sigma^2$ the left. (b) The likelihood is $\sigma^{-n}\exp\{-A(\mu)/(2\sigma^2)\}$ with $A(\mu)=(n-1)s^2+n(\mu-\bar y)^2$, so under the prior $\sigma^{-k}d\sigma\,d\mu$ the posterior is proportional to $\sigma^{-(n+k)}e^{-A(\mu)/(2\sigma^2)}$. The substitution $u=A/(2\sigma^2)$ gives $\int_0^\infty\sigma^{-(n+k)}e^{-A/(2\sigma^2)}d\sigma=\tfrac12(A/2)^{-(n+k-1)/2}\,\Gamma\{(n+k-1)/2\}$, hence
\[
\pi(\mu\mid y)\propto A(\mu)^{-(n+k-1)/2}\propto\Bigl\{1+\frac{n(\mu-\bar y)^2}{(n-1)s^2}\Bigr\}^{-(\nu+1)/2},\qquad\nu=n+k-2.
\]
This is Student $t_\nu$ with scale $c=s\{(n-1)/(n(n+k-2))\}^{1/2}$. For $k=1$, $\nu=n-1$ and $c=s/\sqrt n$, giving the classical central $95$ percent interval $\bar y\pm t_{n-1}s/\sqrt n$. Its coverage is exact because $T_{n-1}=\sqrt n(\bar y-\mu)/s$ has the $t_{n-1}$ sampling law. For $k=2$, $\nu=n$, $c=s\sqrt{n-1}/n$, and the interval is $\bar y\pm t_ns\sqrt{n-1}/n$. It covers when $|T_{n-1}|\le t_n\sqrt{(n-1)/n}<t_{n-1}$, so it is shorter and has coverage below $0.95$ (Exercise~\ref{ex:gr:10}).

\paragraph{Exercise~\ref{ex:gr:10}.} The design is $n=3$ and $n=5$, with $100{,}000$ samples of $n$ independent $N(0,1)$ observations for each $n$ and seed $20261019$. Each sample gives the right Haar interval $\bar y\pm t_{n-1}s/\sqrt n$ and the left Haar interval $\bar y\pm t_ns\sqrt{n-1}/n$. The table gives the fraction of samples in which each contains $\mu=0$, the exact left Haar coverage $\Prob(|T_{n-1}|\le t_n\sqrt{(n-1)/n})$ and the ratio of the two interval widths.
\[
\begin{array}{cccccc}
n & t_{n-1},\ t_n & \text{right Haar} & \text{left Haar} & \text{exact left} & \text{width ratio}\\
3 & 4.303,\ 3.182 & 0.9487 & 0.8760 & 0.8783 & 0.604\\
5 & 2.776,\ 2.571 & 0.9497 & 0.9170 & 0.9170 & 0.828
\end{array}
\]
Monte Carlo standard errors are $0.0007$ for right Haar coverage and $0.001$ for left. Right Haar gives nominal $0.95$ coverage because $\sqrt n(\bar y-\mu)/s$ has the $t_{n-1}$ law (Fraser 1961; Hora and Buehler 1966). Left Haar uses $t_n$ and a scale smaller by $\sqrt{(n-1)/n}$, giving widths $60$ and $83$ percent as large at $n=3$ and $5$, and undercoverage. Simulations are within $0.003$ of exact coverage; $0.917$ at $n=5$ matches Section~\ref{gr:sec:transform}. The right Haar credible interval is also a confidence interval; the left is not.

\section*{Chapter~\ref{ch:dml}: What Double Machine Learning Assumes}

\paragraph{Exercise~\ref{ex:ch04:1}.} Under $Y = \tau(X)T + g(X) + \varepsilon$, partialling out gives $Y - \ell(X) = \tau(X)\nu + \varepsilon$ with $\nu = T - m(X)$. The population regression coefficient of $Y - \ell(X)$ on $\nu$ is $\E\{\tau(X)\nu^2\}/\E\nu^2 = \E\{\tau(X)\,\mathrm{Var}(T \mid X)\}/\E\,\mathrm{Var}(T \mid X)$. For binary $T$, $\mathrm{Var}(T \mid X) = e(X)\{1 - e(X)\}$, so the DML estimand is the overlap weighted average of conditional effects.

\paragraph{Exercise~\ref{ex:ch04:2}.} Perturb the nuisance functions to $\mu_1 + ra_1$, $\mu_0 + ra_0$ and $e + rb$, with bounded functions $a_1, a_0, b$ of $X$ and $e$ bounded away from 0 and 1, and differentiate the mean of the score \eqref{eq:aipw} at $r = 0$, with $\tau$ at its true value. Under the expectation the derivative is
\[
a_1 - \frac{T a_1}{e} - a_0 + \frac{(1 - T)a_0}{1 - e} - \frac{T\{Y - \mu_1\}}{e^2}\,b - \frac{(1 - T)\{Y - \mu_0\}}{(1 - e)^2}\,b.
\]
Given $X$, $\E(T\mid X)=e$ makes the first four terms cancel in pairs. The last two have mean zero since $\E\{T(Y-\mu_1)\mid X\}=\E\{(1-T)(Y-\mu_0)\mid X\}=0$. The derivative is zero in every direction $(a_1,a_0,b)$, so the score is Neyman orthogonal with respect to $\mu_1$, $\mu_0$ and $e$, separately and jointly. Exercise~\ref{ex:ch04:10} derives it from the exact remainder.

\paragraph{Exercise~\ref{ex:ch04:3}.} If $\mu_t$ is correct, $\E[T\{Y - \mu_1(X)\}/e(X) \mid X] = 0$ for any $e$, so the correction terms vanish in expectation and the score reduces to $\E\{\mu_1(X) - \mu_0(X)\} - \tau = 0$. If $e$ is correct but $\mu_t$ is not, $\E[T\{Y - \mu_1(X)\}/e(X) \mid X] = \E\{Y(1) \mid X\} - \mu_1(X)$, which cancels the error in $\mu_1(X)$, and similarly for $\mu_0$.

\paragraph{Exercise~\ref{ex:ch04:4}.} Let $\|\delta_m\| \asymp n^{-a}$ and $\|\delta_\ell\| \asymp n^{-b}$. Then $\|\delta_m\|(\|\delta_m\| + \|\delta_\ell\|) \asymp n^{-a - \min(a, b)}$, so condition (v) of Theorem~\ref{dml:thm:main} holds when $a > 1/4$ and $a + b > 1/2$. With $a = 1/2$ the second inequality asks only that $b > 0$, which is consistency, condition (iv). Take three covariates and a propensity that is a logistic function of a linear index, so that logistic regression gives $a = 1/2$, and an outcome regression $\ell$ that is Lipschitz in $X$ with no further structure. The minimax $L_2$ rate for Lipschitz functions of three variables is $n^{-1/(2+3)} = n^{-1/5}$, attained by local averaging, so $b = 1/5$ and $\|\delta_\ell\|$ is slower than $n^{-1/4}$. Yet $\|\delta_m\|(\|\delta_m\| + \|\delta_\ell\|) \asymp n^{-7/10} = o(n^{-1/2})$. At $n = 9{,}915$ the rates $n^{-1/5}$, $n^{-1/4}$ and $n^{-1/2}$ are 0.159, 0.100 and 0.0100, and the product is 0.0016. The extreme case is a randomized trial with known propensity, where $\delta_m = 0$ and any consistent $\hat\ell$ will do. The roles cannot be exchanged: $\|\delta_m\| \asymp n^{-1/5}$ gives $\|\delta_m\|^2 \asymp n^{-2/5} \gg n^{-1/2}$ whatever $\hat\ell$ is (Exercise~\ref{ex:ch04:8}).

\paragraph{Exercise~\ref{ex:ch04:5}.} We used one data set of the design, replication 1 of Table~\ref{tab:ric} (seed 20261014), and ran DML with the five fold partition shuffled by each of the seeds 1 to 10, with every lasso penalty chosen by five fold cross validation. The ten estimates, sorted, are 1.048, 1.059, 1.060, 1.062, 1.069, 1.070, 1.077, 1.094, 1.099 and 1.102. Their range is 0.054, their standard deviation 0.019 and their median 1.070; the sandwich standard errors run from 0.065 to 0.072, with mean 0.069. On the same data the coupled horseshoe, with the whole sample lasso fit $\hat m(X)$ as a covariate and flat priors on $\tau$ and on its coefficient, run for 1,500 sweeps with the first 500 discarded, has posterior mean 1.075, posterior standard deviation 0.078 and central 95 percent interval $(0.924, 1.223)$. Four further chains with other seeds give standard deviations from 0.076 to 0.081 and means from 1.065 to 1.079. The cross fitted coupling, on the partition with seed 1, gives mean 1.159 and standard deviation 0.068.

The spread over splits is a quarter of the posterior standard deviation, and this data set is typical. In replications 2 to 20, with the same ten seeds, the standard deviation over splits has median 0.022 and maximum 0.029, the posterior standard deviation has median 0.075 and minimum 0.063, and the ratio of the two never exceeds 0.39. Over all 500 replications of the table the median posterior standard deviation is 0.076 and the median sandwich standard error is 0.069, close to the standard deviations of the estimates across data sets, 0.081 for the coupled posterior mean and 0.078 for DML.

The spread over splits is noise from an arbitrary choice by the analyst, with the data held fixed. The posterior standard deviation is uncertainty about $\tau$ given the data, and its DML counterpart is the sandwich standard error, 0.069. The Bayesian counterpart of the split spread is the Monte Carlo error of the chain, here a standard deviation of 0.006 between chain means, a third of the split spread. More partitions stabilize the reported median without shrinking the spread of single splits; a longer chain shrinks the Monte Carlo error itself. On these ten splits the aggregation rule of Chapter~\ref{ch:dml} gives a standard error of 0.0715, about 4 percent above the median single split value of 0.0689. Split noise barely changes an interval's width, yet it can move a single reported estimate by 0.054, about 0.8 of a standard error, which is the reason to report the median over many partitions.

\paragraph{Exercise~\ref{ex:ch04:6}.} Write $D = \tau_\lambda - \tau$ and $d_j = b_j - 1$ for the five active coefficients. The other 45 stay at zero in the limit, because $X_j$ is independent of the residual and the subgradient condition $|\E(X_j R)| = 0 \leq \lambda$ holds. The residual at the lasso solution is $R = -DT - \sum_{j \leq 5} d_j X_j + \varepsilon$, and $\E T^2 = 6$, $\E(TX_j) = 1$ for $j \leq 5$. Because $T$ is not penalized, its score equation $\E(TR) = 0$ gives $\sum_j d_j = -6D$. For an active $X_j$ with $b_j > 0$ the subgradient condition is $\E(X_j R) = \lambda$, that is $-D - d_j = \lambda$. Hence $d_j = -D - \lambda$ for each of the five, so $\sum_j d_j = -5D - 5\lambda = -6D$ and
\[
D = 5\lambda, \qquad b_j = 1 - 6\lambda,
\]
which is valid while $1 - 6\lambda > 0$, that is for $\lambda \leq 1/6$. For $\lambda \geq 1/6$ all five coefficients are zero and $\tau_\lambda = \E(TY)/\E(T^2) = 11/6$, a bias of $5/6$, which is the value of $5\lambda$ at $\lambda = 1/6$. The lasso takes $6\lambda$ from each of the five confounder coefficients, and $\tau$ absorbs the signal they no longer carry through its correlation with $T$, which gives $D = 30\lambda/6$. If $T$ is penalized too, as in the cross validated fit behind Table~\ref{tab:ric}, the same steps give $D = 4\lambda$ for $\lambda \leq 1/5$.

\paragraph{Exercise~\ref{ex:ch04:7}.} Since $\hat V_{Y,i}-\tau_0\hat V_{T,i}=\varepsilon_i-\delta_{\ell,i}+\tau_0\delta_{m,i}$, any subset $I$ with positive denominator satisfies
\[
\hat\tau_I-\tau_0=\frac{\sum_{i\in I}\hat V_{T,i}(\hat V_{Y,i}-\tau_0\hat V_{T,i})}{\sum_{i\in I}\hat V_{T,i}^2}
=\frac{\sum_{i\in I}(s_i+c_i+b_i)}{\sum_{i\in I}\hat V_{T,i}^2},
\]
because $(\nu_i-\delta_{m,i})(\varepsilon_i-\delta_{\ell,i}+\tau_0\delta_{m,i})=s_i+c_i+b_i$ with $s_i=\nu_i\varepsilon_i$, $c_i=\tau_0\nu_i\delta_{m,i}-\nu_i\delta_{\ell,i}-\varepsilon_i\delta_{m,i}$ and $b_i=\delta_{m,i}\delta_{\ell,i}-\tau_0\delta_{m,i}^2$. Taking $I=\{1,\dots,n\}$ and dividing both sums by $n$ gives Lemma~\ref{dml:lem:exact}. Taking $I=I_k$ gives $\hat\tau_k-\tau_0=(S_k+C_k+B_k)/\hat J_k$, with $S_k,C_k,B_k$ the fold means and $\hat J_k=|I_k|^{-1}\sum_{i\in I_k}\hat V_{T,i}^2$. Put $N_k=\sum_{i\in I_k}\hat V_{T,i}\hat V_{Y,i}$ and $D_k=\sum_{i\in I_k}\hat V_{T,i}^2>0$. Then
\[
\hat\tau_{\mathrm{DML2}}=\frac{\sum_kN_k}{\sum_kD_k}=\sum_k\frac{D_k}{\sum_jD_j}\hat\tau_k.
\]
These weights also average the fold errors. DML1 uses $1/K$ instead.

\paragraph{Exercise~\ref{ex:ch04:8}.} With $\hat\ell = \ell$ we have $\delta_\ell = 0$, the cross channel $\langle\delta_m, \delta_\ell\rangle$ vanishes, and \eqref{eq:dml:twochannel} reduces to
\[
-\frac{\tau_0\|\delta_m\|^2}{J_0 + \|\delta_m\|^2}.
\]
The product $\|\delta_m\|\,\|\delta_\ell\|$ is zero, hence $o_p(n^{-1/2})$ whatever $\hat m$ is. Let $\tau_0 \neq 0$ and $\|\delta_m\| = c\,n^{-1/4}$ with $c > 0$. In Lemma~\ref{dml:lem:exact}, $\sqrt n\,S_n \Rightarrow N(0, \E\nu^2\varepsilon^2)$, $\sqrt n\,C_n \to 0$ as in the proof of Theorem~\ref{dml:thm:main}, $\hat J_n \to J_0$, and, if $\E\delta_m^4 \le C\|\delta_m\|^4$ (a fourth moment bound on the fit error, assumed here and not part of Corollary~\ref{dml:cor:channels}), the variance of $\sqrt n\,B_n$ is at most $C\tau_0^2\|\delta_m\|^4 \to 0$, so $\sqrt n\,B_n = -\tau_0\sqrt n\,\|\delta_m\|^2 + o_p(1) \to -\tau_0c^2$. So $\sqrt n(\hat\tau - \tau_0) \Rightarrow N(-\tau_0c^2/J_0, \sigma^2)$, a normal law with the wrong centre, and the interval undercovers. For $\|\delta_m\| \gg n^{-1/4}$ the variable diverges. The product bound misses the own square channel, which is why condition (v) of Theorem~\ref{dml:thm:main} carries $\|\delta_m\|^2$. The two channels cancel when $\langle\delta_m, \delta_\ell\rangle = \tau_0\|\delta_m\|^2$, that is, when $\langle\delta_m, \delta_\ell - \tau_0\delta_m\rangle = 0$: the error in the implied estimate of $g = \ell - \tau_0 m$ is orthogonal to $\delta_m$. The simplest case is $\hat\ell = \tau_0\hat m + g$, which gives $\delta_\ell = \tau_0\delta_m$ and an exact $\hat g$ (Proposition~\ref{dml:prop:orth}). In the design of Figure~\ref{fig:orthbias}, with $\delta_m = cx$, $\tau_0 = 1$ and $\delta_\ell = \kappa c x$, the bias is $(\kappa - 1)c^2/(1 + c^2)$. At $c = 0.2$ it is $-0.038$ for $\kappa = 0$ (exact $\hat\ell$), 0 for $\kappa = 1$ and $+0.038$ for $\kappa = 2$, the perturbation plotted in the figure.

\paragraph{Exercise~\ref{ex:ch04:9}.} The strata are $\pi_g = 0.06 + 0.44(g - 1)/7$ and $\tau_g = 2 - 2(g - 1)/7$, $g = 1, \dots, 8$. The weights $w_g = \pi_g(1 - \pi_g)$ are
\[
0.0564,\ 0.1078,\ 0.1512,\ 0.1868,\ 0.2144,\ 0.2342,\ 0.2460,\ 0.2500,
\]
with total 1.4469, and the weighted effects $w_g\tau_g$ are $0.1128$, $0.1847$, $0.2160$, $0.2135$, $0.1838$, $0.1338$, $0.0703$ and $0$, with total 1.1150. The strata are equal, so the common probability cancels and $\tau_* = 1.1150/1.4469 = 0.7706$. Over the eight strata the mean of $w_g$ is $\bar w = 0.1809$ and $\Cov(\tau_g, w_g) = \overline{w\tau} - \bar w\,\bar\tau = (1.1150 - 1.4469)/8 = -0.0415$, so $\Cov(\tau_g, w_g)/\bar w = -0.2294 = \tau_* - 1$. This is the covariance identity of Proposition~\ref{dml:prop:wate} at $\bar\tau = 1$, since $\tau_* = \overline{w\tau}/\bar w$ and $\Cov(\tau_g, w_g) = \overline{w\tau} - \bar w$. Reversing the order replaces $\tau_g$ by $2 - \tau_g$. The weighted effects now total 1.7787 (unrounded entries), and $\tau_*' = 2 - \tau_* = 1.2294$: the covariance changes sign, and the partially linear score overstates the average effect of 1.0 by 23 percent where it had understated it by 23 percent. The stratum with effect 2.0 now carries 17.3 percent of the weight and the stratum with no effect 3.9 percent.

\paragraph{Exercise~\ref{ex:ch04:10}.} Condition on $X$. Since $\E(TY \mid X) = e\mu_1$ and $\E(T \mid X) = e$, the second term of \eqref{eq:aipw} at $(\hat\mu_1, \hat e)$ has conditional mean $e(\mu_1 - \hat\mu_1)/\hat e$, and the third has conditional mean $(1 - e)(\mu_0 - \hat\mu_0)/(1 - \hat e)$. With $\tau = \E(\mu_1 - \mu_0)$ the mean of the score is
\[
\E\Bigl[(\hat\mu_1 - \mu_1)\Bigl(1 - \frac{e}{\hat e}\Bigr) - (\hat\mu_0 - \mu_0)\Bigl(1 - \frac{1 - e}{1 - \hat e}\Bigr)\Bigr],
\]
and $1 - e/\hat e = (\hat e - e)/\hat e$ and $1 - (1 - e)/(1 - \hat e) = -(\hat e - e)/(1 - \hat e)$ give the stated form. Along the path $(\mu_t + ra_t, e + rb)$ it equals $r^2\,\E[b\{a_1/(e + rb) + a_0/(1 - e - rb)\}]$, which is $O(r^2)$ when $e + rb$ stays in $[c, 1 - c]$ for some $c > 0$ and $a_1, a_0, b$ are square integrable. Its derivative at $r = 0$ is zero in every direction, which is Neyman orthogonality in $(\mu_1, \mu_0, e)$. The denominators are where overlap enters. If $\hat e \in [c, 1 - c]$, both $1/\hat e$ and $1/(1 - \hat e)$ are at most $1/c$, and the Cauchy and Schwarz inequality gives
\[
\bigl|\E\psi\bigr| \le \frac{1}{c}\,\|\hat e - e\|\bigl(\|\hat\mu_1 - \mu_1\| + \|\hat\mu_0 - \mu_0\|\bigr).
\]
Clipping $\hat e$ to $[c, 1 - c]$ enforces this bound by construction, at the price of a constant $1/c = 100$ for $c = 0.01$. It keeps the rate of $\hat e$ only if $e$ itself lies in $[c, 1 - c]$, because clipping to an interval that contains $e$ cannot increase $|\hat e - e|$. If $e < c$ on a set of positive probability, $\|\hat e - e\|$ stays above $\|(c - e)_+\|$ after clipping, the propensity factor does not shrink, and the product rate is lost.

\paragraph{Exercise~\ref{ex:ch04:11}.} The fold numerators and denominators are $(13.50,11.25)$ and $(1.25,0.75)$, giving estimates $1.2000$ and $1.6667$. DML1 is $1.4333$; DML2 is $14.75/12.00=1.2292$, with weights $0.9375$ and $0.0625$. Regrouping into $\{1,4,5\}$ and $\{2,3,6\}$ gives sums $(1.50,1.50)$ and $(13.25,10.50)$, estimates $1.0000$ and $1.2619$, and DML1 $1.1310$. DML2 remains $1.2292$, now with weights $0.125$ and $0.875$, because the two pooled sums are unchanged. DML1 gives each fold half the weight: a fold holding 6 percent of the denominator pulls it up with ratio $1.667$ in the first grouping; one holding 12.5 percent pulls it down with ratio $1.0000$ in the second. Real repartitioning also changes the residuals; holding them fixed isolates the combination rule.

\section*{Chapter~\ref{ch:ipw}: The Inverse Weighting Lineage}

\paragraph{Exercise~\ref{ex:ch05:1}.} The design draws exactly one elephant, since $99/100 + 49/4900 = 1$, and the Horvitz Thompson estimate of the total is $y_I/\pi_I$ for the elephant $I$ drawn. Let $\bar y$ be the average weight, so Sambo weighs $\bar y$ and the total is $50\bar y$. That forces the other 49 elephants to average $\bar y$ too; write $\sigma^2 = \frac{1}{49}\sum_j (y_j - \bar y)^2$ for their variance, the sum running over the other 49. The estimate is $\bar y/0.99 \approx 1.01\,\bar y$ with probability $0.99$ and $4900\,y_j$ with probability $1/4900$ for each other elephant $j$. It is unbiased, and its variance is $\sum_i y_i^2/\pi_i - (50\bar y)^2$:
\[
\begin{aligned}
\mathrm{Var}(\hat Y_{\mathrm{HT}}) &= \frac{\bar y^2}{0.99} + 4900 \cdot 49\,(\sigma^2 + \bar y^2) - 2500\,\bar y^2 \\
&= 240{,}100\,\sigma^2 + \Bigl(237{,}600 + \frac{100}{99}\Bigr)\bar y^2.
\end{aligned}
\]
The second term does not involve $\sigma$. Even if all elephants weigh the same the standard deviation is $487\,\bar y$, almost ten times the total being estimated. With $\bar y = 5$ tonnes and $\sigma = 0.5$ tonnes it is about 2,450 tonnes against a total of 250.

\paragraph{Exercise~\ref{ex:ch05:2}.} Take the treated outcomes independent given $X$, with a common mean and a variance proportional to the propensity,
\[
Y_i \mid X_i, T_i = 1 \sim N\{\mu_1,\ \sigma^2 e(X_i)\},
\]
with $\sigma^2$ known and a flat prior on $\mu_1$. By ignorability this is the likelihood of the treated units, and the assignment mechanism does not involve $\mu_1$. The posterior is normal with precision $\sum_i T_i/\{\sigma^2 e(X_i)\}$ and mean the precision weighted average of the treated outcomes,
\[
\E(\mu_1 \mid \text{data}) = \frac{\sum_i T_iY_i/e(X_i)}{\sum_i T_i/e(X_i)} = \hat\mu_{1,\mathrm{H}}.
\]
With variance $\sigma^2e(x)$, a unit at $e(x)=0.01$ counts as 100 observations of variance $\sigma^2$, as if its outcome averaged $1/e(x)$ units. Since $\E\sum_iT_i/e(X_i)=n$, the posterior variance $\sigma^2/\sum_iT_i/e(X_i)$ is about $\sigma^2/n$. That is hard to defend for the units least likely to be treated. Its common mean makes H\'ajek insensitive to outcome level; its variance model makes it trust extreme weights.

\paragraph{Exercise~\ref{ex:ch05:3}.} With $e(x) = \expit(sx)$, $1/e(x) = 1 + e^{-sx}$, and for $X \sim N(0, 1)$, $\E e^{-sX} = e^{s^2/2}$. The per observation variance of the treated term of \eqref{eq:ipw} is $\E\{Y(1)^2/e(X)\} - \{\E Y(1)\}^2$, which grows at least like $e^{s^2/2}$ times the second moment of $Y(1)$ in the region of small propensity. For the exact variance let $W = TY/e(X) - (1 - T)Y/\{1 - e(X)\}$, so that $T(1 - T) = 0$ gives $\E(W^2 \mid X) = \E\{Y(1)^2 \mid X\}/e(X) + \E\{Y(0)^2 \mid X\}/\{1 - e(X)\}$. In the design $Y(1) = X + 1 + \varepsilon$ and $Y(0) = X + \varepsilon$ with $\varepsilon \sim N(0, 1)$, so the two conditional second moments are $(X + 1)^2 + 1$ and $X^2 + 1$. With $1/e = 1 + e^{-sX}$, $1/(1 - e) = 1 + e^{sX}$ and the Gaussian moments $\E e^{\pm sX} = e^{s^2/2}$, $\E Xe^{-sX} = -se^{s^2/2}$ and $\E X^2e^{\pm sX} = (1 + s^2)e^{s^2/2}$, we get $\E W^2 = 5 + e^{s^2/2}(2s^2 - 2s + 5)$. Since $\E W = 1$, $\mathrm{Var}(W) = 4 + e^{s^2/2}(2s^2 - 2s + 5)$, and the variance of \eqref{eq:ipw} is $\mathrm{Var}(W)/n$. For $n = 1000$ the standard deviation at $s = 0.5, 1, \ldots, 4$ is 0.095, 0.111, 0.155, 0.266, 0.537, 1.24, 3.21 and 9.30, against 0.091, 0.113, 0.145, 0.208, 0.654, 0.581, 0.854 and 0.539 for the 300 replications of the figure: they agree up to $s = 2.5$ and fall two, four and seventeen times below the exact value at $s = 3$, 3.5 and 4.

\paragraph{Exercise~\ref{ex:ch05:4}.} Let $Q$ be the law of $X$. By Exercise~\ref{ex:ch04:1} the DML estimand is $\E\{\nu^2\tau(X)\}/\E\nu^2$ with $\nu = T - m(X)$. For binary $T$, $m(X) = \E(T \mid X) = e(X)$ and $\E(\nu^2 \mid X) = \mathrm{Var}(T \mid X) = e(X)\{1 - e(X)\}$. Conditioning on $X$ in numerator and denominator,
\[
\frac{\E\{\nu^2\tau(X)\}}{\E\nu^2} = \frac{\int \tau(x)\,e(x)\{1 - e(x)\}\,Q(dx)}{\int e(x)\{1 - e(x)\}\,Q(dx)} = \tau_{Q_e},
\]
where $Q_e(dx) \propto e(x)\{1 - e(x)\}\,Q(dx)$ is the overlap weighted law. Partialling out the treatment weights each covariate value by the variance of its treatment, which is largest at $e = 1/2$ and zero at $e \in \{0, 1\}$. As a check, take two equally likely strata with propensities 0.1 and 0.5 and effects 1 and 3. The average effect is 2, the overlap weights are proportional to 0.09 and 0.25, and the DML estimand is $(0.09 \cdot 1 + 0.25 \cdot 3)/0.34 = 2.47$.

\paragraph{Exercise~\ref{ex:ch05:5}.} Take two strata of the covariate. Stratum A holds 30 percent of the population, with propensity $e = 0.02$ and effect $\tau = -4$. Stratum B holds 70 percent, with $e = 0.5$ and $\tau = 1$. Overlap holds, and
\[
\tau_Q = 0.3 \cdot (-4) + 0.7 \cdot 1 = -0.5.
\]
At $\alpha = 0.05$ stratum A lies outside $[\alpha, 1 - \alpha]$, so $\mathcal{X}_\alpha$ is stratum B and $\tau_\alpha = 1$. The sign has changed, and an inverse weighted estimate with the true propensities follows it: it estimates $-0.5$ on the full sample and $1$ on the trimmed one. Trimming here answers a different question, the effect in stratum B. In general $\tau_\alpha$ and $\tau_Q$ differ in sign when the integrals of $\tau$ against $Q$ over the retained and the trimmed regions have opposite signs and the trimmed one is larger in absolute value. The example also shows why trimming tempts: of 1,000 units about 300 are in stratum A and about six of those are treated, so the negative effect rests on very little data.

\paragraph{Exercise~\ref{ex:ch05:6}.} With selection as treatment and $\pi_i$ as the propensity, the augmented estimator of the total is
\[
\hat Y_{\mathrm{AUG}} = \sum_{i=1}^{50}\hat y_i + \frac{y_I - \hat y_I}{\pi_I}
\]
for the elephant $I$ drawn, where $\hat y_i$ is the outcome model. Take the owner's model, $\hat y_i = m$ for every elephant, with $m$ fixed before the draw (Sambo's last recorded weight, say), and suppose the guess is right, $m = \bar y$ in the notation of Exercise~\ref{ex:ch05:1}. If Sambo is drawn the residual is zero and the estimate is $50\bar y$, the owner's answer, with no weight in sight. For the sample in hand, which was Basu's complaint about Horvitz and Thompson, it is not absurd. Its design variance is $\sum_i r_i^2/\pi_i - (\sum_i r_i)^2$ with residuals $r_i = y_i - \bar y$. Sambo's residual is zero and the others sum to zero, so
\[
\mathrm{Var}(\hat Y_{\mathrm{AUG}}) = 4900 \cdot 49\,\sigma^2 = 240{,}100\,\sigma^2,
\]
the Horvitz Thompson variance without the term $237{,}601\,\bar y^2$ that comes from the level. The weight is untouched. If elephant $j$ is drawn the estimate is $50\bar y + 4900(y_j - \bar y)$, so an elephant 20 percent heavier than average moves it by $980\,\bar y$, nearly twenty times the total. With $\sigma = 0.1\,\bar y$ the standard deviation is $49\,\bar y$ against $490\,\bar y$ for Horvitz and Thompson: ten times smaller and still as large as the total. The criticism therefore fails against what the model absorbs, the level, and holds against the correction, which multiplies a residual by $1/\pi_i$ and is guaranteed only on average over the design. A guess that misses by $\bar y - m$ adds about $237{,}601\,(\bar y - m)^2$ to the variance, the level term of Horvitz and Thompson, who guess $m = 0$. The augmented estimator behaves in the circus to the extent that the owner's model is right, which is the chapter's point: the model does the work and the weights only correct it.

\paragraph{Exercise~\ref{ex:ch06:1}.} Take $n \to \infty$ with $n/K \to 0$, the regime of the proposition of Section~\ref{sec:rrw}; for a fixed number of cells every cell is visited infinitely often and the posterior of $\psi$ is consistent. Visited cells then contribute at most $n/K$ to $\psi$, as in the proof of that proposition. Under the exchangeable hierarchical prior, the posterior for the common mean is driven by the selected units, which are drawn with probability proportional to $e(x)$. Since almost every visited cell holds one draw, that draw is Bernoulli with mean $\E\{e(X)\vartheta(X)\}/\E\,e(X)$, and the posterior of the common mean concentrates there. The posterior for each unvisited cell is centred at the posterior of the common mean, so the posterior of $\psi$ converges to this selected unit mean rather than to $\E\,\vartheta(X)$.

\paragraph{Exercise~\ref{ex:ch06:2}.} Every posterior mean of $\psi = K^{-1}\sum_x\vartheta(x)$ lies in $[0, 1]$, whatever the prior, because $\vartheta(x) \in [0, 1]$. The Horvitz Thompson estimate does not. Take the data set in which all $n$ units are selected and every $Y_i = 1$. Then
\[
\hat\psi_{\mathrm{HT}} = \frac{1}{n}\sum_{i=1}^n \frac{1}{e(X_i)} > 1,
\]
since $e(X_i) \leq 1 - \delta < 1$. A prior under which this data set is possible has a posterior mean of at most 1 there and cannot reproduce $\hat\psi_{\mathrm{HT}}$; a prior under which it is impossible has no posterior there, and fails the requirement as well.  Horvitz and Thompson ignore the range of $\psi$, which any posterior mean respects.

\paragraph{Exercise~\ref{ex:ch06:3}.} Cut the range of $\logit e$ into $J$ segments $I_1, \ldots, I_J$ holding shares $W_1, \ldots, W_J$ of the cells, let $B_1, \ldots, B_J$ be a B spline basis on these knots, and replace the line in \eqref{eq:coupled} by the expansion
\[
\logit\vartheta(x) = a + \sum_{j=1}^J \beta_j B_j\{\logit e(x)\} + \eta(x), \qquad \beta_j \mid \lambda_j, \gscale \sim N(0, \lambda_j^2\gscale^2),
\]
with $\eta(x) \sim N(0, \sigma_\eta^2)$ and half Cauchy $\lambda_j$. Take degree zero first, $B_j = \mathbf 1_{I_j}$, Sims's step function, where the algebra is exact. Let $M_j$ and $S_j$ count the selected units and successes in segment $j$, and write $q_j = \E\,\expit(a + \beta_j + \sigma_\eta Z)$ with $Z \sim N(0, 1)$.

Suppose first that no cell is selected twice, which holds with probability tending to one when $n^2/K \to 0$. Integrating out $\eta$, each selected unit in segment $j$ is Bernoulli with probability $q_j$, so the likelihood is $\prod_j q_j^{S_j}(1 - q_j)^{M_j - S_j}$. A cell with no selected unit has posterior mean $\E(q_j \mid \text{data})$, and the others, at most $n$ of them, move the average over cells by at most $n/K$. Hence
\[
\E(\psi \mid \text{data}) = \sum_{j=1}^J W_j\,\E(q_j \mid \text{data}) + O(n/K).
\]
Now let the global scale grow, with $a$ and $\sigma_\eta$ fixed. The substitution $w = \lambda_j\gscale$ gives the marginal prior of $\beta_j$ as
\[
p_{\gscale}(\beta) = \frac{2}{\pi\gscale}\int_0^\infty N(\beta; 0, w^2)\,\frac{\gscale^2}{\gscale^2 + w^2}\,dw.
\]
The integrand is negligible for $w \ll |\beta|$, close to $1/(\sqrt{2\pi}\,w)$ for $|\beta| \ll w \ll \gscale$, and negligible again for $w \gg \gscale$. So $\gscale\,p_{\gscale}(\beta)$ grows like $(2/\pi)(2\pi)^{-1/2}\log(\gscale/|\beta|)$, and $p_{\gscale}(\beta)/p_{\gscale}(\beta') = 1 + \log(|\beta'|/|\beta|)/\log\gscale$ to this order. The prior goes flat, but only at a logarithmic rate. The bound $\gscale\,p_{\gscale}(\beta)/\log\gscale \leq C\{1 + \log^+(1/|\beta|)\}$ for $\gscale \geq e$ allows dominated convergence when $0 < S_j < M_j$, since the likelihood then decays exponentially in $\beta_j$. The limit is a flat prior on $\ell_j = a + \beta_j$, and $a$ drops out. With $\sigma_\eta = 0$ we have $d\ell = dq/\{q(1 - q)\}$, the limiting posterior is $\mathrm{Beta}(S_j, M_j - S_j)$, and its mean is exactly $S_j/M_j$. With $\sigma_\eta > 0$ the posterior density of $q_j$ is proportional to $q^{S_j}(1 - q)^{M_j - S_j}/h(q)$ with $h(q) = dq/d\ell$ smooth and positive, and Laplace's method gives $S_j/M_j + O(1/M_j)$ for $S_j/M_j$ away from 0 and 1. If $S_j = 0$ or $S_j = M_j$ the posterior moves to $q_j = 0$ or 1, which is still $S_j/M_j$. If $M_j = 0$ the likelihood is flat and the prior mean remains, one half in the limit when $\sigma_\eta = 0$, and this is why every segment must contain selected units.

Quadrature on one segment shows the slow rate. With $M = 50$, $S = 15$, $\sigma_\eta = 0$, the posterior mean of $q$ at $\gscale = 1$, 10, $10^3$ and $10^6$ is 0.3168, 0.3058, 0.3017 and 0.3008 when $a = 1$, an effect large against its noise, and 0.2864, 0.2932, 0.2968 and 0.2982 when $a = -1$, an effect within its noise. With $\sigma_\eta = 1$ the flat limit is 0.3008 instead of 0.3000.

Together,
\[
\E(\psi \mid \text{data}) \longrightarrow \sum_{j=1}^J W_j\frac{S_j}{M_j} + R, \qquad R = \sum_{j=1}^J W_j\,O(1/M_j) + O(n/K).
\]
Within a segment $e$ is nearly constant, so the H\'ajek weights cancel and $\sum_jW_jS_j/M_j$ is the stratified H\'ajek estimator with known shares. Replacing $W_j$ by sample shares $N_j/n$ changes it by $\sum_j(N_j/n-W_j)S_j/M_j=O_p(n^{-1/2})$. Since $M_j\approx nW_j\bar e_j$, where $\bar e_j$ is the segment's mean propensity, $\sum_jW_j/M_j\approx n^{-1}\sum_j1/\bar e_j=O(J/n)$ under $e\ge\delta$. Thus the posterior mean approaches the stratified estimator if $J/n\to0$, with a gap negligible against its $n^{-1/2}$ sampling error when $J=o(n^{1/2})$.

The large global scale is an assumption that the hyperprior need not deliver. If the curve is flat, so that every $\beta_j$ is within its noise, a half Cauchy $\gscale$ falls with the data, the segments pool, and the posterior mean moves away from the stratified estimator toward an exchangeable one such as Li's. That is the better answer there, since the selected units are then representative. For degree above zero and $\sigma_\eta = 0$ the score equations of the limit, $\sum_i R_iB_j(\logit e_i)\{Y_i - \hat q(e_i)\} = 0$, balance outcomes over overlapping windows rather than exact strata, so the limit is a smoothed version of the stratified estimator. With $\sigma_\eta > 0$ each term also carries the factor $\hat h_i/\{\hat q_i(1 - \hat q_i)\}$, with $h = dq/d\ell$ as above, which reweights the windows.

\paragraph{Exercise~\ref{ex:ch06:4}.} Take the regularization induced confounding setup, $Y = \tau T + X^\top\beta + \varepsilon$ and $T = X^\top\gamma + \nu$, with $\varepsilon \sim N(0, \sigma^2)$, a flat prior on $\tau$ and independent $N(0, v)$ priors on the $\beta_j$, and condition on $\sigma^2$. For $n$ observations let $Y$, $T$ and $X$ denote the data vectors and matrix. The posterior mean of $(\tau, \beta)$ minimizes $\|Y - \tau T - X\beta\|^2 + (\sigma^2/v)\|\beta\|^2$. Write $\hat\Sigma = X^\top X/n$, $\hat\sigma_{XT} = X^\top T/n$, $\hat\sigma_{TT} = T^\top T/n$ and $\zeta = \sigma^2/(nv)$. The normal equations are
\[
\hat\sigma_{TT}\,\hat\tau + \hat\sigma_{XT}^\top\hat\beta = \frac{T^\top Y}{n}, \qquad \hat\sigma_{XT}\,\hat\tau + (\hat\Sigma + \zeta I)\hat\beta = \frac{X^\top Y}{n}.
\]
They are linear in $Y$, so given $X$ and $T$ the mean of $(\hat\tau, \hat\beta)$ solves them with $Y$ replaced by $\tau T + X\beta$. The errors $b = \hat\tau - \tau$ and $d = \hat\beta - \beta$ then satisfy $\hat\sigma_{TT}\,b + \hat\sigma_{XT}^\top d = 0$ and $\hat\sigma_{XT}\,b + (\hat\Sigma + \zeta I)d = -\zeta\beta$. Eliminating $d$ gives the bias
\[
b = \frac{\zeta\,\hat\sigma_{XT}^\top(\hat\Sigma + \zeta I)^{-1}\beta}{\hat\sigma_{TT} - \hat\sigma_{XT}^\top(\hat\Sigma + \zeta I)^{-1}\hat\sigma_{XT}}.
\]
This is the conditional bias of the proposition of Appendix~\ref{app:theory}, written there with the projected design $X^\top MX$, $M = I - T(T^\top T)^{-1}T^\top$; the two forms agree by the Sherman--Morrison formula. The bias vanishes as $v \to \infty$ when $(T, X)$ has full column rank and tends, as $v \to 0$, to $\hat\sigma_{XT}^\top\beta/\hat\sigma_{TT}$, the omitted variable bias of Proposition~\ref{crit:prop:ric}. In the design of that proposition ($\Sigma = I$, and $\gamma_j = \beta_j = c$ on $s$ coordinates), replacing sample moments by population ones gives $\hat\sigma_{XT} = \gamma$ and
\[
b = \frac{sc^2\sigma^2}{sc^2\sigma^2 + (nv + \sigma^2)\,\sigma_\nu^2}.
\]
With $s = 5$ and $c = \sigma = \sigma_\nu = 1$ the ceiling is 0.83, and at $n = 300$ the bias is 0.56, 0.14 and 0.016 for $v = 0.01$, 0.1 and 1. For fixed $v$ the bias is of order $1/n$; it stays of order one when $nv$ stays bounded, that is, when the prior standard deviation shrinks like $n^{-1/2}$.

\paragraph{Exercise~\ref{ex:ch06:5}.} Let $\logit e_\gamma(x) = c(x)^\top\gamma$ be the propensity model. For $e$ fixed, the coupled posterior has hyperparameters $\xi$, and the outcome factor of the likelihood depends on $\gamma$ only through the prior \eqref{eq:coupled} with $e_\gamma$ in place of $e$. The joint posterior therefore lets the selected outcomes speak about $\gamma$, the feedback that Zigler and coauthors (2013) warn about. The modular recipe of the chapter cuts it: draw $\gamma^{(s)}$ from the posterior given $(X_i, R_i)$ alone, recompute $e_{\gamma^{(s)}}$, draw $\xi^{(s)}$ and the $\vartheta(x)$ given that propensity and the outcomes, and average over cells to get $\psi^{(s)}$. By the law of total variance the spread of these draws is
\[
\E_\gamma\{\mathrm{Var}(\psi \mid \gamma, \text{data})\} + \mathrm{Var}_\gamma\{\E(\psi \mid \gamma, \text{data})\},
\]
and a plug in analysis at $\hat e$ keeps the first term and drops the second. So the uncertainty enters as an average over refitted posteriors, and no term has to be derived.

How large is the second term? It depends on how the fit responds to the propensity. Take the flexible limit of Exercise~\ref{ex:ch06:3}, in which the posterior mean is the average over cells of the selected units' regression on the supplied propensity, and supply $e_\varepsilon = e + \varepsilon d$ with $d$ bounded. Writing $V = e_\varepsilon(X)$ and substituting $e = V - \varepsilon d$ gives the exact identity
\[
\E(Y \mid R = 1, V) = \frac{\E(e\vartheta \mid V)}{\E(e \mid V)} = \E(\vartheta \mid V) - \varepsilon\,\frac{\mathrm{Cov}(d, \vartheta \mid V)}{\E(e \mid V)}.
\]
Averaging over $X$, and assuming that these conditional moments given $V$ tend to those given $e(X)$ as $\varepsilon \to 0$, gives $\psi(\varepsilon) = \psi - \varepsilon\,\E[\mathrm{Cov}\{d, \vartheta \mid e(X)\}/e(X)] + o(\varepsilon)$. The assumption is needed: if $e$ is constant and $d = \vartheta$ then $V$ reveals $\vartheta$, and $\psi(\varepsilon) = \psi$ for every small $\varepsilon \ne 0$. The one arm version of the augmented score \eqref{eq:aipw}, $m(X) + R\{Y - m(X)\}/\hat e(X)$, has the exact mean
\[
\psi + \E\left[\frac{\{m(X) - \vartheta(X)\}\{\hat e(X) - e(X)\}}{\hat e(X)}\right].
\]
For $\hat e=e_\varepsilon$ and $m(X)=\E\{\vartheta(X)\mid e(X)\}$, the augmented estimator has the coupled posterior's first order sensitivity. Using $X$ lets $m$ approach $\vartheta$, making the error a product of two small errors, as in DML. Coupling through $e$ alone leaves the fixed residual $\vartheta-\E(\vartheta\mid e)$ and first order bias unless this residual is uncorrelated with the propensity error given $e$. Under \eqref{eq:coupled}, with $\eta$ independent of $x$, it is uncorrelated with any fixed $d$ and the first order term vanishes. Adding $X$ shrinks the residual and gives a product in the linear model of Proposition~\ref{ss:prop:product}. For selected outcomes, adding $X$ alone does not suffice: logistic regression balances selected residuals against its covariates, while the product needs balance against $1/e$. The augmented score supplies that balance, as would a fit including $1/e$.

The variances differ too. For a parametric $\hat\gamma$ let $\phi = m(X) + R\{Y - m(X)\}/e(X)$ have variance $V_\phi$, let $I = \E[e(1 - e)cc^\top]$ be the information of the propensity score $U = \{R - e(X)\}c(X)$, and let $A = \partial_\gamma\E\phi_\gamma = -\E[(\vartheta - m)(1 - e)c]$. Since $\mathrm{Cov}(\phi, U) = -A$, the estimator at $\hat\gamma$ has $n\,\mathrm{Var}\{\hat\psi(\hat\gamma)\} \approx V_\phi - A^\top I^{-1}A$, the familiar fact that estimating a known propensity can reduce the variance (Hirano, Imbens and Ridder 2003 give the semiparametric form). An analyst of the augmented estimator who treats $\hat e$ as known reports $V_\phi$, and the cut posterior reports $V_\phi + A^\top I^{-1}A$, which adds the term and ignores the negative covariance between the two sampling errors. When $A = 0$, because the outcome model is correct, all three agree to first order. So the augmented estimator takes the uncertainty in $\hat e$ through its sampling distribution, while the posterior carries it explicitly and is conservative in the frequentist sense when its outcome model leaves residual dependence on $c(X)$ given $e$. All of this needs a propensity model that can be estimated at all: if $e$ were as rough over the $K \gg n$ cells as $\vartheta$, neither approach would have a usable $\hat e$.

\paragraph{Exercise~\ref{ex:ch06:6}.} Index the $K = 20{,}000$ cells by $x_k = (k - 1/2)/K$ on $[0, 1]$, with $n = 2{,}000$ draws. Take
\[
\logit\vartheta(x) = -1 + 3(2x - 1) + \sin(2\pi x), \qquad \logit e(x) = 3(2x - 1) + \epsilon(x),
\]
with $\epsilon(x) \sim N(0, 1)$ independently over cells and $e$ clipped to $[0.05, 0.95]$. The outcome surface is smooth in $x$, $\psi = 0.314$, and $\vartheta$ and $e$ have correlation 0.78 across cells, so the selected units overstate $\psi$ just as in the chapter's example. The prior ignores $e$: $\logit\vartheta(x) = \sum_{j=1}^{13}\beta_jB_j(x)$ with cubic B splines on equally spaced knots and independent $\beta_j \sim N(0, 100)$. Then $\psi = K^{-1}\sum_x\expit\{\sum_j\beta_jB_j(x)\}$. We computed its posterior by importance sampling, 1,500 draws from a multivariate $t$ with six degrees of freedom centred at the posterior mode, with the average over cells taken on every fortieth cell. The median effective sample size was 323.

Over 1,000 replications (seed 606, $e$ redrawn each time), the smooth prior had bias 0.002 (Monte Carlo standard error 0.0004), root mean squared error 0.014, 95\% interval coverage 0.95 and mean width 0.055. Horvitz and Thompson had bias 0.000 and root mean squared error 0.019. Independent uniform priors on the cells had bias and root mean squared error 0.185, and Li's exchangeable estimator had 0.109 and 0.110. The smooth prior beats the weights here without using $e$ at all. The reason is that $\E(Y \mid R = 1, X = x) = \vartheta(x)$ whatever $e$ is, and $e \geq 0.05$ everywhere, so the selected units cover the whole of $[0, 1]$ and a low dimensional regression recovers $\vartheta$ on the unvisited cells. Selection changes where the data are dense, which does not bias the fitted curve.

The control is the chapter's own design, in which $e$ is uniform on $[0.05, 0.95]$ independently over the cells and $\logit\vartheta = -1 + 6(e - 0.5) + \eta$, so that $\vartheta$ is rough in $x$. Here the same prior, over 200 replications (seed 606), has bias 0.136, root mean squared error 0.137 and coverage 0 with mean width 0.061. Its result equals Li's (bias 0.137, root mean squared error 0.137), since $x$ carries no information across cells and the curve fits the selected mean. Horvitz and Thompson have root mean squared error 0.014. This is the regime named in Section~\ref{sec:rrw}: the failure needs a $\vartheta$ so rough that $x$ says nothing across cells while $e$ does. One dimension is the favourable case, since in a high dimensional $x$ a sample of this size would leave few visited neighbours for a smooth prior to use.

\paragraph{Exercise~\ref{ex:chweak:1}.} In the cell model the $n$ units are independent: $X_k$ is uniform on the $B$ cells, $R_k \mid X_k = b$ is $\mathrm{Bernoulli}(p_b)$ and, independently of $R_k$, $Y_k \mid X_k = b$ is $\mathrm{Bernoulli}(\theta_b)$. Since $\E(R_k/p_{X_k}) = 1$, $\hat n/n \to 1$ in probability, and the proposition of Section~\ref{sec:weak} gives $\hat\psi_{\mathrm{H}} - \psi \approx n^{-1}\sum_k Z_k$ with $Z_k = R_k(Y_k - \psi)/p_{X_k}$. The $Z_k$ are independent and identically distributed with mean $B^{-1}\sum_b(\theta_b - \psi) = 0$. Because $R_k^2 = R_k$ and $\E R_k = p_b$ in cell $b$,
\[
\E(Z_k^2 \mid X_k = b) = \frac{p_b\,\E(Y - \psi)^2}{p_b^2} = \frac{\theta_b(1 - \theta_b) + (\theta_b - \psi)^2}{p_b},
\]
and averaging over cells gives $n\,\mathrm{Var}(\hat\psi_{\mathrm{H}}) \approx B^{-1}\sum_b\{\theta_b(1 - \theta_b) + (\theta_b - \psi)^2\}/p_b$, the formula of the critique. For Horvitz and Thompson, $W_k = R_kY_k/p_{X_k} - \psi$ has mean zero and, since $Y_k^2 = Y_k$, $\E W_k^2 = B^{-1}\sum_b\theta_b/p_b - \psi^2$, which is exact. The numerators differ by $\psi(\psi - 2\theta_b)$, so H\'ajek's is smaller exactly in the cells with $\theta_b > \psi/2$. In the first design of Table~\ref{tab:weak}, $\theta_b$ is independent of $p_b$ with mean $\psi = 0.5$, the mean of $1/p_b$ is about 4.7, and $\E\{\theta(1 - \theta) + (\theta - \psi)^2\} = \psi(1 - \psi)$. The formulas give $(0.5 \cdot 4.7 - 0.25)/100 \approx 0.021$ for Horvitz and Thompson and $0.25 \cdot 4.7/100 \approx 0.012$ for H\'ajek, which are 2.1 and 1.2 on the scale of the table. A simulation of this design at $n = 2000$ gives $n\,\mathrm{Var}$ of 2.11 and 1.15 against 2.13 and 1.18 from the formulas at the simulated $\theta_b$. At $n = 100$ the simulated H\'ajek variance is 0.96 on the scale of the table, below the first order 1.18, because $\hat n/n$ then has standard deviation 0.19 and second order terms matter.

\paragraph{Exercise~\ref{ex:chweak:2}.} Fix $\lambda$. The denominator $(1 - \lambda)n + \lambda\hat n$ has mean $n$ and $\hat S$ has mean $n\psi$, so expand $\hat S/\{(1 - \lambda)n + \lambda\hat n\}$ about $(n\psi, n)$, where the partial derivatives are $1/n$ and $-\lambda\psi/n$:
\[
\hat\psi_\lambda - \psi \approx \frac{1}{n}\bigl\{(\hat S - n\psi) - \lambda\psi(\hat n - n)\bigr\}.
\]
Every member of the family is consistent to first order, and
\[
\mathrm{Var}(\hat\psi_\lambda) \approx \frac{1}{n^2}\bigl\{\mathrm{Var}(\hat S) - 2\lambda\psi\,\mathrm{Cov}(\hat S, \hat n) + \lambda^2\psi^2\,\mathrm{Var}(\hat n)\bigr\}.
\]
The endpoints are Horvitz and Thompson ($\lambda = 0$) and, with $\hat S - \psi\hat n$ in the braces, H\'ajek ($\lambda = 1$), the proposition of Section~\ref{sec:weak}. The quadratic in $\lambda$ is minimized at
\[
\lambda^* = \frac{\mathrm{Cov}(\hat S, \hat n)}{\psi\,\mathrm{Var}(\hat n)}, \qquad \mathrm{Var}(\hat\psi_{\lambda^*}) \approx \frac{1}{n^2}\Bigl\{\mathrm{Var}(\hat S) - \frac{\mathrm{Cov}(\hat S, \hat n)^2}{\mathrm{Var}(\hat n)}\Bigr\},
\]
the regression coefficient of $\hat S$ on $\hat n$ divided by $\psi$. Since $\psi$ is unknown it has to be estimated, like the variances. H\'ajek is optimal exactly when $\hat S - \psi\hat n$ is uncorrelated with $\hat n$. In the cell model $\mathrm{Var}(\hat n) = n(\overline{1/p} - 1)$ and $\mathrm{Cov}(\hat S, \hat n) = n(\overline{\theta/p} - \psi)$, with bars for averages over cells, so $\lambda^* = 1$ when $\theta_b$ is uncorrelated with $1/p_b$, that is, when selection is unrelated to outcome. In the last design of Table~\ref{tab:weak}, where it is related, $\lambda^* \approx 0.8$, and the first order variances times $n$ are 0.94 for Horvitz and Thompson, 0.69 for H\'ajek and 0.67 at $\lambda^*$.

\paragraph{Exercise~\ref{ex:chweak:3}.} A fixed size design has $\sum_k R_k = n_R$ for every sample, with $n_R$ constant, and so $\sum_k p_k = n_R$. If $Y_k = c\,p_k$ then
\[
\hat S = \sum_k \frac{R_kY_k}{p_k} = c\sum_k R_k = c\,n_R
\]
for every sample, so $\hat\psi_{\mathrm{HT}} = c\,n_R/n$ is constant and equal to $\psi$: zero variance. H\'ajek is $c\,n_R/\hat n$ with $\hat n = \sum_k R_k/p_k$, which varies from sample to sample unless the $p_k$ are equal. With $p = (1/4, 3/4, 1)$ and the sample of size 2 that always contains the third unit and one of the first two, $\hat n$ is 5 or $7/3$ with probabilities $1/4$ and $3/4$, and the H\'ajek estimate is $0.40c$ or $0.86c$ against $\psi = 0.67c$. To first order $\hat\psi_{\mathrm{H}} - \psi \approx n^{-1}(\hat S - \psi\hat n) = \psi(1 - \hat n/n)$, with variance $\psi^2\mathrm{Var}(\hat n)/n^2 > 0$. Under the independent Bernoulli design $\sum_k R_k$ is random, so $\hat S = c\sum_k R_k$ varies and $\mathrm{Var}(\hat\psi_{\mathrm{HT}}) = c^2\sum_k p_k(1 - p_k)/n^2 > 0$. H\'ajek is not exact either, with first order variance $c^2n^{-2}\sum_k p_k(1 - p_k)(1 - \bar p/p_k)^2$ and $\bar p = n^{-1}\sum_k p_k$. The estimator that is exact for $Y_k = c\,p_k$ is $\hat S\sum_k p_k/(n\sum_k R_k)$, which divides by the realized sample size. Horvitz and Thompson are exact for outcomes proportional to $p_k$ under fixed sample size, which supplies the required count. Under Bernoulli sampling that model needs normalization by the realized count, a ratio estimator using $p_k$ as its size variable. H\'ajek's common mean model is exact for constant outcomes $\mu$ in every design, since $\hat S=\mu\hat n$.

\paragraph{Exercise~\ref{ex:chweak:4}.} In the cell model the pairs $(R_kY_k, R_k)$ are independent and identically distributed across draws, and both entries are Bernoulli. With $u = B^{-1}\sum_b\theta_bp_b = \psi\bar p + \sigma_{\theta,p}$ we have $\E R_kY_k = u$ and $\E R_k = \bar p$, and since $R_k^2 = R_k$, $\E(R_kY_k \cdot R_k) = u$. For the numerator $\mathcal{N} = \sum_k R_kY_k + 1$ and the denominator $\mathcal{D} = \sum_k R_k + 2$,
\[
\begin{gathered}
\mathrm{Var}(\mathcal{N}) = n\,u(1 - u), \qquad \mathrm{Var}(\mathcal{D}) = n\bar p(1 - \bar p), \\
\mathrm{Cov}(\mathcal{N}, \mathcal{D}) = n\,u(1 - \bar p),
\end{gathered}
\]
all linear in $n$; no $1/p_b$ appears. Since $0\le\hat\psi_{\mathrm{Li}}=\mathcal N/\mathcal D\le1$ and $\mathcal N/n\to u$, $\mathcal D/n\to\bar p$ almost surely, bounded convergence gives $\E\hat\psi_{\mathrm{Li}}\to\psi_\infty=u/\bar p=\psi+\sigma_{\theta,p}/\bar p$. With $\psi_n=\E\mathcal N/\E\mathcal D\to\psi_\infty$, the ratio delta method gives
\[
\begin{split}
\mathrm{Var}(\hat\psi_{\mathrm{Li}})&\approx\frac{n\{u(1-u)-2\psi_nu(1-\bar p)+\psi_n^2\bar p(1-\bar p)\}}{(n\bar p+2)^2}\\
&\approx\frac{u(\bar p-u)}{n\bar p^3}=\frac{\psi_\infty(1-\psi_\infty)}{n\bar p}.
\end{split}
\]
This is a proportion's variance for $n\bar p$ selected draws, or $V_{\mathrm{Li}}/n$ in Appendix~\ref{app:theory}. In Table~\ref{tab:weak}'s first design, $\bar p=0.5$, $u\approx0.25$, the limit gives 0.50 on the table's scale. Retaining pseudo counts, $\E\mathcal N=26$, $\E\mathcal D=52$, gives 0.46, close to 0.47.

\paragraph{Exercise~\ref{ex:chweak:5}.} The chapter does not fix the priors, so we take Sims's step function with these choices. Cut the range of $p$ into $J$ equal length segments. Within segment $j$ the selected units have outcome probability $\mu_j$, a draw lands in segment $j$ with probability $W_j$, and the priors are $\mu_j \sim \mathrm{Beta}(1, 1)$ and $W \sim \mathrm{Dirichlet}(1, \ldots, 1)$. With $N_j$ draws, $M_j$ selected units and $S_j$ successes in the segment, the likelihood factors as $\prod_jW_j^{N_j}\cdot\prod_j\mu_j^{S_j}(1 - \mu_j)^{M_j - S_j}$ times known terms in $p$. So $W \mid \text{data} \sim \mathrm{Dirichlet}(1 + N_1, \ldots, 1 + N_J)$, independently of $\mu_j \mid \text{data} \sim \mathrm{Beta}(1 + S_j, 1 + M_j - S_j)$. For $\psi = \sum_jW_j\mu_j$ the posterior mean is
\[
\hat\psi_{\mathrm{S}} = \sum_{j=1}^J\frac{1 + N_j}{J + n}\cdot\frac{1 + S_j}{2 + M_j}.
\]
Replacing the Dirichlet by the known shares of the cells gives the bin posterior of Table~\ref{tab:weak}. Binned H\'ajek is as in that table, $\sum_jS_j/\bar p_j$ over $\sum_jM_j/\bar p_j$, where $\bar p_j$ is the mean selection probability of the draws in the segment. It therefore estimates each segment's share through the noisy $M_j/\bar p_j$.

The design is the fifth row of Table~\ref{tab:weak}: $B = 1{,}000$ cells, $p_b$ equally spaced on $[0.1, 0.9]$, $\theta_b = 0.5p_b + 0.5U_b$ redrawn in each replication, $n = 100$ draws, 5,000 replications, seed 2204. For $J = 2$, 5, 10, 20 and 40, mean squared error times 100 is 0.61, 0.66, 0.67, 0.68 and 0.67 for binned H\'ajek and 0.50, 0.38, 0.27, 0.17 and 0.09 for the sieve, with Monte Carlo standard errors of at most 0.013 and paired standard errors of the differences of at most 0.012. With known shares the sieve gives 0.50, 0.38, 0.25, 0.14 and 0.06, so at most 0.026 of its error is the cost of the Dirichlet. At $J = 10$ the binned H\'ajek bias is 0.002 and the sieve bias is 0.015.

The sieve wins by a wide margin, but much of the win comes from the prior. At $J=10$, about five selected units per segment leave weight $2/(2+M_j)$ on the prior centre, one half. Since $\psi=0.5$, stronger shrinkage with larger $J$ removes noise with little bias. Halving the outcome probabilities to $\theta_b=0.25p_b+0.25U_b$, hence $\psi=0.25$, with the rest unchanged, gives binned H\'ajek 0.46, 0.48, 0.49, 0.49, 0.49 and sieve 0.47, 0.63, 1.07, 1.85, 2.85, with sieve bias 0.03, 0.06, 0.09, 0.13, 0.17. The order reverses by $J=5$, as in the table's off centre rows. A fair sieve should learn its prior centre, as in Exercise~\ref{ex:chweak:6}.

\paragraph{Exercise~\ref{ex:chweak:6}.} Let $M_b$ and $S_b$ count the selected draws and successes in cell $b$, let $\mathcal V$ be the set of cells with $M_b > 0$ and $m$ its size, and let $z_b = (1, \logit p_b)^\top$. The prior is the regression on selection of \eqref{eq:coupled} with horseshoe deviations:
\[
\logit\theta_b = z_b^\top\gamma + \eta_b, \qquad \eta_b \mid \lambda_b, \gscale \sim N(0, \lambda_b^2\gscale^2), \qquad \gamma \sim N(0, vI),
\]
with half Cauchy $\lambda_b$ and $\gscale$. As in Section~\ref{sec:hsregression}, write $\lambda_b^2 \mid \nu_b \sim \mathrm{IG}(1/2, 1/\nu_b)$, $\nu_b \sim \mathrm{IG}(1/2, 1)$, $\gscale^2 \mid \xi \sim \mathrm{IG}(1/2, 1/\xi)$ and $\xi \sim \mathrm{IG}(1/2, 1)$. The likelihood is $\prod_{b \in \mathcal V}\theta_b^{S_b}(1 - \theta_b)^{M_b - S_b}$, since the selection factor does not involve $\theta$. A cell outside $\mathcal V$ has no likelihood, so its $(\eta_b, \lambda_b, \nu_b)$ integrate out exactly, and the sampler runs on $\mathcal V$ alone.

Write $\zeta_b=z_b^\top\gamma+\eta_b$ and $\kappa_b=S_b-M_b/2$. The identity of Appendix~\ref{app:toolkit} gives
\[
\frac{e^{S_b\zeta_b}}{(1+e^{\zeta_b})^{M_b}}=2^{-M_b}e^{\kappa_b\zeta_b}\int_0^\infty e^{-\omega\zeta_b^2/2}p_{M_b,0}(\omega)\,d\omega,
\]
where $p_{M,0}$ is the $\mathrm{PG}(M,0)$ density. Add $\omega_b$ for each $b\in\mathcal V$. Exponential tilting gives $\omega_b\mid\zeta_b\sim\mathrm{PG}(M_b,\zeta_b)$. Given $\omega_b$, the Gaussian pseudo observation is $y_b=\kappa_b/\omega_b\sim N(\zeta_b,1/\omega_b)$. Integrating out $\eta_b$ gives $y_b\sim N(z_b^\top\gamma,d_b)$ with $d_b=1/\omega_b+\lambda_b^2\gscale^2$. One scan is then:
\begin{enumerate}
\item Draw $\omega_b \sim \mathrm{PG}(M_b, z_b^\top\gamma + \eta_b)$ for $b \in \mathcal V$.
\item Draw $\gamma \sim N(Cr, C)$ with $C^{-1} = v^{-1}I + \sum_{\mathcal V}z_bz_b^\top/d_b$ and $r = \sum_{\mathcal V}z_by_b/d_b$. This has $\eta$ integrated out, so steps 2 and 3 together are one joint draw of $(\gamma, \eta)$.
\item Draw $\eta_b \sim N(g_b(\kappa_b - \omega_bz_b^\top\gamma), g_b)$ with $g_b^{-1} = \omega_b + 1/(\lambda_b^2\gscale^2)$, for $b \in \mathcal V$.
\item Draw $\lambda_b^2 \sim \mathrm{IG}(1, 1/\nu_b + \eta_b^2/(2\gscale^2))$ and $\nu_b \sim \mathrm{IG}(1, 1 + 1/\lambda_b^2)$ for $b \in \mathcal V$.
\item Draw $\gscale^2 \sim \mathrm{IG}((m + 1)/2, 1/\xi + \sum_{\mathcal V}\eta_b^2/(2\lambda_b^2))$ and $\xi \sim \mathrm{IG}(1, 1 + 1/\gscale^2)$.
\end{enumerate}
To get a draw of $\psi$, draw a half Cauchy $\lambda_b$ and $\eta_b \sim N(0, \lambda_b^2\gscale^2)$ for each $b \notin \mathcal V$ and average $\expit(z_b^\top\gamma + \eta_b)$ over all $B$ cells. Each scan costs $O(m)$ plus the unvisited draws, which are needed only at kept scans. Keeping the unvisited cells in the chain would be valid but would tie $\gscale$ to $B - m$ prior draws of $\eta_b$ and slow it down.

Most visited cells hold one selected unit, a separated observation: for $S_b=M_b$ the likelihood tends to one as $\eta_b\to\infty$, leaving the prior's Cauchy like tail and no posterior mean of $\eta_b$. Chapter~\ref{ch:horseshoe}'s regularized horseshoe addresses this. Report the means of bounded $\theta_b$ and $\psi$, which exist. A cell seen once informs only $\E\{\expit(z_b^\top\gamma+\eta_b)\mid\gamma,\gscale\}$, so without repeats $\gscale$ is informed only through the fitted curve. Table~\ref{tab:weak}'s $n=100$ draws contain about five pairs in the same cell, fewer among selected units. They say little about $\gscale$ directly; as for $\sigma_\eta$ in Section~\ref{sec:rrw}, the curve matters for $\psi$.

As a check on the scan we took five cells with $\logit p_b = -1.5$, $-0.5$, 0, 0.5 and 1.5, $M = (3, 1, 2, 0, 0)$, $S = (3, 0, 1, 0, 0)$ and $v = 4$, and compared it with importance sampling from the prior (3,000,000 draws, effective sample size about 450,000). The sampler used the series of Appendix~\ref{app:toolkit} truncated at 60 terms, with the omitted terms replaced by their conditional mean at the current tilt, and kept 110,000 scans (seed 66). It gave $\E(\psi \mid \text{data}) = 0.4979$ with batch standard error 0.0007, against 0.4981, and $\Prob(\gscale > 1 \mid \text{data}) = 0.504$ with standard error 0.008, against 0.492. These bounded summaries leave the separated tails of $\eta_b$ unchecked.

\section*{Chapter~\ref{ch:critique}: Orthogonality Is Not Enough}

\paragraph{Exercise~\ref{ex:chcrit:1}.} Partialling out gives $Y - \ell(X) = \tau(X)\nu + \varepsilon$ with $\nu = T - e(X)$, so the population Robinson coefficient is $\E\{\nu(Y - \ell)\}/\E\nu^2 = \E\{\tau(X)\,\mathrm{Var}(T \mid X)\}/\E\,\mathrm{Var}(T \mid X)$, and $\mathrm{Var}(T \mid X) = e(X)\{1 - e(X)\}$ for binary $T$ (Proposition~\ref{dml:prop:wate}). The gap to the average effect is $\mathrm{Cov}(\tau(X), w(X))/\E\,w(X)$ with $w = e(1 - e)$, which Cauchy and Schwarz bounds in absolute value by $\mathrm{sd}(\tau)\,\mathrm{sd}(w)/\E\,w$, so an example needs strong heterogeneity and weights that vary. Let $X$ take two values with probability one half each, with $e = 0.5$ and $\tau = -1$ at the first and $e = 0.05$ and $\tau = 3$ at the second. The average effect is 1, the weights are 0.25 and 0.0475, and the partially linear estimand is
\[
\frac{0.25 \times (-1) + 0.0475 \times 3}{0.25 + 0.0475} = -0.361.
\]
It differs from the average effect by 1.361, more than the average effect itself, and has the opposite sign. The stratum that is treated at even odds and is harmed dominates the score, and the stratum that benefits is rarely treated.

\paragraph{Exercise~\ref{ex:chcrit:2}.} If $\tau(x) \equiv \tau_0$ then $\bar\tau = \tau_0$ and $\{\tau(X) - \bar\tau\}^2 = 0$ for every $X$, so the third term of \eqref{eq:dml:bound} is zero. With $e(x) = e$ and $\sigma_1^2 = \sigma_0^2 = \sigma^2$ the first two terms are $\sigma^2/e + \sigma^2/(1 - e)$, and
\[
V_{\mathrm{eff}} = \frac{\sigma^2}{e(1 - e)} + \mathrm{Var}\,\tau(X),
\]
which equals $\sigma^2/\{e(1 - e)\}$ when the effect is constant. When the outcome means do not depend on $X$ this is the large sample variance of $\sqrt n(\bar Y_1 - \bar Y_0)$ in a trial that treats a fraction $e$. It is smallest at $e = 1/2$, where it equals $4\sigma^2$, and rises to $11.1\sigma^2$ at $e = 0.1$ and to $101\sigma^2$ at $e = 0.01$.

\paragraph{Exercise~\ref{ex:chcrit:3}.} With $\delta_t = \hat\mu_t - \mu_t$, the display of Section~\ref{sec:crit:weights} gives the bias of the augmented score at the fitted nuisances as $\E[(\hat e - e)\{\delta_1/\hat e + \delta_0/(1 - \hat e)\}]$, while regression adjustment, the average of $\hat\mu_1 - \hat\mu_0$, has bias $\E(\delta_1 - \delta_0)$. Take two strata of probability one half with $\mu_1 = 1$ and $\mu_0 = 0$ in both, so $\tau = 1$. In stratum a every fit is correct. In stratum b the true propensity is 0.20 and the fitted one 0.05, the control regression is right and the treated regression is too high by one, so both models are misspecified. Regression adjustment returns 1.5, a bias of $+0.5$. In stratum b the augmented bias term is $(0.05 - 0.20)\times 1/0.05 = -3$, the outcome error multiplied by the propensity error and by $1/\hat e = 20$, so the bias is $\tfrac12(-3) = -1.5$ and the estimator returns $-0.5$: three times the bias of regression adjustment, in the other direction. With an error of one in both arms of stratum b, regression adjustment would be unbiased and the augmented bias would be $-1.58$.

\paragraph{Exercise~\ref{ex:chcrit:4}.} We read the coupled horseshoe of Table~\ref{tab:k401} as fixed and let the splits vary. Rerunning its sampler with the same seed and 2,000 kept draws reproduces the printed mean of 9.82 and interval $[7.42, 12.11]$ exactly, and its quartiles give the posterior 50\% interval $[9.01, 10.66]$, half width 0.83. All 20 saved lasso partitions of Chapter~\ref{ch:k401}, from 9.52 to 10.26, lie inside, so the estimated probability is $0/20$, with an exact 95\% upper bound of $1 - 0.05^{1/20} = 0.14$. To tighten the bound we ran 180 more partitions with the same score and nuisances (random states 20 to 199, 200 in all). None falls outside either. Over all 200, the estimates run from 9.51 to 10.47 with mean 9.92 and standard deviation 0.17, and the upper bound is $1 - 0.05^{1/200} = 0.015$.



\paragraph{Exercise~\ref{ex:chcrit:5}.} The net benefit of adopting is $\min(\tau, c) - k$ and that of not adopting is zero, so the Bayes action maximizes the posterior expected utility and adopts exactly when $\E\{\min(\tau, c) \mid \text{data}\} > k$. Since $\min(\tau, c) = \tau - (\tau - c)_+$,
\[
\E\{\min(\tau, c) \mid \text{data}\} = \E(\tau \mid \text{data}) - \E\{(\tau - c)_+ \mid \text{data}\},
\]
and the second term is the part of the posterior above the threshold, which its mean and variance do not fix. Take $c = 2$ and $k = 0.95$. A posterior $N(1, 1)$ has $\E(\tau - c)_+ = \varphi(1) - \Phi(-1) = 0.083$ ($\varphi$ and $\Phi$ the standard normal density and distribution function), so the expected benefit is $0.917 < k$ and the funder does not adopt. A posterior with mass one half on $0$ and on $2$ has the same mean and variance, puts no mass above $c$, and gives the expected benefit $1.000 > k$, so the funder adopts. Over all posteriors with mean $1$ and variance $1$ the expected benefit ranges over $[0.793, 1.000]$: the upper end is the two point law, and the lower end uses the sharp bound $\E(\tau - c)_+ \le \tfrac12\bigl[(m - c) + \sqrt{\sigma^2 + (m - c)^2}\bigr] = 0.207$.

\paragraph{Exercise~\ref{ex:chcrit:6}.} On the exact recovery event of the proof of Proposition~\ref{crit:prop:alignment} (Appendix~\ref{app:dmlproofs}), the Karush, Kuhn and Tucker conditions on $S$ give $\delta_m = u_m - \lambda_{\mathrm L}^m\hat\Sigma^{-1}z$ and $\delta_\ell = u_\ell - \lambda_{\mathrm L}^\ell\hat\Sigma^{-1}z$ there. The sign vector $z$ is the same in both because the signs agree, and $u_m$ and $u_\ell$ are the noise parts. With both penalties of order $\sqrt{\log p/n}$, the noise parts satisfy $\|u\|^2 = O_p(s/n)$, smaller than $\lambda_{\mathrm L}^2 s$ by a factor of order $1/\log p$, and $\hat\Sigma \to I$ in operator norm. Hence
\[
\langle\delta_m,\delta_\ell\rangle = \lambda_{\mathrm L}^m\lambda_{\mathrm L}^\ell s\{1 + o_p(1)\}, \qquad \|\delta_m\|^2 = (\lambda_{\mathrm L}^m)^2 s\{1 + o_p(1)\},
\]
and subtracting $\tau$ times the second from the first gives $b = \lambda_{\mathrm L}^m s(\lambda_{\mathrm L}^\ell - \tau\lambda_{\mathrm L}^m)$ to leading order. It vanishes exactly when $\lambda_{\mathrm L}^\ell = \tau\lambda_{\mathrm L}^m$, where the two shrinkage terms are equal and what is left is of lower order.

Let $\lambda_{\mathrm L}^m = \kappa\sigma_\nu$ and $\lambda_{\mathrm L}^\ell = \kappa\sigma_\ell$ with one constant $\kappa$, where $\sigma_\ell^2 = \tau^2\sigma_\nu^2 + \sigma_\varepsilon^2$ is the noise variance of the outcome equation. For $\tau > 0$,
\[
\lambda_{\mathrm L}^\ell - \tau\lambda_{\mathrm L}^m = \kappa(\sigma_\ell - \tau\sigma_\nu) = \frac{\kappa\sigma_\varepsilon^2}{\sigma_\ell + \tau\sigma_\nu},
\]
which is zero only if $\sigma_\varepsilon = 0$ (for $\tau \le 0$ it is positive whatever $\sigma_\varepsilon$ is). At $\tau = 1$ this gives $b = \kappa^2\sigma_\nu\sigma_\varepsilon^2 s/(\sigma_\ell + \sigma_\nu) > 0$, growing linearly in $s$. In the chapter's designs $\sigma_\nu = 1$ and $\sigma_\ell^2 = 1.9$, so $b = 0.378\,\kappa^2 s$.

\paragraph{Exercise~\ref{ex:chcrit:7}.} By the definition of $U_n$, $\hat\tau - \tau = \bar\Delta_n + U_n$ exactly, and subtracting the constant $\tau$ does not change a variance. The variable $\bar\Delta_n$ depends on the fits alone and satisfies $|\bar\Delta_n| \le \max(1, |\tau|)/J$ when $\|\delta_m\| + \|\delta_\ell\| \le 1$, as assumed in Theorem~\ref{crit:thm:single}, because $|b| \le \max(1, |\tau|)\|\delta_m\|(\|\delta_m\| + \|\delta_\ell\|)$ and the denominator is at least $J$. So $\bar\Delta_n$ and $U_n$ are square integrable when $\hat\tau$ is, and expanding $\mathrm{Var}(\bar\Delta_n + U_n)$ gives \eqref{eq:crit:cfvar}. If the sandwich estimates $\mathrm{Var}(U_n)$, as it is meant to, the variance it omits is $\mathrm{Var}(\hat\tau) - \mathrm{Var}(U_n) = \mathrm{Var}(\bar\Delta_n) + 2\,\mathrm{Cov}(\bar\Delta_n, U_n)$, which is $\mathrm{Var}(\bar\Delta_n)$ when the covariance is zero.

In the numbers, $2.151\times 10^{-4} - 2.365\times 10^{-4} + 3.562\times 10^{-3} = 3.5406\times 10^{-3}$ against $\mathrm{Var}(\hat\tau) = 3.540\times 10^{-3}$, equal to the rounding of the printed entries. The cross term is $-2.365/2.151 = -1.10$ times $\mathrm{Var}(\bar\Delta_n)$, so the omitted variance is about $-0.10\,\mathrm{Var}(\bar\Delta_n) = -2.1\times 10^{-5}$. It is negative because the cross term more than cancels $\mathrm{Var}(\bar\Delta_n)$. The block $\mathrm{Var}(U_n) = 3.562\times 10^{-3}$ is therefore slightly too large for $\mathrm{Var}(\hat\tau)$, by 0.6 percent in variance and 0.3 percent in standard error.

\paragraph{Exercise~\ref{ex:chcrit:8}.} Given the data and the fits, $\hat\nu_i$, $\hat\varepsilon_i$ and $\hat J_n$ are numbers, and only the $\xi_i$ are random. They are independent with mean zero and variance one, so $\E(\hat\tau^* \mid \cdot) = \hat\tau$ and
\[
\mathrm{Var}(\hat\tau^* \mid \cdot) = \frac{1}{\hat J_n^2 n^2}\sum_i\hat\nu_i^2\hat\varepsilon_i^2 = \frac{\sum_i\hat\nu_i^2\hat\varepsilon_i^2}{\bigl(\sum_i\hat\nu_i^2\bigr)^2},
\]
the square of the standard error in step 5 of the algorithm of Chapter~\ref{ch:dml}. With Gaussian multipliers the bootstrap interval is the sandwich interval up to Monte Carlo error in the quantiles, which is what Table~\ref{tab:crit:cmp} shows.

Lemma~\ref{dml:lem:exact} gives $\hat\tau-\tau=(S_n+C_n+B_n)/\hat J_n$. Conditional on the data and fits, $B_n/\hat J_n$ shifts $\hat\tau$ but is constant across bootstrap draws, so it adds nothing to their variance, and centring at $\hat\tau$ removes it. Thus the bootstrap reproduces noise around the realized fit, the first term of Theorem~\ref{crit:thm:single}, and omits the variation across fits, $\sigma_\Delta^2$, which requires refitting the nuisances.

\paragraph{Exercise~\ref{ex:chcrit:9}.} Under every $P_{j,\tau}$ the pair $(X, T)$ has the same law, because $T = m(X) + \nu$ does not involve $\tau$, and given $(X, T)$ the outcome is $N(\tau\nu + h(X), \sigma_\varepsilon^2)$ with $\nu = T - m(X)$. Two normal laws with the same variance and means a distance $\Delta$ apart have Kullback and Leibler divergence $\Delta^2/(2\sigma_\varepsilon^2)$, so by the chain rule
\[
\mathrm{KL}(P_{j,\tau}\,\|\,P_{j,0}) = \E\,\frac{(\tau\nu)^2}{2\sigma_\varepsilon^2} = \frac{\tau^2 j}{2\sigma_\varepsilon^2},
\]
and $n$ independent observations give $n\tau^2 j/(2\sigma_\varepsilon^2)$. Pinsker's inequality, $\|P - Q\|_{\mathrm{TV}} \le \sqrt{\mathrm{KL}(P\,\|\,Q)/2}$ with total variation taken as the supremum over events, gives $\|P_{j,\tau}^{\otimes n} - P_{j,0}^{\otimes n}\|_{\mathrm{TV}} \le |\tau|\sqrt{nj}/(2\sigma_\varepsilon)$. The triangle inequality through $P_{j,0}$ adds the two distances, $(|\tau| + |\tau'|)\sqrt{nj}/(2\sigma_\varepsilon) \le L\sqrt{nj}/\sigma_\varepsilon$. For $n = 300$, $\sigma_\varepsilon = 1$ and $L = 5$ the condition $5\sqrt{300j} \le 0.1$ means $j \le 4\times 10^{-4}/300 = 1.33\times 10^{-6}$, a treatment residual standard deviation $\sqrt j$ of at most 0.00115. At that $j$ the exact distance between $\tau = 5$ and $\tau' = -5$ is $\E\{2\Phi(5\|\nu\|) - 1\} = 0.080$, with $\Phi$ the standard normal distribution function and $\|\nu\|^2 \sim j\chi^2_n$, so the bound is not loose. In the sweep of Section~\ref{sec:crit:repair}, where the smallest residual standard deviation is 0.1, $j = 0.01$ and the bound with the same $n$, $\sigma_\varepsilon$ and $L$ is 8.7, vacuous because a total variation distance is at most 1.

\paragraph{Exercise~\ref{ex:chcrit:10}.} The weights are $0.0564$, $0.1078$, $0.1512$, $0.1868$, $0.2144$, $0.2342$, $0.2460$, $0.2500$, summing to 1.4469; weighted effects sum to 1.1150 (Table~\ref{tab:dml:strata}). Equal stratum sizes give $1.1150/1.4469=0.7706$, 23 percent below the average effect 1.0. Stratum 1, $e=0.06$ with the largest effect, gets 3.9 percent of the weight; stratum 8, $e=0.50$ with no effect, gets 17.3 percent. Reversing effects leaves weights unchanged and gives weighted sum 1.7787 and ratio $2-0.7706=1.2294$, now 23 percent above the average.

\section*{Chapter~\ref{ch:horseshoe}: The Horseshoe and the Coupled Prior}

\paragraph{Exercise~\ref{ex:ch07:1}.} Take $\gscale = 1$ (a general $\gscale$ replaces $p(\beta)$ by $p(\beta/\gscale)/\gscale$). Integrating the normal density against the half Cauchy density of $\lambda$ in \eqref{eq:hs},
\begin{align*}
p(\beta) &= \int_0^\infty \frac{e^{-\beta^2/(2\lambda^2)}}{\sqrt{2\pi}\,\lambda}\,\frac{2}{\pi(1+\lambda^2)}\,d\lambda \\
&= \frac{1}{\sqrt{2\pi}\,\pi}\int_0^\infty \frac{e^{-zu}}{1+u}\,du = \frac{1}{(2\pi^3)^{1/2}}\,e^{z}E_1(z), \qquad z = \frac{\beta^2}{2}.
\end{align*}
The second line substitutes $u = \lambda^{-2}$, which turns $d\lambda/\lambda$ into $du/(2u)$ and $1/(1+\lambda^2)$ into $u/(1+u)$, and then $t = z(1+u)$ in the exponential integral $E_1(z) = \int_z^\infty e^{-t}t^{-1}\,dt$. For the bounds write $\log(1 + c/z) = \int_0^\infty e^{-zu}(1 - e^{-cu})u^{-1}\,du$, and compare $e^zE_1(z) = \int_0^\infty e^{-zu}(1+u)^{-1}\,du$ with $\log(1+1/z)$ and $\tfrac12\log(1+2/z)$ through the integrands. The upper bound needs $(1+u)^{-1} < (1 - e^{-u})/u$, that is $1 + u < e^u$. The lower bound needs $(1+u)^{-1} > (1 - e^{-2u})/(2u)$, that is $(1+u)e^{-2u} > 1 - u$, which is trivial for $u \geq 1$ and for $u < 1$ reads $\log\{(1+u)/(1-u)\} > 2u$, true because the series of the left side is $2u$ plus positive terms. Hence $\tfrac12\log(1+2/z) < e^zE_1(z) < \log(1+1/z)$, and with $z = \beta^2/2$
\[
\frac{1}{2(2\pi^3)^{1/2}}\log\Bigl(1 + \frac{4}{\beta^2}\Bigr) < p(\beta) < \frac{1}{(2\pi^3)^{1/2}}\log\Bigl(1 + \frac{2}{\beta^2}\Bigr).
\]
Both bounds grow like $\log(1/\beta^2)$ at the origin, so $p$ is unbounded there, and both equal $2/\{(2\pi^3)^{1/2}\beta^2\}$ to first order for large $\beta$, the $\beta^{-2}$ tail that gives $\alpha = 2$ in the tail robustness theorem. At $\beta = 1$ they are 0.102 and 0.140, around $p(1) = 0.117$.

\paragraph{Exercise~\ref{ex:ch07:2}.} For the Laplace prior $\pi(\beta) = \frac{\lambda}{2}e^{-\lambda|\beta|}$ and $y \mid \beta \sim N(\beta, 1)$, completing the square on each half line gives the marginal $m(y) = \frac{\lambda}{2}e^{\lambda^2/2}(a_+ + a_-)$ with $a_+ = e^{-\lambda y}\Phi(y - \lambda)$ and $a_- = e^{\lambda y}\Phi(-y - \lambda)$. The two normal density terms of $(a_+ + a_-)'$ are both $(2\pi)^{-1/2}e^{-(y^2 + \lambda^2)/2}$ and cancel, so $\frac{d}{dy}\log m(y) = -\lambda(a_+ - a_-)/(a_+ + a_-)$. Since $a_-/a_+ \to 0$ as $y \to \infty$ and $a_+/a_- \to 0$ as $y \to -\infty$, $\frac{d}{dy}\log m(y) \to -\lambda\,\mathrm{sign}(y)$, a constant shrinkage. For the normal prior $N(0, v)$, $m$ is $N(0, 1 + v)$ and $\frac{d}{dy}\log m(y) = -y/(1 + v)$, shrinkage growing linearly in $y$.

\paragraph{Exercise~\ref{ex:ch07:3}.} Numerical integration of the posterior mean formula of Chapter~\ref{ch:horseshoe} gives $y - \E(\beta \mid y) \approx 0.34$ at $y = 6$, $0.25$ at $y = 8$ and $0.20$ at $y = 10$, against $2/y = 0.33$, $0.25$ and $0.20$. The approximation improves as $y$ grows.

\paragraph{Exercise~\ref{ex:ch07:4}.} The model is $Y\mid\beta,\sigma^2\sim N(Z\beta,\sigma^2I)$, $p(\sigma^2)\propto1/\sigma^2$, $\beta_j\mid\lambda_j^2,\gscale^2,\sigma^2\sim N(0,\lambda_j^2\gscale^2\sigma^2)$, with $\lambda_j^2\mid\nu_j\sim\mathrm{IG}(1/2,1/\nu_j)$, $\nu_j\sim\mathrm{IG}(1/2,1)$, and the same mixture for $\gscale^2\mid\xi$ and $\xi$. Integrating out $\nu_j$ gives $p(x)=1/\{\pi\sqrt x(1+x)\}$ for $x=\lambda_j^2$. Multiplying by $dx/d\lambda_j=2\lambda_j$ gives the half Cauchy density $2/\{\pi(1+\lambda_j^2)\}$. Completing the square in $\beta$ gives mean $A^{-1}Z^\top Y$ and covariance $\sigma^2A^{-1}$, with $A$ as in Section~\ref{sec:hsregression}. Collecting in the joint density the powers $x^{-a-1}e^{-b/x}$ of each variable, where $\nu_j^{-1/2}$ and $\xi^{-1/2}$ are the normalizing constants of the inverse gamma densities of $\lambda_j^2$ and $\gscale^2$, gives
\begin{align*}
\sigma^2\mid\cdot&\sim\mathrm{IG}\Bigl(\frac{n+p}{2},\frac{\|Y-Z\beta\|^2+\sum_j\beta_j^2/(\lambda_j^2\gscale^2)}{2}\Bigr),\\
\lambda_j^2\mid\cdot&\sim\mathrm{IG}\Bigl(1,\frac1{\nu_j}+\frac{\beta_j^2}{2\gscale^2\sigma^2}\Bigr),\qquad
\nu_j\mid\cdot\sim\mathrm{IG}\Bigl(1,1+\frac1{\lambda_j^2}\Bigr),\\
\gscale^2\mid\cdot&\sim\mathrm{IG}\Bigl(\frac{p+1}{2},\frac1\xi+\sum_j\frac{\beta_j^2}{2\lambda_j^2\sigma^2}\Bigr),\qquad
\xi\mid\cdot\sim\mathrm{IG}\Bigl(1,1+\frac1{\gscale^2}\Bigr).
\end{align*}
Here $b$ is the scale in $\mathrm{IG}(a,b)$. These are steps 2 to 4 of Section~\ref{sec:hsregression}.

\paragraph{Exercise~\ref{ex:ch07:5}.} Under \eqref{eq:shared} with the inverse gamma mixture for $\lambda_j$, the full conditional of $\lambda_j^2$ is inverse gamma with shape $3/2$ and rate $1/\nu_j + \beta_j^2/(2\sigma_Y^2\tau_{\mathrm{g},Y}^2) + \gamma_j^2/(2\sigma_T^2\tau_{\mathrm{g},T}^2)$. Both coefficients enter, which is why a large $\gamma_j$ alone, as for an instrument, inflates $\lambda_j$ and relaxes the shrinkage of $\beta_j$.

\paragraph{Exercise~\ref{ex:ch07:6}.} Put $W=Y-Xb$, $c=\hat m^\top T/(\hat m^\top\hat m)$ and $\tilde T=T-c\hat m$, assuming $(T,\hat m)$ has full column rank. Since $\tilde T^\top\hat m=0$, the regression of $W$ on $(\tilde T,\hat m)$ decouples and has fit $a\tilde T+d\hat m$, where
\[
a=\frac{\tilde T^\top W}{\tilde T^\top\tilde T},\qquad d=\frac{\hat m^\top W}{\hat m^\top\hat m}.
\]
The two pairs of columns have the same span and $a\tilde T+d\hat m=aT+(d-ac)\hat m$, so uniqueness makes $a$ the coefficient of $T$ in the original regression, as well as the coefficient from regressing $W$ on $\tilde T$ alone. Given $b,\sigma^2$, flat priors on $\tau,\delta$ make $a$ the Gaussian posterior mean of $\tau$. Also $\tilde T^\top T=\tilde T^\top\tilde T$, the identity used in \eqref{eq:pcbias}.

\paragraph{Exercise~\ref{ex:ch07:7}.} Take $n = 200$ and $p = 50$ standard normal covariates, five of them with $\beta_j = \gamma_j = 0.5$ and every other coefficient zero in both equations, $T = X^\top\gamma + \nu$ and $Y = T + X^\top\beta + \varepsilon$ with standard normal $\nu$ and $\varepsilon$, so that $\tau = 1$. Fit the independent and the shared scale horseshoe by the joint Gibbs sampler of Section~\ref{sec:sharedsim}, 1,500 sweeps with the first 500 discarded, in each of 200 data sets. The independent prior has bias 0.074 (Monte Carlo standard error 0.006), RMSE 0.114 and coverage 0.840. The shared prior has bias 0.046 (0.005), RMSE 0.086 and coverage 0.905, and it is closer to the truth in 78 percent of the data sets. Least squares on all covariates has bias 0.005 and RMSE 0.079. With confounders alone, sharing the scales removes about 40 percent of the bias, because $\gamma_j$ is read from the treatment model and protects $\beta_j$. The bias of either prior is the shrinkage error on the five coefficients projected on $\gamma$ and divided by $\mathrm{Var}(T) = 2.25$. Adding five instruments ($\gamma_j = 1$) and five prognostic covariates ($\beta_j = 1$), the design of Section~\ref{sec:sharedsim}, raises $\mathrm{Var}(T)$ to 7.25, which predicts $0.074 \times 2.25/7.25 = 0.023$ for the independent prior, and reverses the order. Over 100 data sets the independent prior has bias 0.021 (0.003) and RMSE 0.038, the shared prior 0.037 (0.007) and 0.079, because sharing protects the instruments as well.

\paragraph{Exercise~\ref{ex:ch07:8}.} Work in the normal means problem $y_j \mid \beta_j \sim N(\beta_j, 1)$ and write $s = \lambda_j\gscale$ for the prior scale of $\beta_j$, so that $\kappa_j = 1/(1 + s^2)$, $s^2 = (1 - \kappa)/\kappa$ and $|ds/d\kappa| = 1/(2s\kappa^2)$. Under the horseshoe $s \sim C^+(0, \gscale)$. The horseshoe plus of Bhadra, Datta, Polson and Willard (2017) has $\lambda_j \mid \eta_j \sim C^+(0, \eta_j)$ and $\eta_j \sim C^+(0, 1)$, so $u = s/\gscale$ is a product of two independent half Cauchy variables, with density
\[
\frac{4}{\pi^2}\int_0^\infty \frac{\eta\,d\eta}{(\eta^2 + u^2)(1 + \eta^2)} = \frac{2}{\pi^2}\int_0^\infty \frac{dv}{(v + u^2)(1 + v)} = \frac{4\log u}{\pi^2(u^2 - 1)}.
\]
A change of variables gives the two priors on $\kappa$, with $r = s^2/\gscale^2 = (1 - \kappa)/(\gscale^2\kappa)$:
\[
p_{\mathrm{HS}}(\kappa) = \frac{\gscale}{\pi}\,\frac{\kappa^{-1/2}(1 - \kappa)^{-1/2}}{1 - \kappa + \gscale^2\kappa}, \qquad
\frac{p_{+}(\kappa)}{p_{\mathrm{HS}}(\kappa)} = R(r) = \frac{1}{\pi}\,\frac{r + 1}{r - 1}\,\log r.
\]
At $\gscale = 1$ the first is the Beta$(1/2, 1/2)$ density and the second is that density times $\log\{(1 - \kappa)/\kappa\}/\{\pi(1 - 2\kappa)\}$. The factor $R$ satisfies $R(r) = R(1/r)$, has its minimum $2/\pi = 0.64$ at $r = 1$, equals one at $|\log r| = 2.77$ and grows like $|\log r|/\pi$. So the plus prior moves mass from the middle to both ends, by a logarithmic factor. At $\gscale = 1$ it raises $\Prob(\kappa > 0.99)$ from 0.064 to 0.135 and lowers $\Prob(0.1 < \kappa < 0.9)$ from 0.589 to 0.417.

Comparing the two at one value of $\gscale$ misleads, because the extra mass also lies at $s \gg \gscale$. For small $\gscale$ the prior chance of escaping shrinkage is $\Prob(\kappa < 1/2) = (2/\pi)\arctan\gscale \approx 2\gscale/\pi$ under the horseshoe and about $(4\gscale/\pi^2)\{\log(1/\gscale) + 1\}$ under the plus prior. Since $y \mid \kappa \sim N(0, 1/\kappa)$, the posterior density of $\kappa$ is proportional to $\kappa^{1/2}e^{-\kappa y^2/2}p(\kappa)$. Integrating against $1 - \kappa$ at $y = 0$, where the normalizing constant tends to one as $\gscale \to 0$, gives the expected shrinkage of a null
\[
\E(1 - \kappa \mid y = 0) \approx \frac{2\gscale}{\pi} \ \text{(horseshoe)}, \qquad \E(1 - \kappa \mid y = 0) \approx \frac{4\gscale}{\pi^2}\log\frac{1}{2\gscale} \ \text{(plus)},
\]
the second from $\int_0^1 u^{-1/2}\log\{u/(\gscale^2(1 - u))\}\,du = 4\log\{1/(2\gscale)\}$ with $u = 1 - \kappa$. At equal $\gscale$ the plus prior therefore shrinks a null less. A sparse problem pushes its global scale lower to compensate. For a large coefficient, Tweedie's formula \eqref{eq:tweedie} and the tail $p_{+}(\beta) \propto \log(\beta/\gscale)/\beta^2$ give
\[
y\{y - \E(\beta \mid y)\} \approx 2 - \frac{1}{\log(y/\gscale)},
\]
against 2 for the horseshoe: the plus prior shrinks a large confounder less, but the gain fades like $1/\log y$. At $\gscale = 1$ the exact values at $y = 8, 10, 20, 50$ are 1.61, 1.62, 1.68, 1.74 against 2.03, 2.02, 2.005, 2.001.

The size of the effect at a sparsity level of interest takes a calculation. Take two signals of size 5 among 2,000 means and compute the expected total squared error of the posterior mean by quadrature on $\log s$ and $y$, with $\gscale$ chosen to minimize it for each prior. The best scales are 0.0072 (horseshoe) and 0.0015 (plus), a fifth as large, the errors are 10.66 and 10.55, and the shrinkage factors $\E(1 - \kappa \mid y)$ nearly coincide: 0.005 and 0.004 at $y = 0$, 0.050 and 0.049 at $y = 3$, 0.890 and 0.897 at $y = 5$. A null is shrunk almost the same by both once the scales are matched to the sparsity. So is a confounder, and that is what the bias depends on. In the bias formula of Chapter~\ref{ch:horseshoe}, $\gamma^\top\Sigma_X(\beta - \hat\beta)/\mathrm{Var}(T)$, with two orthogonal unit variance confounders, $\gamma_j = 1$, $\mathrm{Var}(T) = 3$, $n = 1{,}000$, $\sigma = 1$ and signals of $\mu = 3$ standard errors, both priors give a bias of 0.047. Neither prior looks at $\gamma_j$, so coupling the scales addresses that bias and changing the local mixing does not. This holds at the best fixed $\gscale$ in squared error; we did not study a learned $\gscale$.

\paragraph{Exercise~\ref{ex:chbaipw:1}.} Write $\eta=(\mu_1,\mu_0,e)$ and $\psi(W;\eta)=\mu_1(X)-\mu_0(X)+h\{Y-\mu_T(X)\}$, with $h=T/e(X)-(1-T)/\{1-e(X)\}$ as in Section~\ref{sec:baipw}. At the truth $\mu_t(x)=\E(Y\mid T=t,X=x)$ and $e_0(x)=\Prob(T=1\mid X=x)$, differentiation along $\eta_0+r(\Delta_1,\Delta_0,\Delta_e)$ gives conditional mean $\Delta_1-\Delta_0-\Delta_1+\Delta_0=0$, since $\E(T\mid X)=e_0$ and both arm residuals have conditional mean zero. This proves orthogonality without using how the nuisances were fitted. For any fixed $\hat\eta$ with $\hat e\in(0,1)$, including fits conditional on training folds,
\[
\E\{\psi(W;\hat\eta)\}-\tau_0=\E\Bigl[\frac{\hat e-e_0}{\hat e}(\hat\mu_1-\mu_1)+\frac{\hat e-e_0}{1-\hat e}(\hat\mu_0-\mu_0)\Bigr].
\]
The identity has no first order term and vanishes if either $\hat e=e_0$ or both outcome regressions are correct. The horseshoe must make these products small; its posterior mean propensity lies in $(0,1)$, so the identity applies as for a lasso or forest.

\paragraph{Exercise~\ref{ex:chbaipw:2}.} Fix a period and drop the subscripts $it$. The chapter leaves $\hat\mu^V$ open, so take the following setting. Let $\mathcal H$ be the history observed before the period, with $R \perp (T, X, Y) \mid \mathcal H$ and $\pi_0 = \Prob(R = 1 \mid \mathcal H)$. Let $V_0 = T - m_0(X)$ and $\varepsilon = Y - \ell_0(X) - \tau_0V_0$, with $\E(V_0 \mid X) = 0$ and $\E(\varepsilon \mid X, T) = 0$. Write $S(\ell, m) = (T - m)\{Y - \ell - \tau_0(T - m)\}$ for the complete data score at $\tau_0$, $\mu_0 = \E\{S(\ell_0, m_0) \mid \mathcal H\}$ and $\bar S = \E\{S(\hat\ell, \hat m) \mid \mathcal H\}$, and treat the nuisances as fixed, as cross fitting allows. The score is $\psi = (R/\hat\pi)S(\hat\ell, \hat m) - (R/\hat\pi - 1)\hat\mu$. Conditioning on $\mathcal H$,
\[
\E(\psi \mid \mathcal H) = \bar S + \frac{\pi_0 - \hat\pi}{\hat\pi}(\bar S - \hat\mu).
\]
With $\Delta_m = \hat m - m_0$ and $\Delta_\ell = \hat\ell - \ell_0$, the complete data score expands exactly as
\[
\begin{split}
S(\hat\ell, \hat m) &= (V_0 - \Delta_m)(\varepsilon - \Delta_\ell + \tau_0\Delta_m)\\ &= V_0\varepsilon - V_0\Delta_\ell + \tau_0V_0\Delta_m - \varepsilon\Delta_m + \Delta_m\Delta_\ell - \tau_0\Delta_m^2.
\end{split}
\]
The first four terms have mean zero, because $\Delta_m$ and $\Delta_\ell$ are functions of $X$. Hence the remainder is, exactly,
\[
\E\,\psi = \E(\Delta_m\Delta_\ell) - \tau_0\E(\Delta_m^2) + \E\Bigl\{\frac{\pi_0 - \hat\pi}{\hat\pi}(\bar S - \hat\mu)\Bigr\}.
\]
It has no first order term. Since $\E(V_0\varepsilon) = 0$, it is the bias of the score at $\tau_0$, and dividing by $\E(V_0^2)$ gives the bias of $\hat\tau$ to leading order, once the fitted Jacobian $\E(T - \hat m)^2$ is close to $\E(V_0^2)$. To find the other products, open the last factor:
\[
\begin{split}
\bar S - \hat\mu = (\mu_0 - \hat\mu) &+ \E\{-\Delta_\ell r_m + \tau_0\Delta_mr_m - \Delta_mr_\varepsilon \mid \mathcal H\} \\
&+ \E(\Delta_m\Delta_\ell - \tau_0\Delta_m^2 \mid \mathcal H),
\end{split}
\]
with $r_m = \E(V_0 \mid X, \mathcal H)$ and $r_\varepsilon = \E(\varepsilon \mid X, \mathcal H)$. If $\hat\pi \geq c > 0$, $|r_m|$ and $|r_\varepsilon|$ are bounded and the nuisance errors are bounded, Cauchy and Schwarz give, with a constant $C$,
\[
\begin{split}
|\E\,\psi| \leq{}& \|\Delta_m\|\,\|\Delta_\ell\| + |\tau_0|\,\|\Delta_m\|^2 \\
&+ C\|\hat\pi - \pi_0\|\bigl(\|\hat\mu - \mu_0\| + \|\Delta_\ell\| + \|\Delta_m\|\bigr).
\end{split}
\]
All five products of the display arise, from different sources. The complete data remainder gives $\Delta_m\Delta_\ell$ and $-\tau_0\Delta_m^2$; missingness adds the product of $\hat\pi-\pi_0$ and $\hat\mu-\mu_0$. Products with $\Delta_\ell$ and $\Delta_m$ arise through $r_m,r_\varepsilon$. If $r_m=r_\varepsilon=0$, those two products vanish, leaving the $\hat\mu-\mu_0$ product and products with quadratic terms. If $\hat\mu$ instead regresses $S(\hat\ell,\hat m)$ on $\mathcal H$, then $\bar S-\hat\mu$ is that regression's error, and its product with $\hat\pi-\pi_0$ involves neither $\Delta_\ell$ nor $\Delta_m$.

Under strict exogeneity, demeaning over all periods preserves the expansion with $X$ the whole covariate path and demeaned nuisance errors. It makes $r_\varepsilon$ nonzero because past outcomes reveal earlier errors in $\varepsilon_{it}-\bar\varepsilon_i$, producing a $(\hat\pi-\pi_0)\Delta_m$ term with coefficient of order the reciprocal of the number of periods. With missing periods, however, the observed unit mean depends on later $R$. Then $R\perp S\mid\mathcal H$ fails and the first display no longer holds; the unit mean is another nuisance. The displayed condition therefore remains unproved for the score as the chapter writes it, and the proof would need $\hat\mu^V$ and the demeaning to be specified. The expansion concerns bias only; variance also needs $\hat\pi$ bounded away from zero and folds by unit.

\paragraph{Exercise~\ref{ex:chbaipw:3}.} Overlap lets a group of units sit at $e_0 = 0.05$, and exactly one of $n_A$ such units is treated with probability $n_A(0.05)(0.95)^{n_A - 1}$. Take one binary covariate, $n_A = 150$ units at $X = 0$ with $e_0 = 0.05$ and $n_B = 50$ at $X = 1$ with $e_0 = 0.5$, so $n = 200$, and fit $\expit(\alpha + \gamma X)$ with the horseshoe prior \eqref{eq:hs} on $\gamma$ and a near flat $N(0, 10^4)$ prior on the intercept. Suppose one unit of group $A$ is treated, an event of probability 0.0036, and 25 units of group $B$. The posterior mean propensity of group $A$, by quadrature over $(\alpha, \gamma)$ and confirmed by importance sampling, is $\hat e_A = 0.0086$, against the sample proportion $1/150 = 0.0067$ and the true 0.05. The horseshoe pulled $\gamma$ toward zero and raised the estimate only from 0.0067 to 0.0086, because the data insist on a difference between the groups and the heavy tails leave $\gamma$ alone. The treated unit then has $h = 1/\hat e_A = 116$, against $1/0.05 = 20$ at the overlap bound, and $\hat e_A < 0.01$ for every number of treated units in group $B$ from 10 to 50. This sample is possible under overlap and has positive probability, so $|h_i| \leq 20$ cannot hold almost surely at $n = 200$.

\paragraph{Exercise~\ref{ex:chbaipw:4}.} Within fold $k$ with index set $I_k$ and $n_k$ units, normalize the weights of each arm to average one. Let $w_i = 1/\hat e(X_i)$, $w_i' = 1/\{1 - \hat e(X_i)\}$, $\bar w_{1k} = n_k^{-1}\sum_{i \in I_k}T_iw_i$ and $\bar w_{0k} = n_k^{-1}\sum_{i \in I_k}(1 - T_i)w_i'$. The score of unit $i \in I_k$ becomes
\[
\hat\psi_i^{H} = \hat\mu_1(X_i) - \hat\mu_0(X_i) + \frac{T_iw_i\{Y_i - \hat\mu_1(X_i)\}}{\bar w_{1k}} - \frac{(1 - T_i)w_i'\{Y_i - \hat\mu_0(X_i)\}}{\bar w_{0k}},
\]
and the estimate and its standard error are the mean of the $\hat\psi_i^{H}$ and their standard deviation divided by $\sqrt n$, as before. Design III has $n = 1{,}000$ and $p = 100$ covariates, the first ten driving treatment, with treatment probability $\expit(2s)$ for $s = 10^{-1/2}\sum_{j \leq 10}X_j$ and $Y = T + s + \varepsilon$. We ran both versions on the same horseshoe fits and folds (600 sweeps, five stratified folds) in 120 data sets. The estimates have variance 0.0152 with Horvitz Thompson weights and 0.0119 with normalized weights, a ratio of 0.78. The bias is $-0.009$ and 0.001 (Monte Carlo standard error 0.011), the mean reported standard error falls from 0.119 to 0.106, and coverage is 0.958 and 0.950. The normalized estimate is closer to the truth in 69 percent of the data sets, and its squared error is lower by 0.0033 on average (paired $t = -4.8$). The gain sits where the weights are large: in the 60 data sets with a largest weight above 50 the RMSE is 0.133 and 0.112, and in the other 60 it is 0.113 and 0.106.

\section*{Chapter~\ref{ch:computation}: Computing the Posterior}

\paragraph{Exercise~\ref{ex:ch14:1}.} Write $\Lambda = \mathrm{diag}(\lambda_j^2\gscale^2)$, so that $A = Z^\top Z + \Lambda^{-1}$ as in Chapter~\ref{ch:horseshoe}, and let $M = I_n + Z\Lambda Z^\top$, an $n \times n$ matrix. The Woodbury identity gives
\[
A^{-1} = \Lambda - \Lambda Z^\top M^{-1} Z\Lambda,
\]
and $A\Lambda Z^\top = Z^\top M$ gives $A^{-1}Z^\top = \Lambda Z^\top M^{-1}$. The posterior mean is therefore $KY$ with $K = \Lambda Z^\top M^{-1}$, which needs an $n \times n$ solve and no $p \times p$ factorization. For the covariance, draw $u \sim N(0, \sigma^2\Lambda)$ and $\delta \sim N(0, \sigma^2 I_n)$ independently, put $v = Zu + \delta$, and let $\beta = KY + (u - Kv)$. This is Gaussian with mean $KY$, and $\mathrm{Cov}(u, v) = \sigma^2\Lambda Z^\top$ and $\mathrm{Var}(v) = \sigma^2 M$ give
\begin{align*}
\mathrm{Var}(u - Kv) &= \sigma^2\bigl\{\Lambda - 2\Lambda Z^\top M^{-1}Z\Lambda + \Lambda Z^\top M^{-1}MM^{-1}Z\Lambda\bigr\}\\
&= \sigma^2\bigl\{\Lambda - \Lambda Z^\top M^{-1}Z\Lambda\bigr\} = \sigma^2 A^{-1}.
\end{align*}
Written out, $\beta = u + \Lambda Z^\top w$ with $Mw = Y - v$: draw $u$ and $\delta$, form $v$, solve for $w$, return $\beta$. This is the sampler of Bhattacharya, Chakraborty and Mallick (2016). Forming $Z\Lambda Z^\top$ costs $O(n^2p)$, the Cholesky factor of $M$ and the solve $O(n^3)$, and everything else $O(np)$, against $O(p^3)$ for the factorization of $A$ in Section~\ref{sec:hsgibbs}.

\paragraph{Exercise~\ref{ex:ch14:2}.} A Polya Gamma variable $\omega \sim \mathrm{PG}(1, 0)$ is, by the definition of Appendix~\ref{app:toolkit}, the infinite convolution
\[
\omega = \frac{1}{2\pi^2}\sum_{k=1}^\infty \frac{g_k}{(k - 1/2)^2}, \qquad g_k \text{ independent standard exponential}.
\]
Its density has no closed form, but its Laplace transform has. An exponential variable has $\E e^{-sg} = 1/(1 + s)$, so independence gives
\[
\E e^{-t\omega} = \prod_{k=1}^\infty \Bigl\{1 + \frac{t/2}{\pi^2(k - 1/2)^2}\Bigr\}^{-1} = \frac{1}{\cosh\sqrt{t/2}},
\]
by the product formula $\cosh z = \prod_{k \ge 1}\{1 + z^2/(\pi^2(k - 1/2)^2)\}$ at $z^2 = t/2$. Taking $t = \zeta^2/2$ gives $\int_0^\infty e^{-\omega\zeta^2/2}p(\omega \mid 1, 0)\,d\omega = 1/\cosh(\zeta/2)$. The right side of the identity of Section~\ref{comp:sec:cost} with $b = 1$ and $\tilde y = a - 1/2$ is then
\[
\frac12\,\frac{e^{(a - 1/2)\zeta}}{\cosh(\zeta/2)} = \frac{e^{(a - 1/2)\zeta}}{e^{\zeta/2} + e^{-\zeta/2}} = \frac{e^{a\zeta}}{1 + e^{\zeta}},
\]
which is the left side. With gamma variables of shape $b$ in place of the exponentials the same product gives $\cosh^{-b}\sqrt{t/2}$, and the identity for every $b$.

\paragraph{Exercise~\ref{ex:ch14:3}.} Let $g(\theta) = \log p(D \mid \theta) + \log\pi(\theta)$ with mode $\hat\theta$, and write $-g''(\hat\theta) = na$, $g'''(\hat\theta) = nb_3$ and $g''''(\hat\theta) = nb_4$. The regularity conditions are these, with constants free of $n$: $0 < c_a \leq a \leq C_a$ and $b_3$, $b_4$ stay bounded; on a fixed neighbourhood $|\theta - \hat\theta| \leq \delta$, $g'' \leq -na/2$ and $|g^{(5)}| \leq Cn$; outside it $\log p(D \mid \theta) - \log p(D \mid \hat\theta) \leq -n\epsilon$; and $\pi$ is proper with $\pi(\hat\theta) \geq c_\pi > 0$. The Laplace approximation of Section~\ref{comp:sec:cost}, integrated, is $\hat Z = e^{g(\hat\theta)}(2\pi/(na))^{1/2}$.

Put $z = (na)^{1/2}(\theta - \hat\theta)$, $c_3 = b_3/(6a^{3/2})$ and $c_4 = b_4/(24a^2)$. On $|z| \leq n^{1/8}$, Taylor's theorem gives
\[
g(\theta) - g(\hat\theta) = -\frac{z^2}{2} + T(z), \qquad T(z) = \frac{c_3z^3}{n^{1/2}} + \frac{c_4z^4}{n} + O\Bigl(\frac{|z|^5}{n^{3/2}}\Bigr),
\]
and $|T| = O(n^{-1/8})$ there, so
\[
e^{T} = 1 + \frac{c_3z^3}{n^{1/2}} + \frac{c_4z^4}{n} + \frac{c_3^2z^6}{2n} + O\Bigl(\frac{|z|^5 + |z|^7 + |z|^9}{n^{3/2}}\Bigr).
\]
Integrating against $e^{-z^2/2}$ over this range differs from the integral over the whole line by a term of order $e^{-n^{1/4}/3}$. The odd term $c_3z^3n^{-1/2}$ integrates to zero, which is why the rate is $n^{-1}$ and not $n^{-1/2}$. The moments $\E z^4 = 3$ and $\E z^6 = 15$ for a standard normal give
\[
\frac{1}{\hat Z}\int_{|z| \leq n^{1/8}} e^{g(\theta)}\,d\theta = 1 + \frac{1}{n}\Bigl(\frac{b_4}{8a^2} + \frac{5b_3^2}{24a^3}\Bigr) + O(n^{-3/2}).
\]
The two outer regions are negligible. For $n^{1/8} < |z|$ and $|\theta - \hat\theta| \leq \delta$, the bound on $g''$ gives $g(\theta) - g(\hat\theta) \leq -z^2/4$, so their contribution relative to $\hat Z$ is at most of order $e^{-n^{1/4}/5}$. For $|\theta - \hat\theta| > \delta$,
\[
\int e^{g(\theta)}\,d\theta \leq p(D \mid \hat\theta)e^{-n\epsilon}\int\pi(\theta)\,d\theta = e^{g(\hat\theta) - n\epsilon}/\pi(\hat\theta),
\]
exponentially small relative to $\hat Z$. Hence $\int p(D \mid \theta)\pi(\theta)\,d\theta = \hat Z\{1 + O(n^{-1})\}$. As a check take $g(\theta) = k\log\theta - b\theta$, whose integral is $\Gamma(k + 1)/b^{k + 1}$. Here $na = k/\hat\theta^2$, $nb_3 = 2k/\hat\theta^3$ and $nb_4 = -6k/\hat\theta^4$, so the correction is $-3/(4k) + 5/(6k) = 1/(12k)$, the first term of Stirling's series, and numerically $\hat Z/Z - 1$ is $-8.33\times 10^{-4}$, $-2.08\times 10^{-4}$ and $-5.21\times 10^{-5}$ at $k = 100, 400, 1{,}600$.

For the density, take $n$ exponential observations with rate $\theta$ and a Gamma$(a_0, b_0)$ prior, so the posterior is Gamma$(k + 1, b)$ with $k = a_0 + n - 1$ and $b = b_0 + \sum_i y_i$. Its mode is $\hat\theta = k/b$ and $-g''(\hat\theta) = k/\hat\theta^2$, so the normal approximation $q$ has standard deviation $\hat\theta/k^{1/2}$. At $\theta = \hat\theta(1 + zk^{-1/2})$ with $z$ fixed,
\[
\log\frac{p(\theta \mid D)}{p(\hat\theta \mid D)} = k\log\Bigl(1 + \frac{z}{k^{1/2}}\Bigr) - k^{1/2}z = -\frac{z^2}{2} + \frac{z^3}{3k^{1/2}} + O(k^{-1}).
\]
The normal has only the first term, and $p(\hat\theta \mid D)/q(\hat\theta) = \hat Z/Z = 1 + O(k^{-1})$, so
\[
\frac{p(\theta \mid D)}{q(\theta)} - 1 = \frac{z^3}{3k^{1/2}} + O(k^{-1}),
\]
of exact order $n^{-1/2}$ for $z \neq 0$. At $z = 2$ the exact values are 0.260, 0.132 and 0.066 at $k = 100, 400, 1{,}600$, halving as $k$ quadruples, against $8/(3k^{1/2}) = 0.267$, 0.133 and 0.067. The two rates do not conflict: the $n^{-1/2}$ term of the density is odd in $z$, the skewness correction, and it integrates to zero.

\paragraph{Exercise~\ref{ex:ch14:4}.} Take $Y_i = \tau T_i + g(X_i) + \varepsilon_i$ and $T_i = m(X_i) + \nu_i$, with target $\tau$. The summary is
\[
S(D) = \Bigl(n,\ \hat\tau,\ \hat J,\ \hat\sigma^2,\ \hat\Omega,\ \hat R^2_Y,\ \hat R^2_T,\ \hat\tau_{\mathrm{plug}},\ \hat\tau_{\mathrm{ols}}\Bigr),
\]
computed with cross fitted residuals $\hat V_i = T_i - \hat m(X_i)$ and $\hat U_i = Y_i - \hat\ell(X_i)$, where $\hat\tau = \sum_i\hat V_i\hat U_i/\sum_i\hat V_i^2$, $\hat J = n^{-1}\sum_i\hat V_i^2$, $\hat\sigma^2 = n^{-1}\sum_i(\hat U_i - \hat\tau\hat V_i)^2$ and $\hat\Omega = n^{-1}\sum_i\hat V_i^2(\hat U_i - \hat\tau\hat V_i)^2$. The last four are the out of fold $R^2$ of the two nuisance fits, the plug in coefficient of $T$ from a regularized outcome regression, and least squares on $T$ and a low dimensional basis of $X$ (or the raw contrast when $p \geq n$). As Section~\ref{comp:sec:summaries} requires, the cross fit runs inside $S$ on every simulated data set. The argument for this choice has three steps.

First, a case with exact sufficiency. Let $g = \Phi\beta$ with $(T, \Phi)$ of full rank and Gaussian errors. The factorization of the likelihood shows that $(T^\top Y, \Phi^\top Y, Y^\top Y)$ is sufficient given the design, and adding $(\Phi^\top T, T^\top T)$ covers the treatment equation when the priors on its coefficients and on $\beta$ are coupled. Under $p(\beta, \tau, \sigma^2) \propto 1/\sigma^2$ the reduction goes further. With $M$ the projection off $\Phi$, the posterior of $\tau$ is Student $t$ with location $T^\top MY/T^\top MT$, squared scale $\mathrm{RSS}/\{(n - \dim\beta - 1)T^\top MT\}$ and $n - \dim\beta - 1$ degrees of freedom, and so it depends on the data only through three numbers, of which $(\hat\tau, \hat J, \hat\sigma^2)$ are the cross fitted analogues.

Second, why the three numbers stay nearly sufficient with a nonparametric $g$, as a first order argument and not a theorem. Fix $\ell_0=\E(Y\mid X)$ and $m_0=\E(T\mid X)$, the least favourable direction. The Gaussian likelihood is proportional to $\exp\{-\sum_i(U_i-\tau V_i)^2/(2\sigma^2)\}$, with sufficient statistics $(\sum_iV_iU_i,\sum_iV_i^2)$ and, under a flat prior, posterior $N(\tilde\tau,\sigma^2/\sum_iV_i^2)$, where $\tilde\tau$ is the oracle estimator. Chapter~\ref{ch:dml}'s rate conditions give $\sqrt n(\hat\tau-\tilde\tau)\to0$ in probability, and $\hat\tau$ reaches the efficiency bound $\sigma^2/\E V^2$, so $(\hat\tau,\hat J,\hat\sigma^2)$ carry this posterior's location and scale. Proposition~\ref{blp:prop:bvm} cannot establish sufficiency here: it needs a fixed dimensional parameter and an injective summary limit of full column rank, which unrestricted nuisances do not provide. Sufficiency must therefore be measured. With heteroskedastic errors $\hat\tau$ is inefficient; $\hat\Omega$ supplies only its variance $\hat\Omega/(n\hat J^2)$, and the summary also needs a weighted estimate.

Third, the score bias $\E(\Delta_m\Delta_\ell)-\tau\E(\Delta_m^2)$ (Section~\ref{sec:baipw}) can in a finite sample be of the same order as the standard error and is invisible in $(\hat\tau,\hat J,\hat\sigma^2)$. The other four components are noisy proxies: nuisance $R^2$ values measure fit, while gaps $\hat\tau_{\mathrm{plug}}-\hat\tau$ and $\hat\tau_{\mathrm{ols}}-\hat\tau$ signal regularization or basis error. Simulations can teach the network the resulting shift and extra width. We have no theorem for this approximation.

The loss can be measured. Since $S$ is a function of $D$,
\[
\E\,\mathrm{Var}(\tau \mid S) - \E\,\mathrm{Var}(\tau \mid D) = \E\,\mathrm{Var}\{\E(\tau \mid D) \mid S\} \geq 0,
\]
so a summary is nearly sufficient for the mean exactly when $\E(\tau \mid D)$ is nearly a function of $S$. In the sparse linear model the Gibbs sampler of Section~\ref{sec:hsgibbs} supplies $\E(\tau \mid D)$, and the regression of it on $S$ estimates this gap. In general, compare held out check loss with and without each component (Section~\ref{comp:sec:checks}): a component that lowers it was not redundant. Rank histograms cannot serve here, because they are flat for any $S$ (Proposition~\ref{comp:prop:ranks}). Nor can a summary repair the model: an interaction of $T$ and $X$ in the simulator changes the meaning of $\tau$, and non Gaussian errors or a model outside the prior break the second step.

\paragraph{Exercise~\ref{ex:ch14:5}.} The design matrix $Z$ is fixed once (seed 20261105): $n = 100$, 50 ones and 50 zeros in the treatment column and 19 independent standard normal columns, so $p = 20$. Each of the 1,000 replications draws $\sigma^2 \sim \mathrm{IG}(2, 1)$, $\tau \mid \sigma^2 \sim N(0, 100\sigma^2)$, $\gscale$ and $\lambda_j$ standard half Cauchy, $\beta_j \mid \lambda_j, \gscale, \sigma^2 \sim N(0, \lambda_j^2\gscale^2\sigma^2)$ for $j = 2, \dots, 20$, and then $Y$ from the regression with error variance $\sigma^2$. Let $\Lambda$ be diagonal with first entry 100 and the others $\lambda_j^2\gscale^2$, and $A = Z^\top Z + \Lambda^{-1}$. The sampler of Section~\ref{sec:hsgibbs}, changed only by the proper priors, draws the coefficients from $N(A^{-1}Z^\top Y, \sigma^2A^{-1})$. With $b = (\tau, \beta^\top)^\top$ it then draws
\[
\sigma^2 \sim \mathrm{IG}\Bigl(2 + \frac{n + p}{2},\ 1 + \frac12\bigl\{\|Y - Zb\|^2 + b^\top\Lambda^{-1}b\bigr\}\Bigr),
\]
which has the prior shape and scale added to those of Exercise~\ref{ex:ch07:4}. The local scales and their auxiliary variables follow as in that exercise, and $\gscale^2$ has shape $(19 + 1)/2 = 10$ with its sum over the 19 horseshoe coefficients only, so $\tau$ does not enter its update. Some prior draws are extreme (the largest $Y^\top Y$ is near $10^{12}$), so the factorization uses $A^{-1} = \Lambda^{1/2}(I + \Lambda^{1/2}Z^\top Z\Lambda^{1/2})^{-1}\Lambda^{1/2}$ and the scales are bounded to $[10^{-30}, 10^{30}]$.

The spacing comes from a pilot on 24 prior draws, 20,000 sweeps each after 2,000 discarded. Mean lag 1 autocorrelation was 0.013 for $\tau$, nearly independent of the scales given the data. For $\log\gscale$, means at lags 1, 10, 50 and 100 were 0.85, 0.23, $-0.01$, $-0.01$; worst values at lags 10, 50 and 100 were 0.51, 0.08, 0.04. We discard 5,000 sweeps and keep 19 draws 100 apart. Ranking the 19 draws within each replication and pooling over replications, the lag 1 autocorrelation of the ranks is $-0.060$ for $\tau$ and $-0.049$ for $\gscale$, against $-1/18=-0.056$ for independent draws.

The histograms count $B$, the number of kept draws below the generating value (Proposition~\ref{comp:prop:ranks}), for $B = 0, \dots, 19$, with 50 expected in each bin:
\begin{align*}
\tau:\ & 45, 51, 58, 51, 44, 55, 47, 65, 52, 61, 44, 51, 42, 64, 55, 41, 43, 49, 39, 43, \\
\gscale:\ & 51, 60, 61, 40, 52, 52, 55, 48, 57, 50, 44, 57, 45, 36, 49, 42, 52, 50, 52, 47.
\end{align*}
Against the flat histogram the chi squared statistics on 19 degrees of freedom are 22.7 ($p = 0.25$) for $\tau$ and 15.9 ($p = 0.66$) for $\gscale$. In bins of four ranks the counts are 205, 211, 208, 202, 174 and 212, 207, 208, 172, 201 (expected 200), with $p = 0.35$ and $p = 0.27$. Neither is U shaped, humped or tilted beyond chance, so the sampler passes at this sample size. The Monte Carlo standard error of a count in a bin of 50 is about 7, so the test would miss a small error. A further 3,000 replications, 4,000 in all, also pass: the statistics are 23.7 ($p = 0.21$) and 12.4 ($p = 0.87$).

To see that the check has power against dependence, keep 19 consecutive sweeps instead (thinning of one, with 1,000 sweeps discarded). The histogram of $\gscale$ turns U shaped, with counts 202 and 153 at the two ends, a chi squared of 763 and a pooled rank autocorrelation of 0.60. That of $\tau$ stays flat, with a chi squared of 14.7 ($p = 0.74$), as its autocorrelation predicts. A dependent chain is too narrow locally and looks like an overconfident posterior, and only the slower variable shows it, which is why the spacing has to be set by $\gscale$.

\paragraph{Exercise~\ref{ex:chgbc:3}.} For level $u$ and summary $s=\bar Y$, the posterior $N(\mu_N(s),\sigma_N^2)$ has quantile $q=\mu_N(s)+\sigma_N\Phi^{-1}(u)$, with $\Phi$ the standard normal distribution function. This minimizes expected check loss by Proposition~\ref{gbc:prop:pinball} (or Exercise~\ref{ex:ch08:2}). The Bayes action is $a^\ast = \mu_N$, and $(a^\ast - \theta)^2 = \sigma_N^2Z^2$ with $Z$ standard normal. The loss is in its upper 10 percent exactly when $|Z| > c = \Phi^{-1}(0.95) = 1.645$. Since $\int_c^\infty z^2\varphi(z)\,dz = c\varphi(c) + 1 - \Phi(c)$ by parts, with $\varphi$ the standard normal density,
\[
\E\{Z^2 \mid |Z| > c\} = \frac{2c\varphi(c) + 0.1}{0.1} = 1 + 20\,c\,\varphi(c) = 4.39.
\]
The mean of the upper 10 percent of the loss is $4.39\,\sigma_N^2$, so the $\CVaR$ at level $0.1$ of the utility $-(a^\ast - \theta)^2$ is $-4.39\,\sigma_N^2$. That is 4.4 times the expected loss $\sigma_N^2$ and 1.6 times the 0.9 quantile of the loss, $2.71\,\sigma_N^2$.

\paragraph{Exercise~\ref{ex:ch08:1}.} Let $F(y \mid x) = \Prob(Y \le y \mid X = x)$ be a regular conditional distribution function, $F_-(y \mid x) = \Prob(Y < y \mid X = x)$, and $G(u, x) = \inf\{y : F(y \mid x) \ge u\}$. Right continuity gives $G(u, x) \le y$ exactly when $u \le F(y \mid x)$, so $G$ is measurable. Enlarge the space by a uniform $V$ independent of $(X, Y)$ and set
\[
U = F_-(Y \mid X) + V\{F(Y \mid X) - F_-(Y \mid X)\}.
\]
Fix $x$ and $u \in (0, 1)$ and let $y_0 = G(u, x)$, so $F_-(y_0 \mid x) \le u \le F(y_0 \mid x)$. The cases $Y < y_0$, $Y = y_0$ and $Y > y_0$ contribute $F_-(y_0 \mid x)$, $u - F_-(y_0 \mid x)$ and zero to $\Prob(U \le u \mid X = x)$, with equality boundaries of probability zero. Thus
\[
\Prob(U \le u \mid X = x) = u,
\]
so $U$ is uniform and independent of $X$. The inverse identity gives $\Prob\{G(U, X) \le y \mid X = x\} = F(y \mid x)$. Also $U \le F(Y \mid X)$ gives $G(U, X) \le Y$. These variables have the same conditional law and are ordered, so $G(U, X) = Y$ almost surely.

\paragraph{Exercise~\ref{ex:ch08:2}.} Let $F$ be the distribution function of $Y$, with $\E|Y| < \infty$, and $g(q) = \E\rho_u(Y - q)$, where $\rho_u(r) = r\{u - \one(r < 0)\}$. The check loss is convex in $q$, with slopes $-u$ for $q < Y$ and $1 - u$ for $q > Y$. Its difference quotients are bounded by one, so the one sided derivatives pass under the expectation:
\[
g'_-(q) = F(q-) - u, \qquad g'_+(q) = F(q) - u.
\]
Thus $g$ is minimized exactly when $F(q-) \le u \le F(q)$, the quantile condition. The smallest minimizer is $Q(u) = \inf\{q : F(q) \ge u\}$; it is unique when the $u$ quantile is, for example when $F$ is strictly increasing. Conditioning on $S = s$ gives Proposition~\ref{gbc:prop:pinball}.

\paragraph{Exercise~\ref{ex:chcomp:1}.} The prior $N(0, 1)$ is the exact posterior given a constant summary, so Proposition~\ref{comp:prop:ranks} applies: $\theta$ and the $L$ network draws are independent from $N(0, 1)$, hence exchangeable, and the rank $B$ of $\theta$ is uniform on $\{0, \dots, L\}$. That is the global check. In the local check the generating values are drawn from the network's posterior at the observed $s$, which is $N(0, 1)$ whatever $s$ is. The data simulated from each value are ignored by the network, which returns fresh $N(0, 1)$ draws, so each rank is again that of one value among $L + 1$ exchangeable ones: uniform at every $s$.

At $u = 1/2$ the check loss is $\rho_{1/2}(r) = |r|/2$. The prior network has expected loss $\frac12\E|\theta| = \frac12\sqrt{2/\pi} = 0.399$. The exact posterior $N(10\bar y/11, 1/11)$ has its median at its mean, so its expected loss is $\frac12\cdot\frac{1}{\sqrt{11}}\sqrt{2/\pi} = 0.120$ at every $\bar y$. The ratio is $\sqrt{11} = 3.32$, the ratio of the standard deviations. 

The held out check loss and the reference check separate the networks: the loss gap is large, and at $\bar y = 1.5$ the prior network has median 0 against the exact 1.36. The global rank check does not: the histograms of both networks are flat. The local rank check does, the wrong way round, passing the useless network and failing the exact posterior. The local histogram of the prior network is flat at every $s$, and that of the exact posterior is not, because its generating values come from the posterior at $s$ and not from the prior: at $\bar y = 1.5$ about two thirds of its ranks fall in the upper half.

\paragraph{Exercise~\ref{ex:chcomp:2}.} For any pair, $\mathrm{Var}(\mu_1 - \mu_0) = \mathrm{Var}\,\mu_1 + \mathrm{Var}\,\mu_0 - 2\Cov(\mu_1, \mu_0)$, so joint draws give $v_1 + v_0 - 2c$. Independent marginal draws set the covariance to zero and give $v_1 + v_0$, overstating the variance by $2c$ when $c > 0$ and understating it when $c < 0$. With $v_1 = v_0 = 1$ and $c = 1/2$ the variance is 2 instead of 1, and the interval is $\sqrt2$ times too wide.

The triangular map of Section~\ref{comp:sec:noise} draws $\mu_1 = G_1(U_1, s)$ and then $\mu_0 = G_0(U_2, s, \mu_1)$ from its conditional quantile function, with independent uniforms. The second network sees the first draw, so the pair has the joint law, covariance $c$ and effect variance $v_1 + v_0 - 2c$.

\paragraph{Exercise~\ref{ex:chcomp:3}.} (a) With $C_{\mathrm{net}}$ negligible, the rule of Section~\ref{comp:sec:algorithm} gives
\[
n_{\mathrm{sets}} > \frac{C_{\mathrm{train}}}{C_{\mathrm{mcmc}}} = \frac{36{,}000}{0.081} \approx 444{,}000
\]
data sets, which the section rounds to 440,000. Training took at least $36{,}000$ seconds, so this is a lower bound on the break-even.

(b) Including checks, the network is cheaper when $C_{\mathrm{train}} + n_{\mathrm{sets}}(C_{\mathrm{net}} + C_{\mathrm{chk}}) < n_{\mathrm{sets}}C_{\mathrm{mcmc}}$, or
\[
n_{\mathrm{sets}} > \frac{C_{\mathrm{train}}}{C_{\mathrm{mcmc}} - C_{\mathrm{net}} - C_{\mathrm{chk}}},
\qquad C_{\mathrm{chk}} < C_{\mathrm{mcmc}} - C_{\mathrm{net}}.
\]
If the denominator is zero or negative, training is never repaid. The upper bound on the check cost is strict: it must cost less than the sampling time the forward pass saves.

\paragraph{Exercise~\ref{ex:chcomp:4}.} For integrable $q$, the identity $\rho_u(y - q) = (q - y)_+ - u(q - y)$ gives
\[
\int_0^1 \rho_u\{y - q(u)\}\,du = \int_0^1 \{q(u) - y\}_+\,du - \int_0^1 u\,q(u)\,du + \frac y2.
\]
The first term depends only on the law of $q(V)$ for $V$ uniform, so rearrangement leaves it unchanged. Hence, by the inequality of Hardy, Littlewood and P\'olya,
\[
\int_0^1 \bigl[\rho_u\{y - q^\ast(u)\} - \rho_u\{y - q(u)\}\bigr]\,du
= -\int_0^1 u\{q^\ast(u) - q(u)\}\,du \le 0.
\]
The gain does not depend on $y$. Taking $q = f_\phi(\cdot, s)$ and $y = \theta$ and averaging over $(\theta, S)$ shows that sorting never raises the network's expected check loss.

\section*{Chapter~\ref{ch:gente}: Generative Treatment Effects and Joint Potential Outcomes}

\paragraph{Exercise~\ref{ex:ch08:3}.} Fix $x$ and drop it from the notation. The difference $\Delta(u) = q_1(u) - q_0(u)$ is a function of the two marginal laws, while the law of $\eta = Y(1) - Y(0)$ also depends on the coupling of $Y(0)$ and $Y(1)$, which the marginals do not determine (Section~\ref{sec:ot}). Let $Y(0)$ be uniform on $\{0, 1\}$ and $Y(1) = 1 - Y(0)$. Both marginals are uniform on $\{0, 1\}$, so $q_0(u) = q_1(u) = \mathbf{1}(u > 1/2)$ and $\Delta(u) = 0$ for every $u$. But $\eta = 1 - 2Y(0)$ is $+1$ or $-1$ with probability one half each, so its $u$ quantile is $-1$ for $u \le 1/2$ and $+1$ for $u > 1/2$. The quantile effect curve is identically zero while every unit moves by one. 

\paragraph{Exercise~\ref{ex:ch08:4}.} In the simulation $Y(0) = e^{Z/2}$ with $Z \sim N(0,1)$, so the control median is 1, and $Y(1) = 1.3\,Y(0)$ if $Y(0) > 1$ and $Y(1) = Y(0)$ otherwise. The gain is $0.3\,Y(0)$ above the median and zero below, and $\E\{e^{\sigma Z}\mathbf{1}(Z > 0)\} = e^{\sigma^2/2}\Phi(\sigma)$, so with $\sigma = 1/2$
\[
\tau = 0.3\,\E\{Y(0)\mathbf{1}(Y(0) > 1)\} = 0.3\,e^{1/8}\,\Phi(1/2) = 0.2351.
\]
The map $y \mapsto y\{1 + 0.3\cdot\mathbf{1}(y > 1)\}$ is nondecreasing, so $q_1(u) = q_0(u)\{1 + 0.3\cdot\mathbf{1}(u > 1/2)\}$ and the true quantile effect curve is $\Delta(u) = 0$ for $u \le 1/2$ and $\Delta(u) = 0.3\,e^{\Phi^{-1}(u)/2}$ above, which is 0.30 just above the median, 0.42 at $u = 0.75$ and 0.68 at $u = 0.95$. Its area is
\[
\int_0^1 \Delta(u)\,du = 0.3\int_0^\infty e^{z/2}\varphi(z)\,dz = 0.3\,e^{1/8}\,\Phi(1/2) = 0.2351,
\]
the same number, because a mean is the integral of its quantile function and $\tau = \int q_1 - \int q_0$. The figure shows only the levels 0.05 to 0.95. The area over that range is 0.1922, and the other 0.0429 comes from the top 5\% of levels, where the curve rises from 0.68 and the lognormal control quantiles are large. In a trial of 1{,}500 units the area under the empirical quantile effect curve over $(0, 1)$ is the difference in arm means, which in 3{,}000 simulated trials averaged 0.236 with standard deviation 0.037.

\paragraph{Exercise~\ref{ex:ch08:5}.} Let $y_1, \ldots, y_n$ be the data, $a$ the concentration and $G_0$ the base law, with $\int |y|\,dG_0 < \infty$. The Dirichlet process posterior of Section~\ref{app:toolkit} has the representation
\[
P = V\,Q + (1 - V)\sum_{i=1}^n W_i\,\delta_{y_i},
\]
where $V \sim \mathrm{Beta}(a, n)$, $Q \sim \mathrm{DP}(a, G_0)$ and $W \sim \mathrm{Dirichlet}(1, \ldots, 1)$ are independent. On any partition with cell counts $n_j$, beta and Dirichlet aggregation give cell masses $\mathrm{Dirichlet}\{aG_0(A_j) + n_j\}$, as required.

Write $m_Q = \int y\,dQ$. Since $\E|m_Q| \le \int |y|\,dG_0$ uniformly in $a$ and $\E V = a/(a + n) \to 0$, both $V m_Q$ and $V$ tend to zero in probability. Thus
\[
\int y\,dP = V m_Q + (1 - V)\sum_i W_i y_i \ \Rightarrow\ \sum_i W_i y_i,
\]
the Bayesian bootstrap posterior of the mean. The cell laws likewise tend to $\mathrm{Dirichlet}(n_1, \ldots, n_k)$.

\paragraph{Exercise~\ref{ex:ch08:6}.} Take $Y > 0$ and a cost $k$ in utility units. Treat a unit with covariates $x$ when
\[
\int \log y\,dF_1(y \mid x) - \int \log y\,dF_0(y \mid x) = \int_0^1 \log\frac{q_1(u \mid x)}{q_0(u \mid x)}\,du > k,
\]
that is, when the geometric mean of $Y(1)$ given $x$ exceeds that of $Y(0)$ by the factor $e^k$. Linear utility treats when $\tau(x) > k$. Write $\mu_t(x) = \E\{Y(t) \mid x\}$ and $J_t(x) = \log\mu_t(x) - \E\{\log Y(t) \mid x\}$, which is nonnegative by Jensen's inequality. The log rule is
\[
\log\frac{\mu_0(x) + \tau(x)}{\mu_0(x)} - \{J_1(x) - J_0(x)\} > k,
\]
the relative effect on the mean, corrected by the change in the spread of the outcome. In the integral a given gain $\Delta(u \mid x)$ is worth more where $q_0$ is small, and a loss that drives $q_1$ toward zero costs without bound, so the lower tail carries the weight. The rules differ when the spread changes. Let $Y(0) = 1$ and $Y(1) \in \{0.1, 2.9\}$ with probability one half each: $\tau = 0.5$, but $\E\log Y(1) - \E\log Y(0) = \tfrac12\log 0.29 = -0.62$, so at $k = 0$ the linear rule treats and the log rule does not. 

\paragraph{Exercise~\ref{ex:ch08:7}.} Centre the covariates, so that $\tau_0$ is the average effect, and write
\[
Y_i = \mu(x_i) + T_i(\tau_0 + x_i^\top\delta) + \varepsilon_i, \qquad \varepsilon_i \sim N(0, \sigma^2),
\]
\[
\delta_j \mid \lambda_j, \gscale, \sigma \sim N(0, \sigma^2\lambda_j^2\gscale^2), \qquad \lambda_j \sim C^+(0, 1), \qquad \gscale \sim C^+(0, 1),
\]
with a flat prior on $\tau_0$ and any prior on the baseline $\mu$, for example $\alpha + x^\top\beta$ with a horseshoe of its own. The modifiers have their own scale $\gscale$, which shrinks the effect toward homogeneity. Given $\mu$, $\tau_0$ and the scales, let $r_i = Y_i - \mu(x_i) - T_i\tau_0$ and let $W$ have rows $T_i x_i^\top$. Completing the square with $A = W^\top W + \mathrm{diag}\{(\lambda_j^2\gscale^2)^{-1}\}$ gives
\[
\delta \mid \text{rest} \sim N(A^{-1}W^\top r, \sigma^2A^{-1}).
\]
Since $T_i^2 = T_i$, $W^\top W = \sum_{i:T_i = 1}x_i x_i^\top$ and $W^\top r = \sum_{i:T_i = 1}x_i r_i$: given the baseline, only treated units inform the modifiers. Adding the column $T_i$ with prior precision zero draws $(\tau_0, \delta)$ as one block. The scale updates are those of Section~\ref{sec:hsregression} with $\delta$ in place of $\beta$.

\paragraph{Exercise~\ref{ex:ch09:1}.} With $F_0 = \Phi$ and $F_1 = \Phi(\cdot - 1)$, the difference $F_0(y) - F_1(y) = \Phi(y) - \Phi(y - 1)$ is maximized at $y = 1/2$ with value $\Phi(0.5) - \Phi(-0.5) \approx 0.383$. Hence, for the gain $\eta = Y(1) - Y(0)$ of the Makarov theorem, $F_\eta(0) \leq 1 - 0.383$ and $\Prob(\eta > 0) \geq 0.383$, while the lower bound on $F_\eta(0)$ is zero, so $\Prob(\eta > 0) \leq 1$.

\paragraph{Exercise~\ref{ex:ch09:2}.} Let $U$ be uniform on $(0, 1)$. For any distribution function $F$, $F^{-1}(u) \le y$ if and only if $u \le F(y)$, so
\begin{align*}
\Prob\{F_0^{-1}(U) \le y_0,\ F_1^{-1}(U) \le y_1\} &= \Prob\{U \le F_0(y_0),\ U \le F_1(y_1)\}\\
&= \min\{F_0(y_0), F_1(y_1)\}.
\end{align*}
The right side is the distribution function of the comonotone coupling, the upper bound in the Fr\'echet and Hoeffding theorem, and a distribution function determines the law. So under the comonotone coupling $(Y(0), Y(1))$ has the law of $(F_0^{-1}(U), F_1^{-1}(U))$, and the gain $\eta$ has the law of $F_1^{-1}(U) - F_0^{-1}(U) = \Delta(U)$. 

\paragraph{Exercise~\ref{ex:ch09:3}.} If $(Y(0), Y(1))$ is bivariate normal with unit variances, means $0$ and $s$, and correlation $\rho$, then $Y(1) - Y(0) \sim N(s, 2(1 - \rho))$ and $\Prob\{Y(1) > Y(0)\} = \Phi\bigl(s/\sqrt{2(1 - \rho)}\bigr)$, which tends to one as $\rho \to 1$ for $s > 0$ and equals $\Phi(s/\sqrt{2})$ at $\rho = 0$.

\paragraph{Exercise~\ref{ex:ch09:4}.} For $y_0 \le y_0'$ and $y_1 \le y_1'$, submodularity is
\[
h(y_1 - y_0) + h(y_1' - y_0') \le h(y_1' - y_0) + h(y_1 - y_0').
\]
Let $a = y_1 - y_0'$ and $b = y_1' - y_0$. The sorted differences lie in $[a, b]$ and sum to $a + b$, so they are $(1 - t)a + tb$ and $ta + (1 - t)b$ for some $t \in [0, 1]$. Adding their convexity bounds gives the inequality.

For laws uniform on $n$ sorted points $x_i$ and $z_i$, Birkhoff's theorem gives an optimal pairing $x_i \leftrightarrow z_{\sigma(i)}$. Swapping any inverted pair does not raise the cost, by submodularity. Finitely many such exchanges give the sorted pairing.

For general laws with finite means, subtract a supporting line of $h$ to make $h \ge 0$ without changing which coupling is optimal. Draw independent pairs $(X_i, Z_i)$ from any coupling $H$ of finite cost; infinite cost needs no proof. With $F_{t,n}$ the empirical distribution functions, the finite case gives
\[
\int_0^1 h\{F_{1,n}^{-1}(u) - F_{0,n}^{-1}(u)\}\,du \le \frac1n\sum_{i=1}^n h(Z_i - X_i).
\]
The right side tends to $\int h\,dH$ by the law of large numbers. The empirical quantiles converge at all but countably many $u$ by Glivenko and Cantelli. Fatou's lemma therefore bounds the cost of the comonotone coupling, $\int_0^1 h\{F_1^{-1}(u) - F_0^{-1}(u)\}\,du$, by $\int h\,dH$.

\paragraph{Exercise~\ref{ex:ch09:5}.} In a crossover trial each subject receives both treatments in turn. With no carryover the second period is unaffected by the first treatment, and with no period effect the outcome under a treatment does not depend on the period, so whatever the order every subject contributes the pair $(Y_i(0), Y_i(1))$, and the pairs are an independent sample from the joint law $H$. A parallel group trial observes one coordinate per subject and cannot do this. The empirical distribution function of the pairs converges to $H$, so $H$ is identified, and so is its copula $C(u, v) = H\{F_0^{-1}(u), F_1^{-1}(v)\}$. With continuous marginals and a Gaussian copula, $Z_0 = \Phi^{-1}\{F_0(Y(0))\}$ and $Z_1 = \Phi^{-1}\{F_1(Y(1))\}$ are standard normal with correlation $\rho$, so
\[
\rho = \mathrm{Corr}(Z_0, Z_1) = 2\sin(\pi r_S/6),
\]
where $r_S = \mathrm{Corr}\{F_0(Y(0)), F_1(Y(1))\}$ is the Spearman correlation of the pair. The second form holds because a Gaussian copula has $r_S = (6/\pi)\arcsin(\rho/2)$, which is increasing in $\rho$ and so can be inverted. Both expressions are functionals of $H$ and are estimated from the pairs, by the within subject correlation of normal scores or of ranks. With normal marginals $\rho$ is the ordinary within subject correlation of the two outcomes. With skewed marginals it is not: for a lognormal $Y(0)$ with log standard deviation 0.5, an exponential $Y(1)$ and $\rho = 0.6$, the Pearson correlation of the raw outcomes is 0.56 (two million simulated pairs), while the rank formula returns 0.60. With 50 subjects the rank based estimate has standard deviation about 0.10 (4{,}000 simulated trials).

\paragraph{Exercise~\ref{ex:ch09:6}.} At $s = 1$ the fraction helped is $\psi(\rho) = \Phi\{1/\sqrt{2 - 2\rho}\}$, increasing in $\rho$, and the limiting posterior is the law of $\psi(\rho)$ with $\rho$ uniform on $[0, 1]$. Its support is $[\psi(0), 1) = [0.760, 1)$ and its median is $\psi(1/2) = \Phi(1) = 0.841$. Solving $p = \psi(\rho)$ gives $\rho = 1 - 1/(2z_p^2)$ with $z_p = \Phi^{-1}(p)$, so
\[
\Prob(\psi \le p) = 1 - \frac{1}{2\{\Phi^{-1}(p)\}^2}, \qquad 0.760 \le p < 1,
\]
and the $q$ quantile is $\Phi\{1/\sqrt{2(1 - q)}\}$. The quartiles are 0.793 and 0.921, the central 95\% interval is $[0.763, 1.000]$, the mean is 0.860 and the standard deviation 0.076, and the probability that more than 90\% of units are helped is 0.304. The Makarov interval is $[2\Phi(1/2) - 1, 1] = [0.383, 1]$ (Exercise~\ref{ex:ch09:1}), of width 0.617. The support of the limiting posterior has width 0.240, which is 39\% of that, and its interquartile range 0.128, which is 21\%. The part below 0.760 is excluded by the prior, because $\rho \ge 0$, and a Gaussian copula never reaches the Makarov lower bound: even $\rho = -1$ gives $\Phi(1/2) = 0.691$. The width that remains comes from $\rho$ alone and does not shrink as $n$ grows.

\section*{Chapter~\ref{ch:bart}: Bayesian Nonparametrics for Causal Inference}

\paragraph{Exercise~\ref{ex:chbart:1}.} Let $p_d = \alpha(1+d)^{-\beta}$, so $p_0 = 0.95$, $p_1 = 0.2375$ and $p_2 = 0.1056$ (Figure~\ref{fig:bartdepth}). Assuming independent node decisions and an available splitting rule at every node, count the tree shapes:
\begin{align*}
\Prob(1) &= 1 - p_0 = 0.05, \\
\Prob(2) &= p_0(1 - p_1)^2 = 0.95 \times 0.7625^2 = 0.552, \\
\Prob(3) &= p_0\cdot 2p_1(1 - p_1)(1 - p_2)^2 \\
&= 0.95 \times 2 \times 0.2375 \times 0.7625 \times 0.8944^2 = 0.275.
\end{align*}
The three sum to 0.878, so four or more terminal nodes have probability 0.122, and two or three together have probability 0.828, which supports the statement that most trees have two or three terminal nodes.

\paragraph{Exercise~\ref{ex:chbart:2}.} Fix the tree structures. At a given $x$ tree $j$ returns the value $\mu_j$ of the one leaf into which $x$ falls, so by \eqref{eq:bart} the regression function is $f(x) = \sum_{j=1}^m \mu_j$. Leaf values are independent $N(0, \sigma_\mu^2)$ draws whatever the structure, and different trees have different leaves, so the $\mu_j$ are independent and
\[
f(x) \mid \mathcal{T}_1, \dots, \mathcal{T}_m \sim N(0, m\sigma_\mu^2), \qquad m\sigma_\mu^2 = m\cdot\frac{0.25}{k^2 m} = \Bigl(\frac{0.5}{k}\Bigr)^2.
\]
This law does not depend on the structures, so it is also the marginal law of $f(x)$: the standard deviation is $0.5/k$ at every $x$ and for every $m$. 

\paragraph{Exercise~\ref{ex:chbart:3}.} Hold the other trees and $\sigma^2$ fixed. In leaf $\ell$, $R_{j,i} = \mu_\ell + \varepsilon_i$ with independent $\varepsilon_i \sim N(0, \sigma^2)$ and independent priors $\mu_\ell \sim N(0, \sigma_\mu^2)$. Let $n_\ell$ be its size and $S_\ell = \sum_{i \in \ell}R_{j,i}$. The conditional kernel is
\[
p(\mu_\ell \mid \cdot) \propto \exp\Bigl\{-\frac12\Bigl(\frac{n_\ell}{\sigma^2} + \frac1{\sigma_\mu^2}\Bigr)\mu_\ell^2 + \frac{S_\ell}{\sigma^2}\mu_\ell\Bigr\}.
\]
Completing the square gives, independently over leaves,
\[
\mu_\ell \mid \mathcal{T}_j, \sigma^2, R_j \sim N\Bigl(\frac{\sigma_\mu^2 S_\ell}{\sigma^2 + n_\ell\sigma_\mu^2},\; \frac{\sigma^2\sigma_\mu^2}{\sigma^2 + n_\ell\sigma_\mu^2}\Bigr).
\]
The mean shrinks the leaf average by $n_\ell\sigma_\mu^2/(\sigma^2 + n_\ell\sigma_\mu^2)$.

\paragraph{Exercise~\ref{ex:chbart:4}.} Fix the tree structures and the variances, so that the sum of trees is a linear model in its leaf values and the argument of Chapter~\ref{ch:horseshoe} can be run exactly. Suppose $K$ trees split on $T$ alone and the others only on $X$. Treatment tree $k$, with independent leaf values $c_{0k}, c_{1k} \sim N(0, \sigma_\mu^2)$, contributes $a_k/\sqrt 2 + d_kD$ with $D = T - 1/2$, $a_k = (c_{0k} + c_{1k})/\sqrt 2$ and $d_k = c_{1k} - c_{0k}$, which are independent with variances $\sigma_\mu^2$ and $2\sigma_\mu^2$. Collect in $B$ the leaf indicators of the other trees and one constant column $1/\sqrt 2$ for each treatment tree, so that its coefficients have the prior $\theta \sim N(0, \sigma_\mu^2I)$. Then $\tau = \sum_kd_k$ has the prior $N(0, w)$ with $w = 2K\sigma_\mu^2$, independent of $\theta$; plain BART has no flat prior on the effect. Let the data be $Y = \tau D + \mu + \varepsilon$ for any surface $\mu = \mu(X)$, and write $H = B(B^\top B + \sigma^2\sigma_\mu^{-2}I)^{-1}B^\top$ for the ridge smoother of this design. Eliminating $\theta$ from the normal equations gives
\[
\hat\tau = \frac{D^\top(I - H)Y}{D^\top(I - H)D + \sigma^2/w},
\]
and hence, exactly,
\[
\hat\tau - \tau = \frac{D^\top(I - H)\mu + D^\top(I - H)\varepsilon - (\sigma^2/w)\tau}{D^\top(I - H)D + \sigma^2/w}.
\]
The first numerator term is regularization induced confounding: the formula of Chapter~\ref{ch:horseshoe} with shrinkage error $\mu - H\mu$. It vanishes under least squares when $\mu$ lies in the span of $B$. The second term is mean zero noise; the third shrinks $\tau$ toward zero and can have either sign relative to the first.

Leaf shrinkage, with $\sigma_\mu$ falling like $m^{-1/2}$, and the depth prior $\alpha(1 + d)^{-\beta}$ leave peaks and troughs of $\mu$ unfitted. Under targeted selection, treated units have high $\mu$, so $D^\top(I - H)\mu > 0$ and the bias is upward. Averaging over structures does not remove this and adds another route: a split on $T$ competes equally with splits on $X$ and reduces residual variance by about $\mathrm{Cov}(T, \mu - H\mu)^2/\mathrm{Var}(T)$. Where several splits on $X$ would be needed for the same fit, the depth prior and marginal likelihood favour $T$, and part of $\mu$ is reported as effect.

\paragraph{Exercise~\ref{ex:chbart:5}.} With no splits, the effect forest sums independent normal leaf constants, so $\tau(x) = \tau \sim N(0, \sigma_\tau^2)$, where $\sigma_\tau^2$ is the sum of their variances. Since the estimated propensity is a fixed function of $x$, write $g(x) = \mu\bigl(x, \hat e(x)\bigr)$. Then \eqref{eq:bcf} becomes
\[
Y_i = \tau T_i + g(x_i) + \varepsilon_i, \qquad \tau \sim N(0, \sigma_\tau^2), \quad g \sim \text{BART},
\]
the partially linear model of Chapter~\ref{ch:dml}. Its nuisance trees may split on $\hat e(x)$ as well as on $x$.

\paragraph{Exercise~\ref{ex:chbart:6}.} The two rules look at different things, so build a design with one region that fools each. Take $n = 1{,}000$ and five independent standard normal covariates. The propensity is $\expit\{1.2\min(x_1, 1) + 2.3\,\mathbf{1}(x_1 > 1)\}$, which is 0.97 for $x_1 > 1$ (region F), except that it is 0.995 where $x_2 > 1.2$ (region P). The outcome is $Y = \mu(x) + \tau(x)T + \varepsilon$ with standard normal $\varepsilon$, $\mu(x) = \min(x_1, 1) + 2\max(x_2 - 1.2, 0) + 0.5x_3$ and $\tau(x) = 1 + \mathbf{1}(x_1 > 1) + 6\max(x_2 - 1.2, 0)$. In F the propensity is extreme but the prognostic surface is flat. P has almost no controls (0.6 per data set on average), but its $x_2$ enters the propensity through a threshold, which a linear logit reads as a moderate slope, and both surfaces are steep there. Fit BART with $T$ as a covariate (\texttt{stochtree} defaults, 10 grow from root sweeps and four chains of 2,000 retained draws each, pooled) and a linear logistic $\hat e$ in each of 200 data sets (seed 20261106). There are about 571 treated units per data set: 114 in P, 136 others in F, and 321 elsewhere, against 4.2 controls in F. Five subsets of the treated are compared: all of them; those kept by the one standard deviation rule of the chapter; those kept by the chapter's 99th percentile variant; those with $\hat e \leq 0.9$; and, as a benchmark no analyst has, those with true $e \leq 0.9$. Each estimate is the posterior mean of the average of $f(x_i, 1) - f(x_i, 0)$ over the subset, judged against the true average $\tau(x_i)$ over the same subset.
\begin{center}
\small
\begin{tabular}{lccccccc}
\toprule
Units used & Kept & Truth & Estimate & Bias & RMSE & Cover & Width \\
\midrule
All & 571 & 1.86 & 1.20 & $-0.66$ & 0.68 & 0.015 & 0.49 \\
One s.d.\ rule & 546 & 1.60 & 1.16 & $-0.44$ & 0.48 & 0.150 & 0.46 \\
99th percentile & 426 & 1.24 & 1.10 & $-0.14$ & 0.19 & 0.650 & 0.38 \\
$\hat e \leq 0.9$ & 477 & 1.65 & 1.16 & $-0.49$ & 0.51 & 0.025 & 0.41 \\
True $e \leq 0.9$ & 321 & 1.00 & 1.05 & $0.05$ & 0.10 & 0.885 & 0.33 \\
\bottomrule
\end{tabular}
\end{center}
The Monte Carlo standard error of each bias is at most 0.013. The rules do discard different units. Trimming on $\hat e$ removes 48 percent of F, where the fitted propensity averages 0.88, and only 24 percent of P, where it averages 0.73. The one standard deviation rule removes 0.1 percent of F and 22 percent of P, and the 99th percentile variant 35 percent of F and 79 percent of P. Their mean $s_i^0$ follows the same order: 0.35 elsewhere, 0.47 in F and 0.62 in P. Of the 94 units trimmed and the 25 discarded by the first rule, only 9 are common. With the second rule, 55 units are common, 39 are discarded only by trimming and 90 only by the rule. The first rule is mild because its cutoff, the largest $s_j^1$ among the treated plus a standard deviation, is 0.73, while the 99th percentile cutoff is 0.49.

Discarding moves the target far more than the estimate. The truths range from 1.00 to 1.86 and the estimates only from 1.05 to 1.20. Over the 250 units in F and P the implied BART average is about 1.39 against a truth of about 2.96. In P the effect is not identified without controls, and in F the jump in the treated outcomes at $x_1 = 1$ can be fitted by a split on $x_1$ shared by both arms, so the flat surface does not rescue the counterfactual there: the small $s_i^0$ is a confident extrapolation. The 95\% intervals, from the 2.5 and 97.5 percent posterior quantiles, cover in only 1.5 to 65.0 percent of replications when they keep any part of those units, and the subset closest to the nominal level, at 88.5 percent, is the one no rule found, which drops both regions. The practical advice of the chapter stands, with a warning: report the discard count and how the effect moves, and do not read a small move as evidence of support. Here the estimate moves by 0.04 to 0.15 while the quantity being estimated changes by up to 0.86, and neither rule flags F and P in full.

\paragraph{Exercise~\ref{ex:chbart:7}.} Index the two sides of the cutoff by $s \in \{-, +\}$. Side $s$ has outcomes $y_s$ at running variable values $r_s$, with $y_i = f_s(r_i) + \varepsilon_i$, $\varepsilon_i \sim N(0, \sigma_s^2)$, and the independent priors $f_s \sim \mathrm{GP}(0, k_s)$, $k_s(r, r') = \eta_s^2\exp\{-(r - r')^2/(2\ell_s^2)\}$. Let $K_s$ be the kernel matrix at $r_s$, $k_s(c)$ the vector of $k_s(r_i, c)$ and $A_s = K_s + \sigma_s^2 I$. Squared exponential paths are continuous, so each one sided limit at $c$ is $f_s(c)$ and $\tau(c) = f_+(c) - f_-(c)$. Conditioning each side on its own data gives $f_s(c) \mid y_s \sim N(m_s, v_s)$ with
\[
m_s = k_s(c)^\top A_s^{-1}y_s, \qquad v_s = \eta_s^2 - k_s(c)^\top A_s^{-1}k_s(c).
\]
The sides share neither prior nor data, so these posteriors are independent and
\[
\tau(c) \mid y, \ell_-, \ell_+ \sim N(m_+ - m_-,\; v_+ + v_-).
\]
The weights $k_s(c)^\top A_s^{-1}$ concentrate near $c$ for small $\ell_s$ and spread out as it grows, so $\ell_s$ acts as a bandwidth. Instead of selecting one value, put a prior on $\ell = (\ell_-, \ell_+)$, and the data weigh the length scales through the marginal likelihood of each side. Holding $\eta_s$ and $\sigma_s$ fixed,
\begin{align*}
p\bigl(\tau(c) \mid y\bigr) &= \int N\bigl(\tau(c); m_+ - m_-, v_+ + v_-\bigr)\,p(\ell \mid y)\,d\ell, \\
p(\ell \mid y) &\propto p(\ell)\prod_s N(y_s; 0, A_s).
\end{align*}
The jump variance is the posterior average of $v_+ + v_-$ plus the variance of $m_+ - m_-$ over $\ell$, carrying uncertainty about how far to look from the cutoff (Section~\ref{sec:bart:gp}).

% Chapter 11 (m08.tex)

\section*{Chapter~\ref{ch:iv}: Instruments}

\paragraph{Exercise~\ref{ex:ch10:1}.} The error $\varepsilon$ is independent of $T$, and $\mathrm{Cov}(T, U) = \gamma\,\mathrm{Var}(U)$ because $Z$ and $\nu$ are independent of $U$. The least squares slope therefore has probability limit
\[
\frac{\mathrm{Cov}(T, Y)}{\mathrm{Var}(T)} = \tau + \delta\,\frac{\mathrm{Cov}(T, U)}{\mathrm{Var}(T)} = \tau + \frac{\delta\gamma\,\mathrm{Var}(U)}{\mathrm{Var}(T)},
\]
so the bias is $\delta$ times the slope of the regression of $U$ on $T$. In the design of Figure~\ref{fig:iv}, $\pi = 0.5$, $\gamma = 1$, $\delta = 2$ and $Z$, $U$ and $\nu$ have unit variance, so $\mathrm{Var}(T) = \pi^2 + \gamma^2 + 1 = 2.25$ and the bias is $2/2.25 = 0.889$. The limit $1.889$ matches the simulated mean of 1.89 within Monte Carlo error. The residual variance of the regression of $Y$ on $T$ is $\delta^2\{1 - \gamma^2/\mathrm{Var}(T)\} + 1 = 3.22$, so at $n = 500$ the least squares estimate has standard deviation $\sqrt{3.22/(500 \times 2.25)} = 0.054$, and the mean of 500 replications has standard error 0.002. 

\paragraph{Exercise~\ref{ex:ch10:2}.} Let $A$ and $B$ be the least squares coefficients of $Y$ and of $T$ on $Z$, so the Wald ratio is $A/B$. Conditional on $Z$ and with normal errors they are jointly normal, with $\mathrm{Var}(B) > 0$ and a correlation of absolute value below one (the correlation of the reduced form errors, which is one only if the structural error is a multiple of the first stage error). Write $A = \beta B + W$ with $\beta = \mathrm{Cov}(A, B)/\mathrm{Var}(B)$. Then $W$ is jointly normal with $B$, uncorrelated with it and therefore independent of it, and $\mathrm{Var}(W) > 0$, so
\[
\frac{A}{B} = \beta + \frac{W}{B}, \qquad \E\Bigl|\frac{W}{B}\Bigr| = \E|W|\;\E\frac{1}{|B|}.
\]
The first factor is positive. The density $f_B$ of $B$ is positive at zero, so $f_B \geq f_0 > 0$ on some interval $[-\epsilon, \epsilon]$ and
\[
\E\frac1{|B|} \geq f_0\int_{-\epsilon}^{\epsilon}\frac{db}{|b|} = \infty.
\]
Hence $\E|A/B| = \infty$. Both the positive and the negative part of $W/B$ have infinite expectation, because $W$ and $B$ each take both signs with positive probability, so the mean is undefined. The tail is Cauchy like: $\Prob(|A/B - \beta| > t) \sim 2f_B(0)\,\E|W|/t$ as $t \to \infty$. In the design of Figure~\ref{fig:iv}, $B$ has mean $0.5$ and standard deviation $0.063$, so $\Prob(B < 0) = 1.3 \times 10^{-15}$. The infinite mean is a fact about the far tail and cannot show in 500 replications.

\paragraph{Exercise~\ref{ex:ch10:3}.} The assumptions are $Y(z, t) = Y(t)$ (exclusion), $Z \indep \{T(0), T(1), Y(0), Y(1)\}$ (random assignment), $T(1) \geq T(0)$ (monotonicity) and $\Prob\{T(1) \neq T(0)\} > 0$ (first stage). The observed data are $T = T(Z)$ and $Y = Y(T)$, the second by exclusion. Random assignment removes the conditioning on $Z$:
\[
\E(T \mid Z = z) = \E\{T(z)\}, \qquad \E(Y \mid Z = z) = \E\bigl[Y\{T(z)\}\bigr].
\]
For the numerator use the identity
\[
Y\{T(1)\} - Y\{T(0)\} = \{Y(1) - Y(0)\}\{T(1) - T(0)\},
\]
which holds stratum by stratum. Both sides are zero for always takers and never takers, whose treatment does not depend on $z$ and, by exclusion, whose outcome does not either. Both sides equal $Y(1) - Y(0)$ for compliers and $-\{Y(1) - Y(0)\}$ for defiers. Monotonicity removes the defiers, so $T(1) - T(0)$ is the complier indicator. With $\pi_c$ the complier share and $\tau_c$ the complier effect,
\begin{align*}
\E(Y \mid Z = 1) - \E(Y \mid Z = 0) &= \pi_c\tau_c, \\
\E(T \mid Z = 1) - \E(T \mid Z = 0) &= \pi_c.
\end{align*}
The first stage assumption makes $\pi_c > 0$, so the ratio is $\tau_c$. Without monotonicity the two differences are $\pi_c\tau_c - \pi_d\tau_d$ and $\pi_c - \pi_d$ for the defier share $\pi_d$ and effect $\tau_d$, and the ratio can lie outside the range of the individual effects. Without exclusion the always takers add $\E\{Y(1, 1) - Y(0, 1) \mid \text{always taker}\}$ times their share to the numerator, and the never takers a similar term.

\paragraph{Exercise~\ref{ex:ch10:4}.} Let $S_i \in \{c, a, n\}$ be the latent stratum (complier, always taker, never taker), with shares $(\pi_c, \pi_a, \pi_n) \sim \mathrm{Dirichlet}(\alpha_c, \alpha_a, \alpha_n)$ independent of $Z$ by random assignment. Exclusion makes the outcome law of always takers and never takers free of $Z$: $Y_i \mid S_i = a \sim N(\mu_a, \sigma_a^2)$ and $Y_i \mid S_i = n \sim N(\mu_n, \sigma_n^2)$, while $Y_i \mid S_i = c, Z_i = z \sim N(\mu_{c,z}, \sigma_{c,z}^2)$, so the local effect is $\mu_{c,1} - \mu_{c,0}$. Each of the four outcome groups $g \in \{a, n, (c,0), (c,1)\}$ has the conjugate prior $\sigma_g^2 \sim \nu_0 s_0^2/\chi^2_{\nu_0}$, $\mu_g \mid \sigma_g^2 \sim N(m_0, \sigma_g^2/n_0)$, and $n_g$, $\bar y_g$ and $\mathrm{SS}_g$ denote the number, mean and within sum of squares of the outcomes of the units currently in group $g$. One sweep has three steps. First, impute the strata. A unit with $(Z_i, T_i) = (1, 0)$ is a never taker and one with $(0, 1)$ an always taker, because there are no defiers. A unit with $T_i = Z_i$ is a complier, or else an always taker if $Z_i = 1$ and a never taker if $Z_i = 0$; call that alternative stratum $s_i$. Then, independently over units,
\[
\Prob(S_i = c \mid \cdot) = \frac{\pi_c\,N(y_i; \mu_{c,z_i}, \sigma_{c,z_i}^2)}{\pi_c\,N(y_i; \mu_{c,z_i}, \sigma_{c,z_i}^2) + \pi_{s_i}\,N(y_i; \mu_{s_i}, \sigma_{s_i}^2)}.
\]
Second, draw $(\pi_c, \pi_a, \pi_n) \sim \mathrm{Dirichlet}(\alpha_c + n_{c,0} + n_{c,1},\, \alpha_a + n_a,\, \alpha_n + n_n)$. Third, for each group draw
\begin{align*}
\sigma_g^2 &= \frac{\nu_0 s_0^2 + \mathrm{SS}_g + n_0n_g(\bar y_g - m_0)^2/(n_0 + n_g)}{\chi^2_{\nu_0 + n_g}}, \\
\mu_g \mid \sigma_g^2 &\sim N\Bigl(\frac{n_0m_0 + n_g\bar y_g}{n_0 + n_g},\; \frac{\sigma_g^2}{n_0 + n_g}\Bigr),
\end{align*}
with always takers of both arms in one group, never takers in another, and compliers split by $Z$. The draws of $\mu_{c,1} - \mu_{c,0}$ are a posterior sample of the local effect. Given the strata the likelihood factorizes over the four groups, which is why the prior is conjugate group by group, and given the parameters the strata are independent across units.

\paragraph{Exercise~\ref{ex:ch10:5}.} Let $\mathcal C$ have eigenvalues $s_k > 0$ and orthonormal eigenfunctions $\phi_k$, and write $h_k = \langle h, \phi_k\rangle$. For the Gaussian prior with covariance $\mathcal C^{2r}$ of finite trace, the Cameron Martin theorem gives
\[
\|h\|_{\mathrm{CM}}^2 = \sum_k \frac{h_k^2}{s_k^{2r}}.
\]
If $f = \mathcal C^r g$, then $f_k = s_k^r g_k$ and $\|f\|_{\mathrm{CM}}^2 = \sum_k g_k^2 = \|g\|^2$. Conversely, finite Cameron Martin norm gives $g_k = f_k/s_k^r$ with $\|g\| < \infty$ and $f = \mathcal C^r g$. Thus the source condition and Cameron Martin membership are equivalent, with equal norms.

\paragraph{Exercise~\ref{ex:ch10:6}.} In the eigenbasis of Exercise~\ref{ex:ch10:5}, write $f_k = s_k^r g_k$ with $\sum_k g_k^2 \le 1$. Ridge multiplies coordinate $k$ by $s_k/(s_k + \alpha)$, so its worst squared bias is
\[
\sup_k\Bigl\{\frac{\alpha s_k^r}{s_k + \alpha}\Bigr\}^2.
\]
For $r \le 1$, the weighted arithmetic and geometric means give $\alpha s^r/(s + \alpha) \le \alpha^r$, with a matching lower bound at an eigenvalue within a constant factor of $\alpha$, as exists under polynomial decay. For $r > 1$ and $\alpha \le s_1$,
\[
\frac{\alpha s_1^{r-1}}2 \le \frac{\alpha s_1^r}{s_1 + \alpha},
\qquad
\frac{\alpha s_k^r}{s_k + \alpha} \le \alpha s_1^{r-1}.
\]
The worst squared bias is therefore of order $\alpha^{2\min(r,1)}$. Since the variance does not depend on $r$, balancing it against this bias gives the same rate for every $r \ge 1$: ridge removes bias only to first order in $\alpha$.

\paragraph{Exercise~\ref{ex:iv:wright}.} Equating demand and supply gives
\[
p = \frac{a - c + u_d - u_s - ew}{b + d}, \qquad q = a - bp + u_d.
\]
Let $w$, $u_d$ and $u_s$ be independent with variances $\sigma_w^2$, $\sigma_d^2$ and $\sigma_s^2$. Then $\mathrm{Var}(p) = (\sigma_d^2 + \sigma_s^2 + e^2\sigma_w^2)/(b + d)^2$ and $\mathrm{Cov}(p, u_d) = \sigma_d^2/(b + d)$, so the least squares slope of $q$ on $p$ converges to
\[
\frac{\mathrm{Cov}(p, q)}{\mathrm{Var}(p)} = -b + \frac{\mathrm{Cov}(p, u_d)}{\mathrm{Var}(p)} = -b + (b + d)\alpha = -(1 - \alpha)\,b + \alpha\,d,
\]
where $\alpha = \sigma_d^2/(\sigma_d^2 + \sigma_s^2 + e^2\sigma_w^2) \in (0, 1)$ is the weight on the supply slope, the share of the price variance that comes from demand shocks. When they dominate, price moves along the supply curve and the scatter traces supply; when supply shocks and the shifter dominate, it traces demand. The shifter identifies demand because $\mathrm{Cov}(w, q)/\mathrm{Cov}(w, p) = -b$. In Figure~\ref{fig:wright}, $b = 1$, $d = 0.8$, $e = 1.5$ and all variances are one, so $\alpha = 1/4.25 = 0.235$ and the limit is $-0.765 + 0.188 = -0.577$. The sample slope of $-0.503$ is 1.3 standard errors above it, the standard error of the least squares slope in 200 markets being 0.055 in this design, so the figure is an ordinary draw. The instrumental variable slope of $-1.01$ is close to $-1$.

\paragraph{Exercise~\ref{ex:iv:card}.} Both parents' education is recorded for 2,220 of the 3,010 men, so 790 (26 percent) are lost, and they are not a random quarter: 44 percent of them are Black against 16 percent of those kept, and they have 12.2 years of schooling on average against 13.6. To separate the two changes, fit the model three times: on the book's sample, on the 2,220 men with the book's controls, and on the 2,220 men with the father's and mother's education added to the controls. The posteriors are those of the chapter, the reduced form posterior of $\pi_Y/\pi_T$ (20,000 draws) and the exact structural posterior.

On the book's sample two stage least squares gives 0.132 (standard error 0.055) with a first stage F of 13.3; the reduced form posterior has median 0.131 and 90\% interval 0.043 to 0.249, and the structural posterior median 0.135 and interval 0.039 to 0.270. On the 2,220 men with the same controls the first stage coefficient falls from 0.32 to 0.27 and its standard error rises from 0.088 to 0.104, so F falls to 6.6. Two stage least squares gives 0.080 (0.074), the reduced form posterior 0.079 (90\% interval $-0.076$ to 0.241) and the structural posterior 0.080 ($-0.139$ to 0.303). Adding the parents' education on the same 2,220 changes little: F 7.1, two stage least squares 0.084 (0.076), reduced form posterior 0.083 ($-0.069$ to 0.248), structural posterior 0.085 ($-0.122$ to 0.304). The probability of a positive return falls from about 0.99 under both posteriors to 0.84 under the reduced form posterior and 0.80 under the structural one on the 2,220 men with the same controls, and is 0.85 and 0.81 with the parents' education added.

So the change comes from the sample, not from the controls. Parents' education correlates with schooling (0.47 for the father, 0.44 for the mother) and college proximity (0.14 and 0.08), yet adding it moves two stage least squares by only 0.004, a small fraction of its standard error, so measured family background does not account for the proximity effect. The fall from 0.132 to about 0.08 comes with dropping 790 men who differ from the rest, and it is less than one subsample standard error, so the data cannot say whether the return differs between the two groups. Losing a quarter of the sample and a sixth of the first stage leaves a weak instrument (F near 7). Two stage least squares lands near least squares (0.080 against 0.077), where weak instrument bias pulls it, with a standard error a third larger and an unreliable normal interval. Both posteriors widen, and unequally: the structural 90\% interval is 13 percent wider than the reduced form interval on the book's sample (0.231 against 0.204) and 34 to 39 percent wider on the subsample (0.43 and 0.44 against 0.32), because the structural posterior has the longer tails and weak identification exposes them (Section~\ref{iv:sec:weak}).

\paragraph{Exercise~\ref{ex:iv:qob}.} Let $e = Y - \tau T$ be the structural error and $P_Z$ the projection on the column space of the $K$ instruments. Then
\[
\hat\tau_{\mathrm{2SLS}} - \tau = \frac{T^\top P_Z e}{T^\top P_Z T}, \qquad \hat\tau_{\mathrm{LS}} - \tau = \frac{T^\top e}{T^\top T}.
\]
Random numbers are independent of $(T, e)$. For independent normal instruments the column space is a uniformly random $K$ dimensional subspace, so for fixed $T$ and $e$ we have $\E P_Z = (K/n)I$, and $T^\top P_Z e$ and $T^\top P_Z T$ are centred at $K/n$ times $T^\top e$ and $T^\top T$. The ratio $T^\top P_Z T/\|T\|^2$ has a Beta$(K/2, (n - K)/2)$ law with coefficient of variation about $\sqrt{2/K}$, and the numerator fluctuates at the same order. So for large $K$ the two stage error is the least squares error times $1 + O(K^{-1/2})$: the fitted treatment $P_ZT$ is a random slice of $T$ itself, endogenous variation included. In the design of Figure~\ref{fig:weakiv} ($n = 2000$, errors correlated 0.6, $\tau = 0.1$, least squares near 0.70) with first stage coefficients set to zero, 300 replications give interquartile ranges for two stage least squares of $(-0.43, 1.4)$ for $K = 1$, $(0.57, 0.83)$ for $K = 20$ and $(0.67, 0.75)$ for $K = 180$, with medians 0.59, 0.70 and 0.71, against a truth of 0.1. For $K = 1$ the law is close to 0.1 plus a Cauchy with location 0.6 and scale 0.8 (quartiles $-0.10$ and 1.50, median 0.70), so the simulated figures of that row carry Monte Carlo error of about 0.1. Bound, Jaeger and Baker (1995) made this replacement with quarters of birth (Section~\ref{iv:sec:unlikely}).

With a single random instrument the reduced form posterior for $\tau = \pi_Y/\pi_T$ is the law of a ratio of two normals centred at noise, since $\hat\pi_T$ is a draw from $N(0, s_T^2)$. It puts mass near $\pi_T = 0$, has Cauchy like tails and no mean (Exercise~\ref{ex:ch10:2}), and so reports that the data say almost nothing about $\tau$. In the 200 data sets of the reduced form study of Section~\ref{iv:sec:weak}, with the first stage set to zero, the median 90\% interval has length 9.7, against 7.3 for the weak instrument of that section, and the median posterior median is 0.69, near least squares and far from 0.1.

\paragraph{Exercise~\ref{ex:iv:samesex}.} Let $Z$ indicate two children of the same sex and $T$ a third child. Compliers continue after two of one sex and stop after a boy and a girl, plausibly because they want both sexes. Always takers continue either way, never takers stop either way, and monotonicity excludes the reverse response. With stratum shares $\pi_g$ and effects $\tau_g$, $g \in \{c, a, n\}$, the Wald ratio identifies $\tau_c$, while the average effect over all mothers of two is $\tau = \pi_c\tau_c + \pi_a\tau_a + \pi_n\tau_n$, so
\[
\tau - \tau_c = \pi_a(\tau_a - \tau_c) + \pi_n(\tau_n - \tau_c).
\]
Always takers are observed only treated and never takers only untreated, so neither $\tau_a$ nor $\tau_n$ is identified.

Even if compliers share everyone else's average effect given $x$, their covariate law $Q_c(dx) \propto \pi_c(x)P(dx)$ differs from that of all mothers of two, $P(dx)$, where $\pi_c(x) = \Prob(\text{complier} \mid X = x)$. The accounting identity of Section~\ref{iv:sec:accounting} then gives
\[
\tau_c - \tau = \frac{\mathrm{Cov}\{\tau(X), \pi_c(X)\}}{\E\,\pi_c(X)},
\]
which is not zero when the effect varies with covariates that predict compliance.

\paragraph{Exercise~\ref{ex:iv:iwscale}.} With $\pi = 0$ and no controls, $T = e_1$ and $Y = \tau e_1 + e_2$, so $(T, Y)$ is bivariate normal with mean zero and covariance
\[
\Omega = \begin{pmatrix} \sigma_{11} & \tau\sigma_{11} + \sigma_{12} \\ \tau\sigma_{11} + \sigma_{12} & \tau^2\sigma_{11} + 2\tau\sigma_{12} + \sigma_{22} \end{pmatrix},
\]
and the likelihood is a function of $\Omega$ alone. Given $\tau$ and $\Omega$ these equations have one solution, $\sigma_{11} = \Omega_{11}$, $\sigma_{12} = \Omega_{12} - \tau\Omega_{11}$ and $\sigma_{22} = \Omega_{22} - 2\tau\Omega_{12} + \tau^2\Omega_{11}$, and along the curve it traces
\[
|\Sigma| = \Omega_{11}(\Omega_{22} - 2\tau\Omega_{12} + \tau^2\Omega_{11}) - (\Omega_{12} - \tau\Omega_{11})^2 = \Omega_{11}\Omega_{22} - \Omega_{12}^2 = |\Omega|.
\]
The inverse Wishart density with 3 degrees of freedom and scale $S$ is proportional to $|\Sigma|^{-3}\exp\{-\mathrm{tr}(S\Sigma^{-1})/2\}$, and the first factor is constant along the curve. The diagonal of $\Sigma^{-1}$ is $(\sigma_{22}, \sigma_{11})/|\Sigma|$, so with $S = c\,\mathrm{diag}(\Omega_{11}, \Omega_{22})$
\[
\mathrm{tr}(S\Sigma^{-1}) = \frac{c\,(\Omega_{11}\sigma_{22} + \Omega_{22}\sigma_{11})}{|\Omega|} = \frac{c\,\Omega_{11}^2(\tau - b)^2 + c\,(2\Omega_{11}\Omega_{22} - \Omega_{12}^2)}{|\Omega|},
\]
with $b = \Omega_{12}/\Omega_{11}$, which gives the factor $\exp\{-c\,\Omega_{11}^2(\tau - b)^2/(2|\Omega|)\}$. At fixed $\tau$ the change of variables from $\Sigma$ to $\Omega$ is triangular with unit diagonal, so its Jacobian is one, and as $n$ grows the likelihood pins $\Omega$ at the sample covariance matrix. The posterior of $\tau$ is then close to normal with mean $b$ and standard deviation $|\Omega|^{1/2}/(\Omega_{11}\sqrt{c})$; the $N(0, 10^4)$ prior on $\tau$ changes nothing at this scale. In the design of Figure~\ref{fig:weakiv} with one instrument, $\Omega_{11} = 1.0009$, $\Omega_{12} = 0.7001$ and $\Omega_{22} = 1.1300$, so $|\Omega| = 0.641$, the standard deviation is $0.80/\sqrt{c}$ and the 90\% interval has length $2 \times 1.645 \times 0.80/\sqrt{c} = 2.63/\sqrt{c}$: 26.3, 8.3, 2.6 and 0.83 for $c = 0.01$, 0.1, 1 and 10. Section~\ref{iv:sec:weak} reports \ivscalesmall, \ivscaledeflen, \ivscaleten\ and \ivscalehund. The reported lengths are shorter because the first stage coefficient is 0.03, not zero: the instrument, though weak, curves the likelihood along the ridge, and its information counts for more the flatter the prior, so the gap is largest at $c = 0.01$ and small at $c = 10$, where the prior dominates.

\paragraph{Exercise~\ref{ex:iv:flatK}.} Write $A = (Z^\top Z)^{1/2}/\sigma_T$ and $\theta = A\pi$. A flat prior on $\pi$ is flat on $\theta$, because the map is linear with a constant Jacobian, and $\mu^2 = \theta^\top\theta = r^2$. In polar coordinates in $\mathbb{R}^K$ the volume element is $r^{K-1}\,dr$ times the element of the unit sphere, so $r$ has prior density proportional to $r^{K-1}$, and with $s = r^2$ and $dr = ds/(2\sqrt{s})$,
\[
p(s) \propto s^{(K-1)/2}\, s^{-1/2} = s^{K/2 - 1}.
\]
For $K = 1$ this is $s^{-1/2}$, which decreases: the prior leans toward a weak first stage, the posterior is free to put $\pi$ near zero, where $\tau$ lies on the ridge of Section~\ref{iv:sec:many}, and the interval is wide. For $K = 20$ it is $s^9$. On any interval $[0, m]$ that density puts a share $(s/m)^{10}$ of its mass below $s$, so a concentration parameter below half the upper limit gets about a thousandth of the mass: the prior asserts a strong first stage. The posterior combines that assertion with a likelihood that is nearly flat in the direction of weakness, takes the overfitted first stage at face value and centres near two stage least squares, which is what Section~\ref{iv:sec:weak} reports: posterior medians at \ivmanybmed\ against \ivmanytsls\ for two stage least squares, and coverage of \ivmanybcov\ with twenty instruments against \ivweakbcov\ with one. A hierarchical prior $\pi_k \sim N(0, v)$, with a prior on $v$ that allows small values, lets $\mu^2$ be small, so the posterior can put the first stage near zero, as the hierarchical and horseshoe rows of Table~\ref{tab:ivmany} show.

\section*{Chapter~\ref{ch:tianpearl}: Bounds and Sensitivity}

\paragraph{Exercise~\ref{ex:ch11:1}.} Here $X$ is independent of the pair $(Y(0), Y(1))$, so the pair has the same law given either value of $X$. Write $p_{ab} = \Prob\{Y(1) = a, Y(0) = b\}$ as in the proof of the experimental bounds on $\PNS$, with $p_1 = \Prob(y_x) = p_{10} + p_{11}$ and $p_0 = \Prob(y_{x'}) = p_{01} + p_{11}$. Among the exposed who fall ill, $Y = Y(1) = 1$, and given $X = 1$ the pair $(Y(1), Y(0))$ has the law $p_{ab}$, so
\[
\PN = \Prob\{Y(0) = 0 \mid X = 1, Y(1) = 1\} = \frac{p_{10}}{p_{10} + p_{11}} = \frac{\PNS}{p_1}.
\]
Among the unexposed who stay well, $Y = Y(0) = 0$, and in the same way
\[
\PS = \Prob\{Y(1) = 1 \mid X = 0, Y(0) = 0\} = \frac{p_{10}}{p_{10} + p_{00}} = \frac{\PNS}{1 - p_0}.
\]
Dividing $\max\{0, p_1 - p_0\} \leq \PNS \leq \min\{p_1, 1 - p_0\}$ by $p_1$ and by $1 - p_0$ gives
\begin{gather*}
\max\Bigl\{0,\, 1 - \frac{p_0}{p_1}\Bigr\} \leq \PN \leq \min\Bigl\{1,\, \frac{1 - p_0}{p_1}\Bigr\}, \\
\max\Bigl\{0,\, \frac{p_1 - p_0}{1 - p_0}\Bigr\} \leq \PS \leq \min\Bigl\{1,\, \frac{p_1}{1 - p_0}\Bigr\}.
\end{gather*}


\paragraph{Exercise~\ref{ex:ch11:2}.} Under exogeneity, $\Prob(y) - \Prob(y_{x'}) = \Prob(x)\{\Prob(y \mid x) - \Prob(y \mid x')\}$ and $\Prob(x, y) = \Prob(x)\Prob(y \mid x)$. The ratio is $1 - \Prob(y \mid x')/\Prob(y \mid x) = 1 - 1/\mathrm{RR}$, so the lower bound is $\max\{0, 1 - 1/\mathrm{RR}\}$, and it is $1 - 1/\mathrm{RR}$ when $\mathrm{RR} \geq 1$.

\paragraph{Exercise~\ref{ex:ch11:3}.} Let $q_{x,a,b} = \Prob\{X = x, Y(1) = a, Y(0) = b\}$ be the eight cells, as in Appendix~\ref{app:graphs}. The program minimizes and maximizes $\PNS = q_{0,1,0} + q_{1,1,0}$ over $q \geq 0$ with $\sum q = 1$, the experimental constraints
\[
\sum_{x,b} q_{x,1,b} = \Prob(y_x), \qquad \sum_{x,a} q_{x,a,1} = \Prob(y_{x'}),
\]
and the observational constraints, which hold because $Y = Y(X)$:
\begin{gather*}
\sum_b q_{1,1,b} = \Prob(x, y), \qquad \sum_b q_{1,0,b} = \Prob(x, y'), \\
\sum_a q_{0,a,1} = \Prob(x', y), \qquad \sum_a q_{0,a,0} = \Prob(x', y').
\end{gather*}
Take $\Prob(y_x) = 0.30$ and $\Prob(y_{x'}) = 0.10$, the population values behind Figure~\ref{fig:pns}.  With 30\% exposed and outcome rates 0.90 and 0.05, so that $\Prob(x, y) = 0.27$, $\Prob(x, y') = 0.03$, $\Prob(x', y) = 0.035$ and $\Prob(x', y') = 0.665$, it returns $[0.205, 0.265]$. The closed forms agree: the lower end is $\Prob(y) - \Prob(y_{x'}) = 0.305 - 0.10$ and the upper end is $\Prob(y_x) - \Prob(y_{x'}) + \Prob(x, y') + \Prob(x', y) = 0.20 + 0.03 + 0.035$. The data narrow the interval here because they are nearly deterministic: the discordant cells, exposed without the outcome and unexposed with it, hold only 6.5\% of the population, and that caps $\PNS$ at the risk difference plus 0.065.

\paragraph{Exercise~\ref{ex:ch11:4}.} Let $p_t = \Prob(Y(t) = 1)$, let $(Z_0, Z_1)$ be standard bivariate normal with correlation $\rho$, and set $Y(t) = 1$ when $Z_t > \Phi^{-1}(1 - p_t)$. The margins are right and the latent pair has the Gaussian copula $C_\rho$. Replacing $Z_t$ by $-Z_t$ shows $\Prob\{Y(1) = 1, Y(0) = 1\} = \Phi_2\{\Phi^{-1}(p_1), \Phi^{-1}(p_0); \rho\}$, and $\PNS = p_1 - \Prob\{Y(1) = 1, Y(0) = 1\}$, so
\[
\PNS(\rho) = p_1 - \Phi_2\{\Phi^{-1}(p_1), \Phi^{-1}(p_0); \rho\}.
\]
The derivative of $\Phi_2$ in $\rho$ is the bivariate normal density, which is positive, so $\PNS$ falls strictly as $\rho$ rises. At $\rho = 1$, $\Phi_2 = \min(p_1, p_0)$ and $\PNS = \max\{0, p_1 - p_0\}$, the lower bound. At $\rho = -1$, $\Phi_2 = \max(0, p_1 + p_0 - 1)$ and $\PNS = \min\{p_1, 1 - p_0\}$, the upper bound. At $\rho = 0$, $\PNS = p_1(1 - p_0)$. The family therefore sweeps the whole interval, and the position $u$ within the bounds is a decreasing function of $\rho$, so a prior on $\rho$ induces a prior on $u$. For $p_1 = 0.30$ and $p_0 = 0.10$, with bounds 0.20 and 0.30, the values $\rho = 0.5, 0, -0.5$ give $\PNS = 0.235, 0.270, 0.294$ and $u = 0.35, 0.70, 0.94$.

\paragraph{Exercise~\ref{ex:ch11:5}.} With $\mathrm{RR} = 10$ the lower bound is $0.9$. To reach a probability near one, one needs either that the risk among the unexposed is produced by causes absent in this plaintiff, which is a statement about her covariates, or a model with dose response showing that her exposure was far above average, which raises her individual risk ratio.

\paragraph{Exercise~\ref{ex:ch11:6}.} Within a stratum let $d = \Prob(y_x) - \Prob(y_{x'}) = p_{10} - p_{01}$, which the experiment identifies, and write $\PNS_{\mathrm{lo}}$ and $\PNS_{\mathrm{hi}}$ for the bounds on $\PNS = p_{10}$. The mean benefit of targeting a customer is
\[
B = g p_{10} - L p_{01} - c = (g - L)\,\PNS + L d - c.
\]
If $g = L$, then $B = g d - c$ and the rule is to target when $d > c/g$: the coupling drops out. If $g \neq L$, $B$ is linear in $\PNS$, so under a uniform prior on $u$, independent of the margins, the expected benefit is $B$ at the mean of $\PNS$, the midpoint of the bounds:
\[
\E(B \mid \text{data}) = (g - L)\,\frac{\PNS_{\mathrm{lo}} + \PNS_{\mathrm{hi}}}{2} + L d - c,
\]
with the midpoint and $d$ replaced by their posterior means when the margins are uncertain. The Bayes rule targets when this is positive. The minimax rule targets when $B$ is positive at the worst coupling, $\PNS_{\mathrm{lo}}$ if $g > L$ and $\PNS_{\mathrm{hi}}$ if $g < L$. The two benefits differ by $|g - L|(\PNS_{\mathrm{hi}} - \PNS_{\mathrm{lo}})/2$, so the Bayes rule targets every stratum the minimax rule targets and also those with $0 < \E(B \mid \text{data}) \leq |g - L|(\PNS_{\mathrm{hi}} - \PNS_{\mathrm{lo}})/2$, where the break even point lies inside the bounds. With $\Prob(y_x) = 0.30$ and $\Prob(y_{x'}) = 0.10$, so that $d = 0.20$ and the bounds are 0.20 and 0.30, and with $g = 1$, $L = 2$ and $c = 0.12$, we get $B = 0.28 - \PNS$. The worst case $\PNS = 0.30$ gives $-0.02$ and the minimax rule does not target; the midpoint gives $0.03$ and the Bayes rule does.

\paragraph{Exercise~\ref{ex:ch12:1}.} For fixed $f_1 > f_0$, Cornfield's formula gives
\[
\mathrm{RR}(R) = \frac{1 + f_1(R - 1)}{1 + f_0(R - 1)}, \qquad
\frac{d\,\mathrm{RR}}{dR} = \frac{f_1 - f_0}{\{1 + f_0(R - 1)\}^2} > 0.
\]
The denominator is positive for $R > 0$. For $f_0 > 0$, dividing numerator and denominator by $R$ gives $\mathrm{RR}(R) \to f_1/f_0$. The limit is never reached: $\mathrm{RR}(R) < f_1/f_0$ is equivalent to $f_0 < f_1$.

\paragraph{Exercise~\ref{ex:ch12:2}.} Write $a = \mathrm{RR}_{EU}$, $b = \mathrm{RR}_{UD}$ and $B(a, b) = ab/(a + b - 1)$ for $a, b \ge 1$. Since $\partial B/\partial a = b(b - 1)/(a + b - 1)^2 \ge 0$, and likewise for $b$, risk ratios bounded by $m$ give $B(a,b) \le m^2/(2m - 1)$, attained at $a = b = m$. This bound increases from one for $m \ge 1$. The smallest $m$ that explains away $\mathrm{RR} > 1$ therefore solves
\[
m^2 - 2\,\mathrm{RR}\,m + \mathrm{RR} = 0, \qquad
m = \mathrm{RR} + \sqrt{\mathrm{RR}(\mathrm{RR} - 1)}.
\]
The other root is below one.

\paragraph{Exercise~\ref{ex:ch12:3}.} The E values are $2 + \sqrt{2} \approx 3.41$ for $\mathrm{RR} = 2$, $9 + \sqrt{72} \approx 17.49$ for $\mathrm{RR} = 9$, and $20 + \sqrt{380} \approx 39.49$ for $\mathrm{RR} = 20$. A risk ratio of 2 is explained away by an unmeasured factor associated with both smoking and lung cancer by risk ratios of about 3.4, which social and occupational exposures could plausibly reach, so an association of that size would not by itself have settled the question. The ratio of about 9 for smokers against nonsmokers in the evidence Cornfield reviewed needs a confounder with risk ratios near 17.5 with both smoking and the disease, far stronger than any factor proposed at the time; for Fisher's constitutional genotype no association of that strength was shown. A ratio of 20 needs risk ratios near 40. 

\paragraph{Exercise~\ref{ex:ch12:4}.} By Proposition~\ref{prop:partial} the limit is the law of $T(\xi)$ with $T = h(\phi_0, \cdot)$ and $\xi$ drawn from its prior. A law is a point mass at $c$ exactly when the variable equals $c$ with probability one, which is the first claim. If $T = c$ on the support, which has prior mass one, then $T(\xi) = c$ almost surely. Conversely, suppose $T = c$ almost surely and $T(\xi^\ast) \neq c$ at a point $\xi^\ast$ of the support. Continuity gives a neighbourhood of $\xi^\ast$ on which $T \neq c$, and it has positive prior mass because $\xi^\ast$ is in the support, a contradiction. Without continuity the second equivalence fails: with $\xi \sim N(0, 1)$ and $T(\xi) = \mathbb{1}\{\xi = 0\}$ the limit is a point mass at $0$ although $T$ is not constant on the support.

\paragraph{Exercise~\ref{ex:ch12:5}.} By Proposition~\ref{prop:partial}, with $\tau = \phi - ab_U$, $\phi$ the identified parameter with limit $\tau_0$ and $\xi = (a, b_U)$, the limiting posterior of $\tau$ is the law of $\tau_0 - a b_U$ with $(a, b_U)$ drawn from the prior. Write $a = s_a Z_1$ and $b_U = s_b Z_2$ with $Z_1, Z_2$ independent standard normal, and $t = \tau_0/(s_a s_b)$. Then $\tau > 0$ exactly when $Z_1 Z_2 < t$, and conditioning on $Z_1 = z$,
\begin{align*}
\Prob(\tau > 0) = \Prob(Z_1 Z_2 < t) &= \int_{-\infty}^{\infty} \varphi(z)\,\Phi\bigl(t/|z|\bigr)\,dz \\
&= 2\int_0^\infty \varphi(z)\,\Phi(t/z)\,dz,
\end{align*}
where $\varphi$ and $\Phi$ are the standard normal density and distribution function. For $z > 0$ the event is $Z_2 < t/z$, and for $z < 0$ it is $Z_2 > t/z$, which has probability $\Phi(-t/z) = \Phi(t/|z|)$. The scales and $\tau_0$ enter only through $t$, so the probability is one increasing function of $t$: it equals $1/2$ at $t = 0$ and tends to one as $t \to \infty$. For $\tau_0 = 0.30$ and $s_a = s_b = 1$ we have $t = 0.30$ and the integral is 0.723 (0.72 in Figure~\ref{fig:sens}, which also carries the standard error 0.05). The pair $(s_a, s_b) = (0.5, 2)$ gives the same value, as does any pair with the same product.

\paragraph{Exercise~\ref{ex:ch12:6}.} The book has no size coefficient from the real Lanarkshire trial, so we use the simulated trial of Section~\ref{sec:sensitivity} at correlation 0.8 between recorded size and the hidden score $U$; the recipe is the same with real coefficients. Regress the gain on treatment and size to get $\hat\tau$, and let $U$ have standard deviation one. Let $a$ be the coefficient of $T$ in the regression of $U$ on $T$ and size, and $b_U$ the coefficient of $U$ when $U$ joins the regression. The omitted variable identity for least squares gives $\hat\tau - \tau = a\,b_U$ up to estimation error, and the product does not change if $U$ is rescaled. Neither factor is identified, so each gets a prior scaled by what size does in the data: $a_S = -0.708$, the coefficient of $T$ in a regression of size on $T$, and $b_S = -0.527$, the coefficient of size in the regression of the gain on $T$ and size. The story fixes the signs, since teachers favoured children who looked poorly and those children grow more. Hence
\[
a \sim N(k a_S, k^2a_S^2), \qquad b_U \sim N(k b_S, k^2b_S^2), \qquad \tau = \hat\tau - a\,b_U + 0.025\,z,
\]
with $a$, $b_U$ and $z \sim N(0, 1)$ independent and $k$ the strength of the hidden score relative to size. Independence gives $\E(a\,b_U) = k^2a_Sb_S = 0.373k^2$ and $\mathrm{Var}(a\,b_U) = (2k^2a_S^2)(2k^2b_S^2) - k^4a_S^2b_S^2 = 3k^4a_S^2b_S^2$, so the prior mean of the bias is $0.373k^2$ with standard deviation $0.647k^2$.

These formulas read the family before any sampling. With $\hat\tau = 0.351$, the prior mean of $\tau$ is zero at $k = (0.351/0.373)^{1/2} = 0.97$, so on average a hidden score as strong as size removes the whole estimate, and it reaches the true 0.2 at $k = 0.64$. The means $0.351 - 0.373k^2$ at $k = 1/4, 1/2, 1$ are 0.33, 0.26 and $-0.02$, the column printed in Section~\ref{sec:sensitivity}.

The choice of $k$ is the judgement. Only the part of $U$ that size does not explain can bias $\hat\tau$. If $U$ has correlation $r$ with size, that part has standard deviation $(1 - r^2)^{1/2}$, which suggests $k \approx (1 - r^2)^{1/2}$ as a rough guide. It is a heuristic, because it assumes the hidden score acts per standard deviation as size does. At the true $r = 0.8$ it gives 0.6, close to the hidden score's own ratios $0.39/0.71 = 0.55$ and $0.30/0.53 = 0.57$. An analyst does not know $r$, so she can draw $r$ uniformly on $(0.5, 0.95)$, which puts $k$ between 0.31 and 0.87, and then $a$, $b_U$ and $z$ as above. With 200,000 draws (seeds 1933 and 1934) the posterior of $\tau$ has mean 0.18, 95\% interval $-0.68$ to 0.66, and $\Prob(\tau > 0) = 0.80$. The interval contains the truth, which the classical interval $[0.30, 0.40]$ misses. The price is width, and that width is honest: the data say nothing about $a$ and $b_U$, so the prior on them decides the answer. Two cautions apply. The benchmark $a_S$ is not conditioned on any other covariate, while $a$ is conditioned on size, so a value of $k$ below one is the natural reading. And size was recorded because someone thought of it, which is weak evidence about what was not recorded, so the report should show the whole family of $k$ and not a single value.

\section*{Chapter~\ref{ch:smoking}: Smoking and Lung Cancer}

\paragraph{Exercise~\ref{ex:ch13:1}.} Write $G$ for the graph of Figure~\ref{fig:dagfd}, with arrows $S \to M$, $M \to C$, $U \to S$ and $U \to C$. Summing over $M$,
\[
\Prob(c \mid do(s)) = \sum_m \Prob(m \mid do(s))\,\Prob(c \mid do(s), m).
\]
For the first factor, $M \indep S$ in $G_{\underline{S}}$, because the only path between them, $S \leftarrow U \to C \leftarrow M$, is blocked at the collider $C$. Rule 2 then gives $\Prob(m \mid do(s)) = \Prob(m \mid s)$. For the second factor, $C \indep M \mid S$ in $G_{\overline{S}\underline{M}}$, which keeps only $S \to M$ and $U \to C$, so rule 2 gives $\Prob(c \mid do(s), m) = \Prob(c \mid do(s), do(m))$. Next, $C \indep S \mid M$ in $G_{\overline{M}\,\overline{S}}$, which keeps only $M \to C$ and $U \to C$; here $S(W) = S$ because $W$ is empty, so rule 3 removes the action: $\Prob(c \mid do(s), do(m)) = \Prob(c \mid do(m))$. Summing over $S$,
\[
\Prob(c \mid do(m)) = \sum_{s'} \Prob(s' \mid do(m))\,\Prob(c \mid do(m), s').
\]
In $G_{\overline{M}}$ the path $S \leftarrow U \to C \leftarrow M$ is blocked at $C$, so $S \indep M$ and rule 3 gives $\Prob(s' \mid do(m)) = \Prob(s')$. In $G_{\underline{M}}$ the only path from $M$ to $C$ is $M \leftarrow S \leftarrow U \to C$, which conditioning on $S$ blocks, so rule 2 gives $\Prob(c \mid do(m), s') = \Prob(c \mid m, s')$. Substituting the three results into the first display gives \eqref{eq:frontdoor}.

\paragraph{Exercise~\ref{ex:ch13:2}.} Change the tar equation to
\[
\Prob(M = 1 \mid S, U) = \expit(-2 + 4S + \gamma U)
\]
and keep the rest of the design. The derivation of Exercise~\ref{ex:ch13:1} then breaks twice: $S \leftarrow U \to M$ is open, so $\Prob(m \mid do(s)) \neq \Prob(m \mid s)$, and adjusting for $S$ no longer blocks $M \leftarrow U \to C$ in the effect of tar on cancer. The two quantities to compare are exact sums over $U$, $S$, $M$ and $C$, so no sampling is needed: the true risk difference, and the value of \eqref{eq:frontdoor} computed from the joint law of $(S, M, C)$, which is what the plug-in estimator converges to. At $\gamma = 0$ both are 0.145. At $\gamma = 0.5, 1, 2$ the truth is 0.143, 0.137, 0.118 and the front door value is 0.161, 0.176, 0.198, a bias of 0.018, 0.039, 0.080; at $\gamma = -1$ the bias is $-0.027$. The bias has the sign of $\gamma$ and grows with $|\gamma|$, and a larger sample does not remove it. For $\gamma \geq 0$ the naive contrast stays near 0.23, so the front door value remains the closer of the two but is no longer consistent.

\paragraph{Exercise~\ref{ex:ch13:3}.} Let $\mu(s, x)$ be the fitted probability of dying of lung cancer in the follow-up window for smoking status $s$ and covariates $x$, and let $Q_{SX}$ be the law of $(S, X)$ in the population of interest, whose covariate margin is $Q$, so that it records who smokes as well as age and sex. With $\mu$ read causally, the risk if no one smoked is $\E_{Q_{SX}}\mu(0, X)$ and the observed risk is $\E_{Q_{SX}}\mu(S, X)$, so
\begin{align*}
\mathrm{PAF}(\mu, Q_{SX}) &= \frac{\E_{Q_{SX}}\mu(S, X) - \E_{Q_{SX}}\mu(0, X)}{\E_{Q_{SX}}\mu(S, X)} \\
&= \frac{\E_{Q_{SX}}\bigl[S\{\mu(1, X) - \mu(0, X)\}\bigr]}{\E_{Q_{SX}}\mu(S, X)}.
\end{align*}
Each posterior draw of $\mu$ gives a PAF draw, so it is an accounting functional like $\tau_Q$, with $\mu$ from the likelihood and $Q_{SX}$ chosen by the analyst. Only smokers contribute excess risk, so the PAF depends on smoking prevalence as well as the age and sex distribution that sets baseline risk. The covariate margin $Q$ alone is insufficient; the conditional risk ratio $\mu(1,x)/\mu(0,x)$ instead describes the model at one $x$. With constant risk ratio $\theta$ and smoking prevalence $p$ at every $x$, the PAF is $p(\theta - 1)/\{1 + p(\theta - 1)\}$: for $\theta = 9$, it is 0.44, 0.67 and 0.80 at $p = 0.1, 0.25, 0.5$.

\paragraph{Exercise~\ref{ex:ch13:4}.} Twin data bound how strongly any one genotype can be tied to smoking, and so restrict the pair $(f_1, \mathrm{PR})$ of the chapter's illustration. They do not fix PR alone, and they say nothing about $R$. Use the threshold model. Each twin smokes when a liability crosses a threshold, $S = \mathbf 1\{\delta U + e > c\}$, where $U$ is the binary genotype with prevalence $f$, and the residuals $(e_1, e_2)$ of a pair are bivariate normal with correlation $r \geq 0$, independent of the genotypes. Let $p_u = \Prob(S = 1 \mid U = u)$ and let $\bar p$ be the share of smokers. Split the covariance within a pair as
\[
\mathrm{Cov}(S_1, S_2) = \rho_U f(1 - f)(p_1 - p_0)^2 + \E\,K(U_1, U_2),
\]
where $\rho_U$ is the correlation of $U$ within the pair and $K(u, v)$ is the covariance of $S_1$ and $S_2$ given $U_1 = u$ and $U_2 = v$. Identical twins have $\rho_U = 1$. For fraternal twins under random mating, a function of one autosomal locus has $\rho_U$ between 1/4 and 1/2. Take $r$ no smaller for identical twins; then the residual term is at least as large for them. With a common $r$ the two expectations differ by $\Prob(U_1 = 1, U_2 = 0)\{K(1, 1) + K(0, 0) - 2K(1, 0)\} \geq 0$, because $K$ is the integral over $(0, r)$ of the bivariate normal density in the two thresholds, which is a positive definite kernel for correlations in $(0, r)$ by Mehler's formula. A larger $r$ for identical twins only raises $K$. Let $c$ be the concordance, the probability that a twin smokes given that the co-twin does, so that $\mathrm{Cov}(S_1, S_2) = \bar p(c - \bar p)$. Dividing by $\bar p(1 - \bar p)$ gives the twin correlation $\phi = (c - \bar p)/(1 - \bar p)$ and
\[
\phi_{US}^2 \leq 2(\phi_{\mathrm{MZ}} - \phi_{\mathrm{DZ}}), \qquad \phi_{US}^2 = \frac{f(1 - f)(p_1 - p_0)^2}{\bar p(1 - \bar p)} = \frac{(f_1 - f_0)^2\,\bar p(1 - \bar p)}{f(1 - f)},
\]
with $f = \bar pf_1 + (1 - \bar p)f_0$ and $f_0 = f_1/\mathrm{PR}$. The last equality uses $\mathrm{Cov}(U, S) = f(1 - f)(p_1 - p_0) = \bar p(1 - \bar p)(f_1 - f_0)$.

The prior follows. Take as an illustration, from no particular study, $\bar p = 0.5$ and concordances 0.80 and 0.69, so that $\phi_{\mathrm{MZ}} = 0.60$, $\phi_{\mathrm{DZ}} = 0.38$ and the bound is 0.44. Start from the chapter's sceptical prior ($f_1$ uniform, $\log\mathrm{PR}$ and $\log R$ half normal with scale 2) and keep only draws with $\phi_{US}^2 \leq w \times 0.44$, where $w$ is the share of the heritable resemblance that this one genotype carries. The twins give the 0.44, with its sampling error if the registry is small, and $w$ is a judgement. From 2,000,000 draws (seed 1958), which reproduce the chapter's $\Prob(B \geq 9) = 0.022$ when nothing is removed:
\begin{center}
\begin{tabular}{cccc}
\toprule
$w$ & Median PR & 95th percentile of PR & $\Prob(B \geq 9)$ \\
\midrule
no restriction & 3.9 & 50 & 0.022 \\
1 & 3.1 & 38 & 0.011 \\
1/4 & 2.0 & 23 & 0.0018 \\
1/10 & 1.6 & 16 & 0.0004 \\
\bottomrule
\end{tabular}
\end{center}
Moving $\bar p$ to 0.3 or 0.7 at $w = 1$ gives medians 3.1 and 3.3 and 95th percentiles 34 and 42. The tail of PR stays long because a rare genotype can have a huge prevalence ratio and still explain little of the smoking: at $f_1 = 0.1$ and $\bar p = 0.5$ even $\mathrm{PR} = \infty$ gives only $\phi_{US}^2 = 0.053$. The bound therefore cuts the middle of the prior, common genotypes with large PR, and leaves the rare ones, which then need a very large $R$ to reach Cornfield's region. This limits what twins can do here. With $w = 1/10$ the median matches the chapter's twin informed 1.6, but its 95th percentile of 4 and its probability of $4.4 \times 10^{-5}$ rest on a further judgement against rare major genotypes, one that the twins do not supply, which is why the chapter calls the 4 illustrative. Two biases act in opposite directions. If identical twins share more environment than fraternal twins the bound is too generous, so it stays valid. Assortative mating raises $\rho_U$ above 1/2 in fraternal twins and makes the bound too tight, so the analyst should relax it if mating on smoking is plausible. And since concordance in smoking is silent on cancer, the prior on $R$ stays as it was, as the chapter says.

\paragraph{Exercise~\ref{ex:ch13:5}.} Under the null and the assumption $\Prob(Y = 1 \mid A, U, H) = r(U, H)$, the risk in a group is the average of $r$ over that group: $\Prob(Y = 1 \mid A = a, H = h) = \sum_u r(u, h)\Prob(U = u \mid A = a, H = h) = r(0, h) + \{r(1, h) - r(0, h)\}\Prob(U = 1 \mid A = a, H = h)$. Since $r(1, h) > r(0, h)$, quitters and continuing smokers with history $h$ differ in risk if and only if $\Prob(U = 1 \mid A, H = h)$ differs between them, that is, only if $U$ is associated with quitting in that stratum. The assumption is needed: if some other cause of quitting also raised cancer risk, say quitting when $V = 1$ and risk $0.05 + 0.10U + 0.20V$ with $U$ and $V$ independent, the groups would differ in risk although $U$ has the same prevalence in both. This can happen even when $U$ has no effect on the decision to quit. Let $U$ and $V$ be independent Bernoulli(1/2), let a person start smoking if $U = 1$ or $V = 1$, quit if $V = 1$, and let cancer depend on $U$ alone. Among smokers, $U$ has prevalence 1/2 among quitters and 1 among continuing smokers, so continuing smokers have higher risk although smoking does nothing and $U$ does not affect quitting; the association comes from selection into smoking. If $U \indep A \mid H$, the null gives $\Prob(Y = 1 \mid A = 1, H) = \Prob(Y = 1 \mid A = 0, H)$, and an observed fall in risk among quitters with the same history contradicts the constitutional hypothesis under that assumption, which is the test. The independence involves the unmeasured $U$ and remains an assumption; what the data test is its implication, equal risks for quitters and continuing smokers within each history stratum. Risks over a fixed horizon serve better than hazards here, because the ratio of group hazards drifts as high risk people are depleted.

\paragraph{Exercise~\ref{ex:ch13:6}.} A smoker at dose $d$ falls ill with probability $\mathrm{RR}(d)\,r_0$, where $r_0$ is the common risk of nonsmokers, so among the smokers who fall ill the share at dose $d$ is $w_d\,\mathrm{RR}(d)/\mathrm{RR}$ with $\mathrm{RR} = \sum_d w_d \mathrm{RR}(d)$. The average of the dose specific lower bounds over those smokers is
\begin{align*}
\sum_d \frac{w_d\,\mathrm{RR}(d)}{\mathrm{RR}}\Bigl\{1 - \frac{1}{\mathrm{RR}(d)}\Bigr\} &= \sum_d \frac{w_d\{\mathrm{RR}(d) - 1\}}{\mathrm{RR}} \\
&= \frac{\mathrm{RR} - 1}{\mathrm{RR}} = 1 - \frac{1}{\mathrm{RR}},
\end{align*}
where the second equality uses $\sum_d w_d = 1$. The dose specific bound $1 - 1/\mathrm{RR}(d)$ exceeds $1 - 1/\mathrm{RR}$ exactly when $\mathrm{RR}(d) > \mathrm{RR}$, that is for plaintiffs whose dose carries more than the average risk ratio of smokers. When the risk ratio rises with dose, heavy smokers are in this group and light smokers are not, so a single ratio understates the bound for the first and overstates it for the second.

\section*{Chapter~\ref{ch:did}: Difference in Differences with the Horseshoe}

\paragraph{Exercise~\ref{ex:chdid:1}.} Number the periods 1 to 8, call the cohorts A (adopting in period 3) and B (period 6), and let A have a share $\pi$ of the units. In a balanced panel the residual of $D_{it}$ on unit and period dummies is $\tilde D_{it} = D_{it} - \bar D_{i\cdot} - \bar D_{\cdot t} + \bar D_{\cdot\cdot}$. The unit means are $3/4$ for A and $3/8$ for B, the period means are $0$, $\pi$ and $1$ in periods 1 and 2, 3 to 5 and 6 to 8, and the grand mean is $\tfrac38(1 + \pi)$. In periods 6 to 8 every unit is treated, so $\tilde D_{it} = \bar D_{\cdot\cdot} - \bar D_{i\cdot}$, which is $-\tfrac38(1 - \pi)$ for A and $\tfrac38\pi$ for B; for A in periods 3 to 5, $\tilde D_{it} = \tfrac58(1 - \pi)$. Weighting each cell by its cohort share and normalizing to sum to one over the treated cells, the common factor $\pi(1 - \pi)/8$ cancels and
\[
\begin{gathered}
\omega_{A,e} = \tfrac13\ (e = 0, 1, 2), \qquad \omega_{A,e} = -\tfrac15\ (e = 3, 4, 5), \\
\omega_{B,e} = \tfrac15\ (e = 0, 1, 2),
\end{gathered}
\]
for every $\pi$ strictly between 0 and 1; the weights sum to $1 - \tfrac35 + \tfrac35 = 1$. The negative weights are on the last three cells of the early cohort, in periods 6 to 8 after the later cohort has adopted: cohort A has been treated for a larger fraction of the panel than average, so its residual treatment is negative there. These are the cells with the largest effects when effects grow with exposure. With $\tau_{g,e} = 1 + e$, $\beta_{\mathrm{TWFE}} = \tfrac13(1 + 2 + 3) - \tfrac15(4 + 5 + 6) + \tfrac15(1 + 2 + 3) = 0.2$, against an average of 3 over treated unit periods when the cohorts are equal in size.

\paragraph{Exercise~\ref{ex:chdid:2}.} Write $\bar Y_{g,s}$ for the mean outcome at time $s$ of cohort $g$ and $\bar Y_{\infty,s}$ for that of the never treated units, and take the base period $g - 1$ to be untreated (no anticipation). For $t = g + e$ the estimator is
\[
\hat\tau_{g,e} = (\bar Y_{g,t} - \bar Y_{g,g-1}) - (\bar Y_{\infty,t} - \bar Y_{\infty,g-1}).
\]
By \eqref{eq:hsdid0}, the change of a unit in cohort $g$ is $Y_{it} - Y_{i,g-1} = \tau_{g,e} + \gamma_t - \gamma_{g-1} + \delta_i(\tilde t_t - \tilde t_{g-1}) + \varepsilon_{it} - \varepsilon_{i,g-1}$, since $\alpha_i$ cancels, and $\tilde t_t - \tilde t_{g-1} = (t - g + 1)/s_t$ because standardizing subtracts the same mean from both times. A never treated unit has the same change without $\tau_{g,e}$. Averaging within each group, subtracting, and using $\E\varepsilon_{it} = 0$, the period effects cancel and
\[
\E\,\hat\tau_{g,e} = \tau_{g,e} + \frac{t - g + 1}{s_t}\,(\bar\delta_g - \bar\delta_\infty),
\]
where $\bar\delta_g$ and $\bar\delta_\infty$ average $\delta_i$ over cohort $g$ and over the never treated units. Under parallel trends, \eqref{eq:twoway}, every $\delta_i = 0$ and the estimator is unbiased. Under the trend departure model with $\bar\delta_\infty = 0$ the bias is $(t - g + 1)\bar\delta_g/s_t$, which grows linearly with the distance from the base period. 

\paragraph{Exercise~\ref{ex:chdid:3}.} Keep the model \eqref{eq:hsdid0} to \eqref{eq:hsdid2} and add two things. Covariates enter with a linear term and a flat prior: the twelve time varying controls in the distributed file (black and white men aged 15 to 24 and 25 to 44, log public welfare and subsidy spending, log police, unemployment, poverty, log income, log prisoners and its lag), standardized. Population weights enter as a variance model, $\mathrm{Var}(\varepsilon_{it}) = \sigma^2/w_i$, where $w_i$ is the state's average population over the panel (the file's \texttt{popwt}, constant within a state) divided by the mean of these averages, so the largest weight is 6.1. This needs only the existing sampler applied to the data and design multiplied by $w_i^{1/2}$. In a posterior a weight is a statement about the noise, so it can be checked. The logarithm of the squared residual from a least squares fit with state effects, year effects, state trends and the post indicator falls with slope $-1.08$ on log population across state years, and $-0.97$ across state means, close to the $-1$ that weighting assumes and that a count based rate implies. Each fit below runs four chains of 20,000 sweeps from dispersed starts (generators from SeedSequence(20261101)), discards the first 5,000 of each, and has split $\hat R$ at most 1.001. The chains differ in the positive probability by at most 0.02. The summary keeps the chapter's target, the average over treated unit periods.
\begin{center}
\small
\small
\begin{tabular}{@{}p{4.4cm}cccc@{}}
\toprule
Fit & Two way & Horseshoe & 95\% interval & $\Prob(\mathrm{ATT} > 0)$ \\
\midrule
As in the equal noise rows of Table~\ref{tab:castle} & 0.069 & 0.025 & $[-0.047, 0.095]$ & 0.75 \\
Covariates & 0.069 & 0.023 & $[-0.050, 0.096]$ & 0.73 \\
Population weights & 0.076 & 0.079 & $[0.027, 0.129]$ & 0.998 \\
Weights and covariates & 0.080 & 0.084 & $[0.031, 0.136]$ & 0.999 \\
Weights, covariates, region by year & 0.082 & 0.073 & $[0.016, 0.129]$ & 0.993 \\
\bottomrule
\end{tabular}
\end{center}
The first row is a rerun of the chapter's model at the same length with other seeds, and its 0.75 sits within the spread of its own chains, 0.743 to 0.762, of the chapter's 0.76. The two way column is weighted least squares with the same covariates, and it comes close to the published eight percent once the weights are in. Region by year effects add 30 identified columns, four regions by eleven years less those collinear with the state and year effects. Averaging over treated population instead of treated unit periods moves the horseshoe estimate from 0.084 to 0.090 in the fourth row.

The covariates do almost nothing, and the weights do almost everything: the horseshoe rises from a third of the two way coefficient to match it, suggesting that the downward trend departures belong mostly to small states. The residuals do not support the equal noise model behind the chapter's three to one verdict. To allow a noise floor for large states, we fitted $\mathrm{Var}(\varepsilon_{it}) = \sigma^2\{(1 - \eta) + \eta/w_i\}$ with $\eta$ uniform on $(0, 1)$, which nests both models (same chain lengths, $\eta$ drawn on a grid of 41 points). The posterior of $\eta$ has mean 0.92 and 95\% interval 0.80 to 1. The effect is 0.076, interval 0.023 to 0.128, probability 0.997, or 0.073 and 0.992 with covariates and region by year, so the conclusion does not rest on the exact variance law.

With weights the evidence is strong, and the posterior is close to Cheng and Hoekstra's estimate. It leans on large adopters. In the weighted fit without covariates, dropping the largest adopter (18 percent of adopting population) raises the estimate to 0.094; dropping the single 2006 adopter (14 percent) lowers it to 0.049, interval $-0.007$ to 0.103, probability 0.96. Dropping both gives 0.064, interval 0.003 to 0.122, probability 0.98. Population weights set noise precision, while the cohort effect still spreads a large state's experience to its cohort. Police and prisoners may respond to the law, though the covariates move the estimate little. No refit tests an adopter specific shock starting at adoption.

\paragraph{Exercise~\ref{ex:chdid:4}.} With $N$ units and $T$ periods, write the untreated outcome as $Y_{it}(\infty) = \alpha_i + \gamma_t + f_t^\top\delta_i + \varepsilon_{it}$ with $f_t, \delta_i \in \R^2$; the trend departure of \eqref{eq:hsdid0} is the case of one known factor, $f_t = \tilde t$. The loadings get the horseshoe, $\delta_{ik} \sim N(0, \lambda_{ik}^2\tau_{\mathrm{g},\delta}^2\sigma^2)$ with half Cauchy $\lambda_{ik}$ and $\tau_{\mathrm{g},\delta}$, the factors get $f_t \sim N(0, I_2)$, and the other priors stay. Let $b$ collect $\alpha_i$, $\gamma_t$, $\theta_e$, $u_{g,e}$ and the lead coefficients, and let $r_{it} = Y_{it} - z_{it}^\top b$ with $z_{it}$ the design row. A sweep of the sampler of Section~\ref{sec:hsregression} has five steps.
\begin{enumerate}
\item Draw $b$ as in step 1 of that sampler, with response $Y_{it} - f_t^\top\delta_i$ and the prior precisions of the unchanged blocks (zero for $\alpha_i$, $\gamma_t$ and $\theta_e$).
\item Draw each $\delta_i$. With $F$ the $T \times 2$ matrix with rows $f_t^\top$, $r_i = F\delta_i + \varepsilon_i$, so $\delta_i \mid \cdot \sim N(A_i^{-1}F^\top r_i,\ \sigma^2A_i^{-1})$ with $A_i = F^\top F + \tau_{\mathrm{g},\delta}^{-2}\,\mathrm{diag}(\lambda_{i1}^{-2}, \lambda_{i2}^{-2})$.
\item Draw each $f_t$. With $\Delta$ the $N \times 2$ matrix with rows $\delta_i^\top$, $r_{\cdot t} = \Delta f_t + \varepsilon_{\cdot t}$, so $f_t \mid \cdot \sim N(B^{-1}\Delta^\top r_{\cdot t},\ \sigma^2B^{-1})$ with $B = \Delta^\top\Delta + \sigma^2I_2$.
\item Draw $\sigma^2$ as in step 2, with shape $(n + p)/2$, where $n = NT$ and $p$ counts all penalized coefficients including the $2N$ loadings, and with the residuals $r_{it} - f_t^\top\delta_i$ and the loading term $\sum_{ik}\delta_{ik}^2/(\lambda_{ik}^2\tau_{\mathrm{g},\delta}^2)$ in the rate.
\item Draw the scales of the cohort deviations and of the lead coefficients as in steps 3 and 4, each block from its own coefficients, and the loading scales in the same way with the $2N$ loadings as the coefficients: $\lambda_{ik}^2 \sim \mathrm{IG}(1,\ \nu_{ik}^{-1} + \delta_{ik}^2/(2\tau_{\mathrm{g},\delta}^2\sigma^2))$, $\nu_{ik} \sim \mathrm{IG}(1,\ 1 + \lambda_{ik}^{-2})$, $\tau_{\mathrm{g},\delta}^2 \sim \mathrm{IG}((2N + 1)/2,\ \xi^{-1} + \sum_{ik}\delta_{ik}^2/(2\lambda_{ik}^2\sigma^2))$ and $\xi \sim \mathrm{IG}(1,\ 1 + \tau_{\mathrm{g},\delta}^{-2})$.
\end{enumerate}
The likelihood is unchanged when $c$ is added to every $\delta_i$ and $f_t^\top c$ is subtracted from $\gamma_t$, and it does not fix the scale, order, signs or rotation of the factors. The prior on $f_t$ fixes the scale, and the horseshoe, which is not rotation invariant, pins the shift and favours rotations in which loadings are sparse. Order and signs stay arbitrary, so report quantities that do not depend on them: the counterfactual paths $\alpha_i + \gamma_t + f_t^\top\delta_i$ and the effect on the treated built from them.

% Chapter 15 (m11.tex, ch:k401). Solutions for ex:chk401:1, 3, 4, 5 and ex:chlalonde:5 are in ch15-t.tex and ch15-ch16b.tex; ex:chk401:2 and ex:chlalonde:1, 2 have none (see the log).

\paragraph{Exercise~\ref{ex:chdid:5}.} Use the chapter's design: 40 units over periods 0 to 15, each adopting in period 4, 8 or 12 or never with probability one quarter, a share $p = 0.1$ of units departing from parallel trends by $0.15(t - 8)$, noise standard deviation 0.5, and effects $\tau_{g,e} = 0.5 + 0.15e$. The average true effect is 1.12. Fit the same 500 panels (seeds as in the chapter's simulation) with two parametrizations, 1,500 sweeps with 500 discarded, and the same leads and unit trends. The first has one flat common effect $\bar\theta$ with $\tau_{g,e} = \bar\theta + u_{g,e}$. The second is the chapter's path, $\tau_{g,e} = \theta_e + u_{g,e}$ with $\theta_e$ flat. Both have the horseshoe on $u_{g,e}$.
\begin{center}
\small
\begin{tabular}{@{}lcccc@{}}
\toprule
Parametrization & Bias & RMSE & Covered of 500 & Mean width \\
\midrule
Single effect, 40 units & $-0.266$ & 0.319 & 268 & 0.53 \\
Path, 40 units & 0.000 & 0.092 & 488 & 0.42 \\
Single effect, 200 units & $-0.031$ & 0.052 & 468 & 0.19 \\
Path, 200 units & 0.000 & 0.040 & 487 & 0.18 \\
\bottomrule
\end{tabular}
\end{center}
At 40 units the single effect fails on every count, and its intervals are wider than those of the path and still miss. Sparsity explains it. The cohorts are observed at event times 0 to 11, 0 to 7 and 0 to 3, which gives $12 + 8 + 4 = 24$ cells. With one common effect, $u_{g,e} = 0.5 + 0.15e - \bar\theta$ can vanish at one event time at most, which holds at most three cells, so at least 21 of the 24 deviations are nonzero. With the path $\theta_e = 0.5 + 0.15e$ all 24 are zero. The horseshoe is built for many zeros and a few large signals, and the single effect hands it a dense vector instead. Many of its entries are also small against the noise: adjacent event times differ by 0.15, a cell mean over about 10 units has standard deviation $0.5/\sqrt{10} = 0.16$, and at $\bar\theta = 1.12$ nine of the 24 deviations lie within two such standard deviations of zero. The prior pulls these deviations toward zero and the estimate toward the flat level, and the unit trends take up part of the growth, which lowers the effect further because a unit's fitted trend carries past adoption into its counterfactual. At 200 units a cell mean has standard deviation about 0.07 and only four deviations lie that close, so most stand clear of the noise and escape the shrinkage, and the failure nearly disappears.

To check that sparsity, and not the sampler or the panel, is the cause, replace the effects by the constant $\tau_{g,e} = 1.12$, with all else as above, and fit both parametrizations to 200 new panels of 40 units (SeedSequence(20261101)). Now the single effect has all 24 deviations zero, and the path has them zero as well but spends twelve flat coefficients on a constant. The single effect has bias 0.002 (standard error 0.006), root mean squared error 0.078, 192 of 200 intervals covering and mean width 0.34. The path has bias 0.006 (standard error 0.007), root mean squared error 0.094, 195 of 200 covering and mean width 0.42. The single effect now wins, so the same prior and sampler fail when the vector they shrink is dense and not when it is sparse. The path costs little when it is not needed and saves the fit when it is, which is why the flat coefficients belong on the common shape and the shrinkage on what is left.

\section*{Chapter~\ref{ch:k401}: Two Benchmarks}

\paragraph{Exercise~\ref{ex:chk401:1}.} Write the error as a scale mixture, $e_i = \sigma z_i/\sqrt{w_i}$ with $z_i \sim N(0, 1)$ and $w_i \sim \mathrm{Gamma}(\nu/2, \nu/2)$ in the shape and rate form. Given the coefficients and $\sigma$, the weight of household $i$ enters the joint density through $w_i^{\nu/2 - 1}e^{-\nu w_i/2}$ from the prior and $w_i^{1/2}\exp\{-w_i r_i^2/(2\sigma^2)\}$ from the likelihood, where $r_i = y_i - x_i^\top b$, so
\[
w_i \mid \text{rest} \sim \mathrm{Gamma}\Bigl(\frac{\nu + 1}{2}, \frac{\nu + r_i^2/\sigma^2}{2}\Bigr), \qquad \E(w_i \mid \text{rest}) = \frac{\nu + 1}{\nu + r_i^2/\sigma^2},
\]
independently over households. The other steps are those of the Gaussian sampler with $y_i$ and $x_i$ multiplied by $\sqrt{w_i}$: a weighted least squares draw inside the horseshoe step and a draw of $\sigma^2$ from the weighted residual sum of squares. At $\nu = 3$ a household whose residual is ten scales has expected weight $4/103 \approx 0.04$, so it barely moves the coefficients.

Given $\sigma$, the $t$ likelihood mode solves $\sum_i\psi_\nu(r_i)x_i = 0$, where $\psi_\nu(r) = \E(w \mid r)\,r = (\nu + 1)r/(\nu + r^2/\sigma^2)$ falls toward zero for large $|r|$. For outcomes symmetric about the regression line, every odd $\psi$ gives the same line, and the effect is the shift of the mean. Assets are right skewed: the Gaussian limit $\nu = \infty$, with $\psi(r) = r$, estimates the mean, while smaller $\nu$ discounts the right tail and moves the fitted location toward the centre. The common effect is 9.82 in the Gaussian limit, 7.85 at $\nu = 30$, 6.47 at 10, 5.28 at 5, 4.22 at 3, 3.36 at 2 and 2.05 at 1. Median regression gives 2.92; the lower quartile and ninetieth percentile give 2.00 and 13.19. The effect grows along the distribution, so discounting the right tail lowers the estimate; at one degree of freedom it falls below the median shift. Each answer is a correct summary of a different location, and the average effect of Table~\ref{tab:k401} is the Gaussian one.

\paragraph{Exercise~\ref{ex:chk401:2}.} The model is the chapter's coupled horseshoe with the quintile effects replaced by interactions with the whole basis: $y_i = a + \delta\hat e_i + z_i^\top\beta + T_i(\tau_0 + z_i^\top\gamma) + \varepsilon_i$, with $z_i$ the standardized basis, $\hat e_i$ the cross validated propensity, $T_i$ eligibility, and $a$, $\delta$, $\tau_0$ flat. The main effects keep the chapter's horseshoe on 54 terms. The interactions get a second one on the 49 effective terms (the squares of the five binary indicators are affine in them), with its own half Cauchy global scale $\gscale$ of scale one, learned from the data (posterior mean 0.012 in units of $\sigma$). The basis is centred, so $\tau_0$ is the average effect over all households, 9.04 [6.58, 11.45], and $\tau(x) = \tau_0 + z(x)^\top\gamma$. Four chains of 11,000 sweeps of the chapter's Gibbs sampler with one global scale per block, 1,000 discarded from each, seeds 4010100 to 4010103; split $\hat R \le 1.003$, effective sample size at least 1,260.

The effect $\tau(x)$ at income $g$ is reported as the average of $\tau(x_i)$ over the 300 households with income nearest $g$, which averages the other covariates over those households. At 15, 25, 35, 50 and 75 thousand dollars the posterior mean and 95\% interval are 2.45 [$-1.98$, 6.95], 6.85 [3.85, 9.87], 10.47 [7.83, 13.13], 13.98 [10.97, 17.00] and 16.23 [11.72, 20.79]. Averaged within the income quintiles the effects are 1.87, 5.67, 9.24, 12.50 and 15.90, against 1.41, 1.45, 6.76, 12.80 and 20.78 in Table~\ref{tab:k401q}. Both rise steeply from a small effect at low incomes, but the horseshoe line is smoother: higher in the second and third quintiles and 4.9 lower, two posterior standard deviations, in the top one. Three interactions have 95\% intervals that exclude zero: eligibility times income (6.19 per standard deviation of the term), times income and two earners ($-6.61$), and times income and a defined benefit pension (3.48, lower limit 0.005). The other 46 are shrunk, 43 of them to posterior means below 1.5 in absolute value.

\paragraph{Exercise~\ref{ex:chk401:3}.} The score is the one of Exercise~\ref{ex:chlalonde:3} with eligibility as the treatment. With $r = Y - \mu_0(X)$, $w = e(X)/\{1 - e(X)\}$ and $p = \Prob(T = 1)$,
\[
\psi = \frac1p\bigl[T(r - \tau_T) - (1 - T)\,w\,r\bigr],
\]
and the estimate sets the sample mean of $\psi$ to zero, with standard error $\{n^{-1}\,\mathrm{mean}(\psi^2)\}^{1/2}$. The computation uses the 9,915 households with the chapter's five fold cross fitting, the partition of Table~\ref{tab:k401} (seed 1) and seeds 0 to 4 for the spread. In each fold $\mu_0$ is fitted on the training controls and $e$ on all training households, by forests (200 trees, minimum leaf 5, four of the nine standardized covariates per split) or, on the 54 term basis, by a lasso with five fold cross validation for $\mu_0$ and a cross validated logistic regression for $e$; $e$ is clipped to $[0.01, 0.99]$. Forests give 9.91, standard error 1.75, interval [6.47, 13.34]; the lasso version gives 10.76, standard error 1.74, interval [7.35, 14.17]. Over the five partitions the estimates run from 9.91 to 11.13 for forests (median 10.78) and from 10.76 to 11.62 for the lasso (median 11.23), and the tabled partition is the lowest for both, as it was for the partially linear score.

The accounting value 11.65, interval [9.27, 14.16], is 0.9 above the forest median and 0.4 above the lasso median, less than one DML standard error. Most of the gap is in the top income quintile, which carries weight 0.32: the same score within quintiles gives 16.84 there with forests and 18.90 with the lasso, against 20.78 in Table~\ref{tab:k401q}, with a standard error of 5.1 against a posterior standard deviation of 2.2. At the tabled partition that quintile accounts for about 1.3 of the 1.7 gap with forests and 0.6 of 0.9 with the lasso. DML is less precise overall, 1.7 against 1.25, because it weights the controls by the odds $e/(1 - e)$, up to 12.7 with forests. Both estimates exceed the partially linear ones of Table~\ref{tab:k401}, 8.59 and 9.52, because the treated are richer: the score weights the quintiles by the treated shares, 0.07 to 0.32, not by $e(1 - e)$. The two analyses agree within the noise of the DML estimate.

\paragraph{Exercise~\ref{ex:chk401:4}.} The propensity is the cross validated logistic regression on the 54 term basis that the coupled prior uses, with range 0.020 to 0.889. Only 17 of the 9,915 households fall below 0.05, 15 in the lowest income quintile and 2 in the second, and one is eligible: overlap at the low end is adequate (Section~\ref{sec:k401overlap}). Rerunning the chapter's coupled horseshoe with quintile effects (its sampler, seed 401, 2,000 draws kept) reproduces Table~\ref{tab:k401q} and the accountings of Section~\ref{sec:k401acc}, and averaging the quintile draws with the quintile shares of the 9,898 retained households gives 8.65, interval [6.22, 11.18], against 8.64 for all households, with 11.65 for the retained treated (11.65 before) and 6.87 for the untreated (6.86). For DML, the partially linear score of Table~\ref{tab:k401} refitted to the retained households with the same learners and five folds gives, on the partition with seed 1, 8.99 with forests (standard error 1.29, interval [6.46, 11.51]) and 9.80 with the lasso (1.32, [7.23, 12.38]), against 8.59 and 9.52. The medians over seeds 0 to 4 are 8.95 and 9.80 against 9.00 and 9.82, so the shifts at seed 1 come from the different fold assignment of a smaller sample.

Trimming at 0.05 changes nothing a reader could use: the 17 households have $e(1 - e)$ of at most 0.047, against a mean of 0.199 among the others. 

\paragraph{Exercise~\ref{ex:chk401:5}.} The argument is plausible for the plan and weak for the worker. Employers set eligibility, and at a given income a worker seldom takes a job for its plan. But eligibility is far from random before adjustment (a raw difference of 19.56 against adjusted estimates of 5.90 to 9.84), the adjustment depends on functional form, and what remains is a claim about nine covariates. A taste for saving is the obvious omitted variable: it raises assets at any eligibility, and savers may seek out employers with plans. Both paths push the estimate up, so it is more likely too high than too low. IRA ownership and pension coverage proxy the taste but may respond to eligibility, so controlling for them is not innocent. The argument supports the sign and the order of magnitude, not the second digit.

For the sensitivity analysis take the linear form of Section~\ref{sec:sensitivity} (Figure~\ref{fig:sens}). Let $U$ be a standardized taste for saving, $a$ its mean difference between eligible and ineligible households with the same covariates, and $b$ its effect on assets in thousands of dollars per standard deviation. Omitting $U$ shifts the estimated effect by $ab$, so the effect given $(a, b)$ is $\tau - ab$. Fit the chapter's coupled horseshoe once and keep the draws of the average effect, the effect on the treated and the five quintile effects. The data cannot update $(a, b)$, so draw them from the prior, subtract $ab$ from each draw of the effect and report the result. Because the worry is the sorting of savers into plans, take $a = m s_a |z_a|$ and $b = m s_b |z_b|$ with $z_a$, $z_b$ standard normal. Here $s_b$ is the coefficient of standardized income in the outcome model, since a taste for saving may predict assets as well as income does (partial $R^2$ 0.088), and $s_a$ is half the coefficient of eligibility in the regression of standardized income on eligibility, since employers set eligibility. Report $m \in \{0.5, 1, 2\}$, with the symmetric normal priors of Figure~\ref{fig:sens} as a second row for sorting the other way: the posterior mean and 95\% interval of the average effect and of the effect on the treated, the probability that each is positive, the multiple $m$ at which it falls to 0.95, and the same by quintile, where the two lowest effects are already indistinguishable from zero. The posterior of $(a, b)$ is its prior, so every figure is conditional on the stated beliefs.

\paragraph{Exercise~\ref{ex:chlalonde:1}.} Replace the 15,992 CPS controls by the 2,490 men of Dehejia's PSID comparison file and keep the 185 treated, the ten covariates, the logistic propensity and every fitting choice of Section~\ref{sec:lalonde}. Tree posteriors keep 1,000 draws. We ran four seeds (7, 11, 13, 17) because the section shows that the tree rows move with the seed; DML has forests of 300 trees, five folds and five partitions (seeds 1 to 5). Intervals are from the first seed or partition. All values are thousands of dollars, and the experimental benchmark is 1.79.
\begin{center}
\small
\begin{tabular}{@{}p{3.7cm}ccc@{}}
\toprule
Method & CPS (Table~\ref{tab:lalonde}) & PSID, by seed & PSID interval \\
\midrule
Difference in means & $-8.50$ & $-15.20$ & \\
Least squares & 1.07 & 0.00 & $[-1.67, 1.68]$ \\
DML, augmented weighting for the treated, forests & 1.51 & 1.86, 1.69, 2.14, 2.08, 2.15 & $[0.24, 3.49]$ \\
Hill's BART & 1.54 & 1.05, 1.26, 0.98, 1.07 & $[-1.11, 3.23]$ \\
BART with estimated propensity & 1.51 & 1.52, 1.54, 1.53, 1.29 & $[-0.95, 3.84]$ \\
BART after common support check & 1.39 & 1.10, 1.26, 1.02, 1.07 & $[-1.06, 3.30]$ \\
Bayesian causal forest & 1.68 & 1.35, 1.51, 1.50, 1.25 & $[-1.06, 3.81]$ \\
\bottomrule
\end{tabular}
\end{center}
The spread is wider on every measure. Excluding the raw difference, the rows run from 0.00 to 2.15, against 1.07 to 1.68 with the CPS. The flexible methods alone run from 0.98 to 2.15 across seeds, against a range of 0.29 at the printed CPS values and 0.35 once the seed ranges the section reports are included. Hill's BART has an interval of width 4.3 against 1.3, and DML one of width 3.3 against 2.9. The raw gap is nearly twice as large, because PSID men earn more than CPS respondents do, and least squares falls by about one thousand dollars to zero. Propensity information still separates the tree rows: the two that use it, BART with the estimated propensity (1.29 to 1.54) and the causal forest (1.25 to 1.51), sit above plain BART (0.98 to 1.26), and each moves by about a quarter of a unit across seeds. The seeds repeat the same data with 1,000 kept draws and no convergence check, so that movement is computational, and these runs do not settle any tree row's posterior mean more finely than about a quarter of a unit.

Weak overlap is the likely reason. Of the PSID controls, 55 percent have estimated propensity below one in a thousand and 41 percent lie below the lowest treated propensity, while the treated median is 0.55, against 0.32 with the CPS. There are also fewer controls, so only a few resemble the trainees. The largest odds weight in DML ranges from 9 to 33 across partitions, against 8 with the CPS, which fits the movement of its estimate from 1.69 to 2.15. The common support check is not much help: it drops at most one treated man, since its rule compares each man's posterior spread with that of the controls. The flexible rows all contain 1.79 in their intervals, but with intervals this wide that is a weak test. Section~\ref{sec:crit:applied} runs the placebo tests on this same comparison group, and they fail by 2.4 to 7.3 thousand dollars for every procedure. Closeness to the benchmark here should not be read as evidence that the PSID comparison is ignorable given these covariates, a point Smith and Todd (2005) make about such comparisons.

\paragraph{Exercise~\ref{ex:chlalonde:2}.} The sample is the chapter's, 185 NSW trainees and 15,992 CPS controls, with 1978 earnings in thousands of dollars as the outcome. The basis is the quadratic expansion of the ten covariates (squares and pairwise products of the standardized variables, less the six squares of binary variables), with each year's linear and squared earnings terms replaced by a cubic B-spline in that year's earnings: interior knots at 1, 2.5, 5, 10 and 17.5 thousand dollars, boundary knots at zero and above the maximum, the function that is nonzero at zero earnings dropped, 71 standardized terms in all. The model is the coupled horseshoe of Chapter~\ref{ch:horseshoe}, $y_i = a + \tau T_i + \delta\hat m_i + z_i^\top\beta + \varepsilon_i$, with $a$, $\tau$, $\delta$ flat, a half Cauchy global scale, and $\hat m_i$ the cross validated logistic propensity on the basis. Each fit reported below runs four chains of 6,000 sweeps of the Gibbs sampler of Chapter~\ref{ch:horseshoe}, 1,000 discarded from each, with four seeds of its own (19860100 to 19860103 for this fit; the seven variants use 19860200 to 19860803, four to a fit); split $\hat R$ is below 1.001 and every effective sample size above 10,000. Two readings of the effect on the treated are reported: $\tau$, whose interval carries the sampling variability of the trainees' earnings as DML's does, and the average over the 185 trainees of observed minus fitted control earnings, which holds those earnings fixed as the BART rows do.

With the propensity itself as $\hat m$, $\tau$ is 1.01 [$-0.34$, 2.35] and the averaged difference 1.01 [0.11, 1.93]. Table~\ref{tab:lalonde} gives 1.54 [0.85, 2.16] for Hill's BART, 1.51 [0.77, 2.28] with the propensity and 1.39 [0.70, 2.03] after the support check; the experiment gives 1.79, and BART on its data 1.55 [0.74, 2.32]. So the coupled horseshoe sits about half a unit below the BART rows, inside the experimental band, and most variants leave it there: with no $\hat m$ at all $\tau$ is 1.00 [$-0.19$, 2.19], with the logit of the propensity in its place 1.03 [$-0.18$, 2.25], and with the quadratic basis in place of the splines 1.04 [$-0.31$, 2.37], so the splines change nothing. Two choices move it. Clipping the logit at 0.001 and 0.999 gives 1.43 [0.18, 2.69] (averaged difference 1.43 [0.68, 2.16]; 1.45 with the quadratic basis), within 0.11 of all three BART rows. But 79 percent of the controls have a propensity below 0.001, so the clip gives all of them the same $\hat m$, and the agreement with BART rests on where the clip sits; no other level was tried. Letting the effect vary linearly in $\hat m$ gives 1.87 [0.34, 3.40] with the propensity and 1.68 [0.42, 2.93] with the clipped logit. All of these posterior means, from 1.00 to 1.87, lie inside the experimental band, so this sample cannot choose among them.

\paragraph{Exercise~\ref{ex:chlalonde:3}.} Let $p = \Prob(T = 1)$ and $\mu_0(x) = \E(Y \mid X = x, T = 0)$, so that under unconfoundedness $\tau_T = \E[T\{Y - \mu_0(X)\}]/p$. With a candidate $\tilde\mu_0$ in place of $\mu_0$ the treated term is off by $\E[T\{\mu_0(X) - \tilde\mu_0(X)\}] = \E[e(X)\{\mu_0(X) - \tilde\mu_0(X)\}]$. A correction from the controls with weight $w(X)$ has mean $\E[(1 - T)\,w(X)\{Y - \tilde\mu_0(X)\}] = \E[(1 - e(X))\,w(X)\{\mu_0(X) - \tilde\mu_0(X)\}]$, which cancels that error for every $\tilde\mu_0$ exactly when $(1 - e)\,w = e$. The score is therefore
\[
\psi = \frac1p\Bigl[T\{Y - \mu_0(X) - \tau_T\} - (1 - T)\,\frac{e(X)}{1 - e(X)}\,\{Y - \mu_0(X)\}\Bigr],
\]
with mean zero at $\tau_T$ for any candidate $\tilde\mu_0$ substituted for $\mu_0$ when $e$ is right. When $\mu_0$ is right, the mean is zero for any control weights because $\E[(1 - T)\{Y - \mu_0(X)\} \mid X] = 0$. The same score with eligibility as $T$ is the one Exercise~\ref{ex:chk401:3} asks for. The estimate is $n_1^{-1}\{\sum_i T_i r_i - \sum_i (1 - T_i)\hat w_i r_i\}$ with $r_i = Y_i - \hat\mu_0(X_i)$ and $\hat w_i = \hat e_i/(1 - \hat e_i)$.

The weights are the odds of being a trainee. Since $\E\{(1 - T)\,w(X)\} = \E\,e(X) = p$, they sum to about $n_1 = 185$ over the 15,992 controls, an average of 0.012. Half the controls have propensity below 0.001, hence odds of at most 0.0010 and a joint weight of at most $7{,}996 \times 0.001/0.999 \approx 8$, 4 percent of the total. The 1 percent above 0.18, about 160 households, carry at least $160 \times 0.18/0.82 \approx 35$, 19 percent. Under the forest propensity that DML uses the largest odds reached 8, so one household counts as much as eight trainees and nearly 700 times the average control. The control term is an average over the few households that resemble trainees, and a few earnings residuals drive its variance.

\paragraph{Exercise~\ref{ex:chlalonde:4}.} For each method, count the fraction of its 1,000 effect draws above 1, in thousands of dollars. A BART draw averages observed minus imputed earnings over the treated; a causal forest draw averages $\tau(x)$. Without the draws, use $\Prob(\tau_T > 1) \approx \Phi\{(m - 1)/s\}$ with printed mean $m$ and $s$ the 95\% interval width divided by 3.92. This gives 0.95 for Hill's BART ($m = 1.54$, $s = 0.33$), 0.91 for BART with the estimated propensity ($1.51$, $0.39$), 0.87 for BART after the common support check ($1.39$, $0.34$) and 0.86 for the Bayesian causal forest ($1.68$, $0.63$). The printed intervals are nearly symmetric, so the draws should refine these figures rather than change them. The causal forest has a larger mean than Hill's BART but a lower probability, because its interval is twice as wide. That width includes the sampling variability of the treated earnings, which BART holds fixed (Table~\ref{tab:lalonde}), so the two kinds of rows do not answer the same question.

% Chapter 16 (m12.tex, ch:nhefs). Solutions for ex:chnhefs:1, 2, 4 are in ch15-ch16b.tex; ex:chnhefs:3 and ex:chthornton:2 have none (see the log).

\paragraph{Exercise~\ref{ex:chlalonde:5}.} Two kinds of dependence need different tools. Averaging over a subset of the treated, the 179 with support against all 185 (1.39 against 1.54), changes the target and is an accounting of one set of draws. Changing the comparison group, or the rule that trims it, changes the data and needs its own fit. Let $s = 1, \ldots, S$ index such subsamples: the CPS and the PSID comparison groups (Exercise~\ref{ex:chlalonde:1}), each whole and with the comparison units below the lowest treated propensity dropped. Fit one model to all of them, say BART with the estimated propensity, so that the subsample is not mixed with the method, and keep the posterior mean $m_s$ and standard deviation $v_s$ of the effect on the treated, as the chapter does for its sensitivity curve (Figure~\ref{fig:lalonde_sens}). The second level, in thousands of dollars, is
\[
m_s \mid \tau_s \sim N(\tau_s, v_s^2), \qquad \tau_s \mid \mu, \sigma \sim N(\mu, \sigma^2),
\]
with a flat prior on $\mu$ and a half normal prior with scale 1 on $\sigma$. Report the posterior of $\sigma$, the dependence on the subsample as a number, with $\Prob(\sigma > 0.5 \mid \text{data})$; the posterior of $\mu$; the predictive distribution of the effect in a subsample not on the list, whose variance adds $\sigma^2$ to that of $\mu$; and the shrunken $\tau_s$ beside the raw $m_s$, which shows which subsamples pull the answer. A $\sigma$ near zero says the subsamples agree and Smith and Todd's warning does not bite for this model; a $\sigma$ comparable to the experimental 1.79 says it does, and then a single row of Table~\ref{tab:lalonde} understates the uncertainty. The subsamples share their treated participants, so the errors of the $m_s$ are positively correlated, which the model treats as independent: the posterior of $\sigma$ is pulled down and the interval for $\mu$ is too narrow. Repeating the analysis with the Bayesian causal forest gives a second $\sigma$, and agreement between the two says the dependence belongs to the data.

\section*{Chapter~\ref{ch:nhefs}: Two Applications in Health}

\paragraph{Exercise~\ref{ex:chnhefs:1}.} The data are the 1,566 smokers with complete data. The propensity model is the one of Hern\'an and Robins: logistic regression of quitting on sex, race, age and age squared, education (five levels), cigarettes per day and its square, years of smoking and its square, exercise and usual activity (three levels each), and 1971 weight and its square, by maximum likelihood. The weights are stabilized, $\Prob(T = t)/\Prob(T = t \mid X)$ with $\Prob(T = 1) = 0.257$; they reach 4.30 (propensities from 0.051 to 0.777) and the effective sample size is 1,446. Weighted least squares of weight change on quitting equals the difference of the weighted means and gives 3.44 kilograms, the value of Table~\ref{tab:nhefs} and of Hern\'an and Robins (3.4). For the interval the bootstrap resamples the 1,566 rows with replacement 1,000 times (generator seed 1971) and refits the propensity model and the weights in each resample. The 1,000 estimates have standard deviation 0.50 and a 95\% percentile interval of [2.41, 4.40] (normal interval [2.45, 4.43]). A sandwich variance that treats the estimated weights as known gives standard error 0.53 and [2.41, 4.47], the published 2.4 to 4.5; holding the weights fixed in the resamples gives 0.55, so ignoring the estimation of the weights is conservative.

Against Table~\ref{tab:nhefs}, the weighting interval is 2.0 kilograms wide, against about 1.7 for least squares, the horseshoe, BART and the causal forest and 2.25 for DML. The estimate lies within 0.21 of every regression and Bayesian row (3.23 to 3.48) and 0.27 below the DML estimate of 3.71. The bootstrap standard deviation, 0.50, exceeds the posterior standard deviations of 0.43 to 0.45 because weighting uses only the propensity and ignores what the covariates explain of weight change, which an outcome model uses.

\paragraph{Exercise~\ref{ex:chnhefs:2}.} Dropping the four interactions of quitting with age, cigarettes per day, years of smoking and weight leaves the design $(T, 1, \hat e, \text{basis})$ of the chapter's horseshoe, with the coefficient of $T$ a constant effect, and the same sampler (3,000 Gibbs iterations, 1,000 burn in, generator seed 1972). The posterior of the effect has mean 3.48, standard deviation 0.44 and 95\% interval [2.62, 4.33], against 3.48, 0.44 and [2.60, 4.34] for the average effect of Table~\ref{tab:nhefs}, which rerunning the chapter's seed (1971) reproduces. The posterior means differ by 0.001 and the standard deviations by less than 1 percent. The modifiers are standardized, so in the larger model the average effect is the coefficient of $T$ and the interactions drop out of it, and the horseshoe shrinks them: their posterior means are $-0.03$, 0.19, $-0.03$ and 0.14 kilograms per standard deviation, with standard deviations 0.25 to 0.29, and every 95\% interval contains zero. The unit level effects of the larger model have posterior means between 2.92 and 4.40, and its effects on the treated and untreated, 3.45 and 3.49, differ by 0.04, a difference the constant effect model sets to zero. The interactions cost nothing in the average and the data do not ask for them, as Figure~\ref{fig:nhefshet} suggests.

\paragraph{Exercise~\ref{ex:chnhefs:3}.} The data are the chapter's 1,566 smokers (403 quitters) followed from January 1983 to December 1992, 120 months; 291 died, 91 of them quitters. The NHEFS file that Hern\'an and Robins distribute has 1,629 smokers, so the counts here need not match theirs: the chapter keeps the 1,566 whose weight was measured in 1982. Each person contributes one record per month until death or month 120 (171,076 records), with the event indicator 1 in the month of death. The hazard is the pooled logistic model
\[
\logit h(t) = \beta_0 + \beta_q Q + B(t)^\top(\alpha + \delta Q) + \gamma^\top z,
\]
where $Q$ indicates quitting, $B(t)$ is a natural cubic spline of the month (boundary knots 0 and 119, interior knots 30, 60 and 90) and $z$ holds the chapter's baseline covariates, standardized: the four continuous ones and their squares, and dummies for sex, race, education, exercise and activity. The 28 coefficients have independent normal priors, $N(0, 10^2)$ for the intercept and $N(0, 2.5^2)$ for the rest. Hamiltonian Monte Carlo ran four chains of 600 warm up and 1,500 kept iterations (seed 1603), with no divergence after warm up, split $\hat R$ below 1.001 and effective sample sizes above 8,800, more than the 6,000 draws because successive Hamiltonian draws are negatively correlated. For each draw we set $Q = 1$, then $Q = 0$, for all 1,566 people, computed each person's survival $\prod_{t=0}^{119}\{1 - h(t)\}$ and averaged.

Ten year survival is 0.814 (95\% interval [0.783, 0.844]) under quitting and 0.818 ([0.798, 0.837]) without. The difference has posterior mean $-0.004$, standard deviation 0.018, 95\% interval $[-0.041, 0.030]$ and $\Prob(\text{difference} > 0) = 0.41$, so the interval rules out a gain above three points and a loss above about four, and the data say nothing sharper. The raw gap is $-5.4$ points (77.4\% alive against 82.8\%); the adjustment removes almost all of it, and quitters were 3.4 years older. The hazard ratio for quitting runs from 1.42 ([0.74, 2.46]) at month 12 to 0.49 ([0.21, 0.89]) at month 108. A quadratic in the month gives a difference of $-0.004$, prior standard deviation 10 gives $-0.003$ ($[-0.044, 0.034]$), and dropping the quitting by month terms gives $-0.010$ ($[-0.046, 0.026]$), each variant from two chains of 500 warm up and 1,000 kept iterations. All of it rests on unconfoundedness given the baseline covariates.

\paragraph{Exercise~\ref{ex:chnhefs:4}.} The computation uses the 2,000 draws of the chapter's horseshoe with effect modifiers, including the draw of the error variance $\sigma^2$ at each iteration. The woman is coded sex 1 (mean 1971 weight 64.7 kilograms, against 77.3 for sex 0), aged 45, smoking 20 cigarettes a day for 25 years, with the sample medians for the rest: race 0, education level 3, exercise level 1, usual activity level 1 and 1971 weight 69.23 kilograms. Her propensity of quitting from the Hern\'an and Robins logistic model is 0.204. For each draw and each action $t$ the mean function is $\mu_t = \beta_T t + \beta_0 + \beta_e \hat e + t\,\gamma^\top z_{1:4} + \beta^\top z$, with $z$ her standardized basis vector and $z_{1:4}$ its four continuous terms, and then $Y_t \sim N(\mu_t, \sigma^2)$, with 200 error draws for each posterior draw. Without quitting the mean function is 1.86 kilograms (posterior standard deviation 0.49, interval [0.89, 2.79]) and the predictive distribution has standard deviation 7.33, quartiles ($-3.1$, 1.9, 6.8) and 95\% interval [$-12.5$, 16.3]. With quitting the mean function is 5.31 (0.60, [4.13, 6.44]) and the predictive distribution has standard deviation 7.35, quartiles (0.4, 5.3, 10.3) and 95\% interval [$-9.1$, 19.7]. Her expected effect is 3.45, interval [2.58, 4.32], close to the average effect because the heterogeneity is small. The two predictive distributions overlap heavily, since the residual standard deviation of 7.3 is twice the effect. How her two potential changes relate to each other, and so whether she would gain more by quitting than by not quitting, depends on a correlation that the data do not identify (Chapter~\ref{ch:gente}).

\paragraph{Exercise~\ref{ex:chthornton:1}.} Constant baseline risk in distance and age gives $\beta_1 = \beta_2 = 0$. With $\delta = 0$ as well, every respondent has $p_{0i} = \expit(\alpha)$ and $p_{1i} = \expit(\alpha + \tau)$. The sample averages equal these probabilities, so
\[
\frac{\bar p_1/(1 - \bar p_1)}{\bar p_0/(1 - \bar p_0)} = \frac{e^{\alpha + \tau}}{e^\alpha} = e^\tau,
\]
the conditional odds ratio. The conditions remove both baseline variation and the interaction with distance that otherwise make the marginal odds ratio depend on averages of nonlinear probabilities.

\paragraph{Exercise~\ref{ex:chthornton:2}.} The model is the chapter's, on the same 2,816 respondents, with a village intercept: $\logit\Prob(\text{learned}_{ij} = 1) = x_{ij}^\top\beta + u_j$ with $u_j \sim N(0, \sigma^2)$. Here $x_{ij}$ holds the five terms of Section~\ref{sec:thornton} with independent normal priors of variance 10, and $\sigma$ has a half Cauchy prior with scale 1. There are 119 villages of 2 to 127 respondents (median 17), and the four respondents with no village number each get an intercept of their own, so there are 123 effects. The Polya Gamma Gibbs sampler draws $(\beta, u)$ in one block and $\sigma$ through the inverse gamma form of the half Cauchy. Four chains of 4,000 iterations with 1,000 dropped give 12,000 draws (seed 1602), with split $\hat R$ at most 1.004 and effective sample sizes of about 1,000 for $\sigma$ and above 7,000 for everything else. Dropping the village term (three chains of 3,000 iterations, 1,000 dropped) reproduces Table~\ref{tab:thornton}: risk difference 0.446, interval [0.405, 0.487].

The village scale $\sigma$ has posterior mean 0.60 and 95\% interval [0.43, 0.79], an intraclass correlation of about 0.10 on the logit scale. The sample risk difference, averaged over the respondents with their own village effects as in the chapter, has posterior mean 0.438 and interval [0.396, 0.480], against 0.447 and [0.404, 0.486]. It falls by 0.009, under half a posterior standard deviation (0.021), and the interval keeps its width, 0.084 against 0.082: the incentive was randomized and $\sigma$ is small, so the village term has little variance to remove, and 16 of the 123 groups have one arm and give no comparison within the village. The mean falls because the comparison now runs within villages: a linear probability model with village fixed effects gives 0.432 (robust standard error 0.022, or 0.023 clustered by village) against the raw 0.448, and the random effects model lands between them. Averaging over a new village drawn from $N(0, \sigma^2)$ gives 0.425 ([0.382, 0.469]) and fixing the typical village, $u = 0$, gives 0.450. The conditional odds ratio rises to 8.01 ([6.45, 9.91]) and the marginal one falls to 6.94 ([5.68, 8.45]), so the village effects widen the chapter's gap from 0.23 to 1.07.

\paragraph{Exercise~\ref{ex:chthornton:3}.} For $a \in \{0\} \cup [a_{\min}, a_{\max}]$, write $G(a) = p(a) - p(0)$ and $c(a) = a\,p(a)/G(a)$, the payout per additional learner. With $V$ in the units of $a$,
\[
N(a) = V\,G(a) - a\,p(a) = G(a)\{V - c(a)\}.
\]
The chapter's benefit to payout ratio is largest at $a_{\min}$ in every draw, so $c$ is smallest there. No incentive is best when $V \le c(a_{\min})$, and a positive amount above that threshold, which has its own posterior. In a draw where $p$ increases with amount, a differentiable interior maximizer satisfies
\[
(V - a)\,p'(a) = p(a),
\]
so $a < V$. Also check the endpoints and the logit kink at one unit.

For $V < V'$ and maximizers $a,a'$, adding their optimality inequalities gives $(V' - V)\{p(a') - p(a)\} \ge 0$. In each draw where $p$ increases with the amount, $a' \ge a$, so the posterior of the best amount moves up with $V$, from $a_{\min}$ toward $a_{\max}$. For each draw of $p(\cdot)$, maximize $N$ on a grid of amounts and summarize the 2,000 maximizers at each $V$ on a grid.

\paragraph{Exercise~\ref{ex:chthornton:4}.} Let $Z = 1$ if the respondent was offered any incentive, $T = 1$ if she learned her status and $Y$ the later behaviour, so that the theorem of Section~\ref{iv:sec:late} applies. Random assignment holds by design and the first stage is large: 78.9 percent of those offered an incentive learned their result against 34.1 percent of the others, 44.8 points more. Two assumptions must be added. Exclusion: the incentive moves $Y$ only through $T$. It is the weak point, because the voucher is paid on return to the centre, so a respondent who learns her result under $Z = 1$ also receives money; if the money itself changes the behaviour studied, the Wald ratio credits that income effect to the information. Monotonicity: nobody who would learn her status without an incentive declines with one. A payment makes this plausible, though an offer read as pressure could produce defiers. No interference between respondents is assumed as well.

With these assumptions the Wald ratio is the effect of learning for compliers, the 44.8 percent who learn only when offered an incentive. The 34.1 percent who learn anyway and the 21.1 percent who never learn contribute nothing, and the effect for them is not identified. The posterior of the effect is that of the ratio of the two reduced form differences, which is stable here, with a first stage near 0.45, unlike the weak instrument case of Section~\ref{iv:sec:weak}.

\section*{Chapter~\ref{ch:cdt}: Causal Decisions}

\paragraph{Exercise~\ref{ex:chcdt:1}.} Write $O$ for taking the opaque box only, $B$ for taking both, and $H$ for the event that the opaque box holds $M$. The predictor is right with probability $p$ whichever act is chosen, so $\Prob(H\mid O)=p$ and $\Prob(H\mid B)=1-p$. The payoffs are $M$ or $0$ for $O$ and $M+m$ or $m$ for $B$, according as $H$ holds or not. Evidential expected utility, \eqref{cdt:eq:eu} with the ordinary conditional probability, is
\[
\mathrm{EU}(O) = pM,\qquad \mathrm{EU}(B) = (1-p)(M+m) + pm = m + (1-p)M,
\]
so $O$ is preferred exactly when $p > (M+m)/(2M) = \tfrac12 + m/(2M)$. The causal calculation uses $\Prob(H\mid\mathrm{do}(A)) = q$, the agent's credence that the box is full, which does not depend on the act because the contents were fixed before the choice. Then
\[
\mathrm{EU}_{\mathrm{CDT}}(O) = qM,\qquad \mathrm{EU}_{\mathrm{CDT}}(B) = qM + m,
\]
and $B$ is better by $m>0$ for every $q$. The theories disagree exactly when $p>(M+m)/(2M)$, a nonempty range because $m<M$. For $M=10^6$ and $m=10^3$ it is $p>0.5005$, so a predictor barely better than a coin flip is enough. At $p=0.9$ the evidential utilities are $900{,}000$ for $O$ and $101{,}000$ for $B$.

\paragraph{Exercise~\ref{ex:chcdt:2}.} Take the payoffs to the agent: $R$ if both cooperate, $S$ if only she does, $T$ if only the twin does, $P$ if neither does. Causal reasoning holds the twin's act fixed. If the twin cooperates with probability $q$, a belief that does not depend on what the agent does, then
\[
\begin{split}
\mathrm{EU}_{\mathrm{CDT}}(D) - \mathrm{EU}_{\mathrm{CDT}}(C) &= \{qT + (1-q)P\} - \{qR + (1-q)S\}\\
&= q(T-R) + (1-q)(P-S) > 0
\end{split}
\]
for every $q\in[0,1]$, because $T>R$ and $P>S$. Defection dominates, so no belief about the twin makes cooperation the causal choice. Evidential reasoning conditions on the act. If the twin matches the agent's act with probability $p$ whichever act that is, $\mathrm{EU}(C) = pR + (1-p)S$ and $\mathrm{EU}(D) = pP + (1-p)T$. Cooperation is better exactly when $p(R-P) > (1-p)(T-S)$, that is, when
\[
p > p^\ast = \frac{T-S}{(T-S) + (R-P)} .
\]
Since $T-S = (T-R) + (R-P) + (P-S) > R-P > 0$, the threshold lies strictly between $\tfrac12$ and $1$. For $T=5$, $R=3$, $P=1$, $S=0$ it is $p^\ast = 5/7 \approx 0.71$.

\paragraph{Exercise~\ref{ex:chcdt:3}.} On the finite action set $A = \prod_i A_i$, correlated equilibria are the probability vectors $\sigma$ satisfying
\[
\sum_{a_{-i}}\sigma(a_i,a_{-i})\{u_i(a_i',a_{-i}) - u_i(a_i,a_{-i})\} \le 0
\quad\text{for all }i,a_i,a_i'.
\]
These finitely many linear inequalities intersect the probability simplex in a bounded polyhedron, hence a convex polytope.

For a mixed Nash equilibrium $x = (x_1,\dots,x_n)$, put $\sigma(a) = \prod_j x_j(a_j)$ and $U_i(a_i,x_{-i}) = \sum_{a_{-i}}\prod_{j\ne i}x_j(a_j)u_i(a_i,a_{-i})$. The constraint's left side is
\[
x_i(a_i)\{U_i(a_i',x_{-i}) - U_i(a_i,x_{-i})\} \le 0:
\]
it is zero off the support, and on the support $a_i$ is a best response. Thus every mixed Nash equilibrium induces a correlated equilibrium, and existence of a mixed Nash equilibrium makes the polytope nonempty.

\paragraph{Exercise~\ref{ex:chcdt:4}.} The net benefit is normal with mean $0.2\lambda-5000$ and variance $(0.1\lambda)^2+2000^2$, so
\[
\Prob(\lambda\Delta E-\Delta C>0) = \Phi\!\left(\frac{0.2\lambda-5000}{\sqrt{0.01\lambda^2+4\times10^6}}\right).
\]
The curve rises with $\lambda$: $0.006$ at $\lambda=0$, $0.09$ at $10{,}000$, $0.36$ at $20{,}000$, $0.50$ at $25{,}000$, $0.61$ at $30{,}000$, $0.82$ at $50{,}000$ and $0.93$ at $100{,}000$. It approaches $\Phi(2)=0.977$, the probability that $\Delta E>0$, as $\lambda\to\infty$. The ratio of posterior means is $5000/0.2=25{,}000$, so the ratio rule adopts the treatment for $\lambda>25{,}000$. For a normal net benefit this is also where the posterior mean turns positive and the probability passes one half, so the ratio rule is the Bayes decision under linear utility and agrees with the acceptability curve at level $1/2$. It says nothing about how sure the adoption is. At $\lambda=30{,}000$ it adopts although the net benefit is positive with probability only $0.61$. A probability of $0.8$ needs $\lambda\approx46{,}000$, $0.9$ needs $74{,}000$ and $0.95$ needs $147{,}000$, nearly six times the ratio, and no $\lambda$ gives certainty because $\Prob(\Delta E<0)=0.023$. The ratio of the posterior means is finite, but the posterior mean of $\Delta C/\Delta E$ does not exist, since $\Delta E$ has positive density at zero.

\paragraph{Exercise~\ref{ex:chcdt:5}.} A reduced form under rule $\rho_0$ has parameters $g(\theta,\rho_0)$ because private decisions depend on the expected rule. Changing to $\rho_1$ changes these parameters, so the reduced form fails modularity (Assumption~\ref{blp:ass:mod}): an intervention alters an equation it was meant to leave alone. This is the Lucas critique. In a modular structural model with invariant parameters $\theta$, prediction under the new rule is $\int P(Y\mid\theta,\rho_1)p(\theta\mid D,\rho_0)\,d\theta$ (Section~\ref{cdt:sec:lucas}).

For the New Keynesian model
\[
\begin{aligned}
x_t &= \E_t x_{t+1} - \sigma^{-1}(i_t - \E_t\pi_{t+1}),\\
\pi_t &= \beta\,\E_t\pi_{t+1} + \kappa x_t + u_t,\\
i_t &= \phi\,\pi_t,
\end{aligned}
\]
with $u_t = \nu u_{t-1} + e_t$ and $\theta = (\beta,\sigma,\kappa,\nu)$, modularity requires $\mathrm{do}(\phi:=\phi_1)$ to replace $i_t = \phi_0\pi_t$ by $i_t = \phi_1\pi_t$ while leaving the first two equations, $\theta$ and the shock law unchanged. Expectations are computed under the new rule. The assumption fails if price setting or the shock variance responds to the rule.

\paragraph{Exercise~\ref{ex:chgbc:1}.} Let $F$ be the distribution function of integrable $Z$ and $Q_Z(u) = \inf\{z : F(z) \ge u\}$. Right continuity gives $Q_Z(u) \le z$ exactly when $u \le F(z)$, so for $V$ uniform on $(0,1)$, $\Prob\{Q_Z(V) \le z\} = F(z)$. Thus $Q_Z(V)$ has the law of $Z$ and
\[
\E Z = \E Q_Z(V) = \int_0^1 Q_Z(u)\,du.
\]
For $\beta(u) = \eta u$, $\eta \in (0,1]$, substitute $u = \eta v$ in \eqref{gbc:eq:yaari}:
\[
\Vbeta = \int_0^1 Q_Z(\eta v)\,dv = \eta^{-1}\int_0^\eta Q_Z(u)\,du,
\]
the mean of the worst fraction $\eta$. Both integrals are finite since $\int_0^1 |Q_Z(u)|\,du = \E|Z| < \infty$.

\paragraph{Exercise~\ref{ex:chgbc:2}.} For $\eta < 0$, monotonicity of $\Phi$ gives $\beta(u) = \Phi(\Phi^{-1}(u) + \eta) < u$ (Table~\ref{gbc:tab:distortions}) on $(0,1)$, with $\beta(0) = 0$ and $\beta(1) = 1$ by limits. Thus $\beta(u) \le u$. For $Z \sim N(m,s^2)$,
\[
Q_Z\{\beta(u)\} = m + s\{\Phi^{-1}(u) + \eta\},
\qquad
\Vbeta = m + \eta s,
\]
because $\Phi^{-1}(V)$ is standard normal for $V$ uniform on $(0,1)$. The value is the mean lowered by $|\eta|$ standard deviations. For programme A, $N(1,0.35^2)$, and $\eta = -0.75$, this gives $1 - 0.75\times0.35 = 0.7375$, the $0.74$ of Figure~\ref{gbc:fig:distort}.

\paragraph{Exercise~\ref{ex:chgbc:4}.} Draw $200{,}000$ utilities per programme, A from $N(1,0.35^2)$ and B from $N(-0.9,0.5^2)$ with probability $0.12$ and from $N(1.35,0.35^2)$ otherwise, so the means are $1.00$ and $1.08$, and form both quantile functions on a grid of $4000$ levels. A loop over $\eta$ that evaluates $\Vbeta$ with $\beta(u)=\eta u$ for both programmes gives $\CVaR_\eta$ of each. At $\eta=0.25$ the values are $0.56$ and $-0.02$, as in the text, and at $\eta=1$ they are the means $1.00$ and $1.08$. The difference $\CVaR_\eta(B)-\CVaR_\eta(A)$ changes sign once, at $\eta^\ast\approx0.76$: the planner prefers A for $\eta<\eta^\ast$ and B for $\eta>\eta^\ast$.

The exact distributions give $\eta^\ast=0.752$, where both values are $0.85$. For a normal law $\CVaR_\eta=m-s\,\varphi(\Phi^{-1}(\eta))/\eta$, with $\varphi$ the standard normal density; for B, the partial mean $m\Phi(a)-s\varphi(a)$, $a=(q-m)/s$, of each component below the $\eta$ quantile $q$ of the mixture is weighted by the mixture probabilities and divided by $\eta$. The simulated value is $0.007$ higher because of Monte Carlo error, whose standard deviation at $200{,}000$ draws is about $0.006$, so the simulation fixes $\eta^\ast$ only to within about $0.01$.

\paragraph{Exercise~\ref{ex:chgbc:5}.} Use $\CVaR$ at level $\eta=\tfrac12$ (the mean of the worst half), no discounting and two periods. The second period has no decision, and $U_2$ is $0$ or $4$ with probability $\tfrac12$ each, independently of the first period, whatever the first action. In the first period action $a$ gives $U_1=1.5$ for sure and action $b$ gives $U_1=0$ or $4$ with probability $\tfrac12$ each.

Applied to total utility, $a$ has total $1.5$ or $5.5$ with probability $\tfrac12$ each, so its value is $1.5$. For $b$ the total is $0$, $4$ or $8$ with probabilities $\tfrac14,\tfrac12,\tfrac14$. The worst half is the atom at $0$ and a quarter of probability at $4$, so the value is $(0\cdot\tfrac14+4\cdot\tfrac14)/\tfrac12=2$. The whole horizon criterion chooses $b$.

The stagewise recursion of Section~\ref{gbc:sec:dynamic} gives $V_2 = \CVaR_{1/2}(U_2) = 0$ after every history and first period values $1.5$ for $a$ and $0$ for $b$, so it chooses $a$. For $b$, good outcomes in one period offset bad outcomes in the other under the total criterion: $\CVaR_{1/2}(U_1 + U_2) = 2$ although $\CVaR_{1/2}(U_1) + \CVaR_{1/2}(U_2) = 0$.

\section*{Chapter~\ref{ch:blp}: The Causal Effect of a Merger}

\paragraph{Exercise~\ref{ex:chapter_blp:1}.} Row $j$ of the first order conditions, for the firm $f$ that owns $j$, is $s_j+\sum_{k\in f}(\partial s_k/\partial p_j)m_k=0$. Under logit $\partial s_j/\partial p_j=-\alpha s_j(1-s_j)$ and $\partial s_k/\partial p_j=\alpha s_js_k$ for $k\neq j$, so
\[
s_j\Bigl\{1-\alpha(1-s_j)m_j+\alpha\sum_{k\in f,\,k\neq j}s_km_k\Bigr\}=s_j\Bigl\{1-\alpha m_j+\alpha\sum_{k\in f}s_km_k\Bigr\}=0 .
\]
Since $s_j>0$, $m_j=\alpha^{-1}+\sum_{k\in f}s_km_k$, which is the same for every $j\in f$. Call it $\bar m_f$. Then $\bar m_f=\alpha^{-1}+s_f\bar m_f$, so $\bar m_f=1/\{\alpha(1-s_f)\}$, with $1-s_f>0$ because $s_0>0$. Joint monopoly is one firm owning all $J$ products, with $s_f=1-s_0$, and every markup equals $1/(\alpha s_0)$. The diversion ratio is $-(\partial s_k/\partial p_j)/(\partial s_j/\partial p_j)=\alpha s_js_k/\{\alpha s_j(1-s_j)\}=s_k/(1-s_j)$; summing it over the partner's products, each with markup $1/\{\alpha(1-s_n)\}$, and dividing by $p_j$ gives $\mathrm{GUPPI}_j$. The price $p_j$ enters the condition only through the shares: utility is linear in price, so every share derivative is $\alpha$ times a polynomial in shares, and dividing by $s_j$ removes everything else. Given the shares, the markup is the same for all products of the firm whatever their prices. With a log price term the derivatives would carry $1/p_j$ and the condition would give $m_j=p_j/\alpha+\sum_{k\in f}s_km_k$, which does depend on own price.

\paragraph{Exercise~\ref{ex:chapter_blp:2}.} Row $j$ of $D_p^\top$ is $(\partial s_k/\partial p_j)_k$. With no outside good $\sum_ks_k=1$ at every price vector, so the row sums to $\partial(\sum_ks_k)/\partial p_j=0$, that is $D_p^\top\bm 1=0$. A common price shift $a\bm 1$ lowers every $V_{ij}$ by the same $\alpha_ia$ and, with $s_{i0}=0$, leaves every share unchanged, so also $D_p\bm 1=0$, that is $\bm 1^\top D_p^\top=0$. Joint monopoly has $H^M$ equal to the matrix of ones, so $H^M\odot D_p^\top=D_p^\top$, which is singular. Multiplying the first order conditions $\bm s+D_p^\top\bm m=0$ on the left by $\bm 1^\top$ gives $\bm 1^\top\bm s=1=0$, so no markup vector satisfies them. Economically, industry demand is perfectly inelastic because every consumer buys something. A monopolist who raises all prices by $a$ keeps all shares and gains $a\sum_js_j=a$ in profit, so profit has no maximum and markups are not defined.

In the proof of Lemma~\ref{blp:lem:dd} the strict inequality fails, because it used $s_{i0}>0$. With $s_{i0}=0$, $\sum_{k\neq j}\int\alpha_is_{ij}s_{ik}\,dF=\int\alpha_is_{ij}(1-s_{ij})\,dF$: the rows of $-D_p^\top$ sum to zero, dominance is weak, and $A=-(H^M\odot D_p^\top)$ is a singular M matrix, so $A^{-1}$ and $\bm m>0$ are lost. For $H^B$ with at least two firms the entrywise product deletes a strictly positive off diagonal entry from every row, strict dominance returns, and Bertrand Nash markups still exist.

\paragraph{Exercise~\ref{ex:chapter_blp:3}.} By Lemma~\ref{blp:lem:logit}, $\mathrm{GUPPI}_j=s_n/\{\alpha p_j(1-s_j)(1-s_n)\}=0.15/(3\times0.90\times0.85)=0.0654$. By Lemma~\ref{blp:lem:pressure}, $\Delta m_j/p_j=s_n/\{\alpha p_j(1-s_f)(1-s_f-s_n)\}=0.15/(3\times0.80\times0.65)=0.0962$. GUPPI understates the markup pressure by 3.1 points of the price, a share $1-0.0654/0.0962=0.32$ of it; in money $\Delta m_j=0.288$ against $p_j\mathrm{GUPPI}_j=0.196$. The ratio is
\[
\rho(s_n)=\frac{\Delta m_j}{p_j\mathrm{GUPPI}_j}=\frac{(1-s_j)(1-s_n)}{(1-s_f)(1-s_f-s_n)},\qquad \rho(0.15)=1.471 .
\]
Its logarithmic derivative is $-1/(1-s_n)+1/(1-s_f-s_n)=s_f/\{(1-s_n)(1-s_f-s_n)\}>0$, so the ratio rises with the partner's share: from $0.9/0.64=1.41$ as $s_n\to0$ to $1.88$ at $s_n=0.5$, and without bound as $s_n\to1-s_f$, when the merging firms take in the whole market. GUPPI keeps the partner's premerger markup, which stays finite, while the merged firm's markup at fixed prices is $1/\{\alpha(1-s_f-s_n)\}$.

\paragraph{Exercise~\ref{ex:chapter_blp:4}.} Write $\sigma=\sigma_p$, $\zeta_j=\alpha p_j$ with $\zeta_0=0$ for the outside good, $\bar\zeta=\sum_{k=0}^Js^0_k\zeta_k$ and $V_{ij}=\delta_j-\alpha e^{-\sigma\nu_i}p_j$. At $\sigma=0$, $\partial V_{ij}/\partial\sigma=\nu_i\zeta_j$ and $\partial^2V_{ij}/\partial\sigma^2=-\nu_i^2\zeta_j$. The first derivative is the one in the proof of Proposition~\ref{blp:prop:flat}, so $\partial s_{ij}/\partial\sigma=s^0_j\nu_i(\zeta_j-\bar\zeta)$ and $\partial s_j/\partial\sigma=s^0_j(\zeta_j-\bar\zeta)\E\nu_i=0$. In the second derivative the curvature of $V$ adds $s^0_j\{\partial^2V_{ij}/\partial\sigma^2-\sum_ks^0_k\partial^2V_{ik}/\partial\sigma^2\}=-s^0_j\nu_i^2(\zeta_j-\bar\zeta)$ to the terms of the proposition, and with $\E\nu_i^2=1$
\[
\frac{\partial^2s_j}{\partial\sigma_p^2}\Big|_{\sigma_p=0}=s^0_j\Bigl\{(\zeta_j-\bar\zeta)^2-\sum_{k=0}^Js^0_k(\zeta_k-\bar\zeta)^2-(\zeta_j-\bar\zeta)\Bigr\}.
\]
The extra term equals $\alpha\,\partial s^0_j/\partial\alpha$: the mean price coefficient $\E\alpha_i=\alpha e^{\sigma_p^2/2}\approx\alpha(1+\sigma_p^2/2)$ rises at order $\sigma_p^2$. It changes the curvature vector $\bm a$ and leaves the conclusion that the aggregate moments identify $\sigma_p^2$.

A micro moment weights choices by $\nu_i$. Let $\mathcal M$ be the set of minivans, $s^0_{\mathcal M}=\sum_{j\in\mathcal M}s^0_j$ and $\bar\zeta_{\mathcal M}=\sum_{j\in\mathcal M}s^0_j\zeta_j/s^0_{\mathcal M}$. From $s_{ij}\approx s^0_j\{1+\sigma_p\nu_i(\zeta_j-\bar\zeta)\}$, $\E[\nu_is_{ij}]=\sigma_ps^0_j(\zeta_j-\bar\zeta)+O(\sigma_p^2)$: the weight turns the vanishing $\E\nu_i$ into $\E\nu_i^2=1$. Hence
\[
\mu^{\mathrm{mic}}_1=\frac{\E[\nu_i\sum_{j\in\mathcal M}s_{ij}]}{\E[\sum_{j\in\mathcal M}s_{ij}]}=\sigma_p(\bar\zeta_{\mathcal M}-\bar\zeta)+O(\sigma_p^2),
\]
and $\mu^{\mathrm{mic}}_2=\sigma_p\sum_{j\in\mathcal M}(s^0_j/s^0_{\mathcal M})\,p_j(\zeta_j-\bar\zeta)+O(\sigma_p^2)$. The slope of $\mu^{\mathrm{mic}}_1$ is nonzero whenever the minivans are priced differently from the share weighted average (the outside good counting as price zero), which gives the Jacobian column for $\sigma_p$ that the aggregate moments lack at zero.

\paragraph{Exercise~\ref{ex:chapter_blp:5}.} With $\ell_2=49\ell_1$ the rule of Proposition~\ref{blp:prop:decision} challenges if and only if $\Prob(\Delta p/p>r\mid\mathcal D)>\ell_1/(\ell_1+\ell_2)=1/50=0.02$. In Table~\ref{blp:tab:cereal}, with $r=5\%$, the GBC probability is 1.00 for the merger as proposed and 0.01 with the divestiture of C2. The first exceeds 0.02, so the agency challenges the merger; the second is below it, so the agency accepts the remedy. Of the \numcerngbcdraws{} posterior draws behind the table, \numcerprgbcnabove{} exceed 5 percent, so the second entry is 0.007 before rounding. A normal approximation to the printed mean and standard deviation, 3.0 and 1.3 percent, would give 0.06 and the opposite action; it is wrong here, because a few extreme draws inflate the standard deviation while 90 percent of the draws lie between \numcerdprgbcqlo{} and \numcerdprgbcqhi{} percent. The BLP GMM entries, 1.00 and 0.00, give the same actions. 

\paragraph{Exercise~\ref{ex:chapter_blp:6}.} By Proposition~\ref{blp:prop:bma} the posterior mean is $\bar\mu=0.7\times10+0.3\times0=7$ points. The within conduct variance is $0.7\times1.5^2+0.3\times0.5^2=1.65$ and the between conduct variance is $0.7\times(10-7)^2+0.3\times(0-7)^2=21$, so $\mathrm{Var}(e\mid\mathcal D)=22.65$ and the standard deviation is $4.76$ points. A procedure that maintains Bertrand Nash has posterior expected squared error $22.65+(10-7)^2=31.65$, a root of $5.63$ points; directly, $0.7\times1.5^2+0.3\times(0.5^2+10^2)=31.65$. It reports the standard deviation $1.5$, which understates its own root mean squared error by a factor $3.75$ and the posterior variance by a factor of ten ($22.65/2.25$). The between conduct term supplies 93\% of the posterior variance.

\paragraph{Exercise~\ref{ex:chapter_blp:7}.} Let $H^{\mathrm{post}}$ be $H^B$ with the products of A and B given a common owner. The effect on the outsider is the C1 entry of $\Delta\bm p_t$,
\[
\Delta p_{C1,t}=p_{C1,t}(H^{\mathrm{post}})-p_{C1,t}(H^B),
\]
and the row of $H$ for C1 is the same in both terms. The effect is nonzero because the equilibrium in \eqref{blp:eq:scm} determines all prices jointly: the potential price of C1 depends on the conduct of other products, so the stable unit treatment value assumption fails. A and B raise prices, C1's share rises, and C1's best response is a higher markup.

Put $N = \nabla\bm m^{\mathrm{post}}(\bm p)$. Proposition~\ref{blp:prop:firstorder} gives $\Delta\bm p_t = (I-N)^{-1}\Delta\bm m + O(\lVert\Delta\bm m\rVert^2)$. The block of firm C is unchanged by the merger, so $\bm s_C + (D_p^\top)_{CC}\bm m_C = 0$ under both ownership matrices and $\Delta\bm m_C = 0$ at fixed prices. Yet C's shares depend on A and B's prices: under logit, $m_{C1} = 1/\{\alpha(1-s_C)\}$ and $N_{C1,k} = s_ks_C/(1-s_C)^2 > 0$ for $k$ in A or B. Thus the first order outsider effect is
\[
\sum_{k\in A\cup B}[(I-N)^{-1}]_{C1,k}\Delta m_k,
\]
the merging parties' markup pressure transmitted through best responses.

\paragraph{Exercise~\ref{ex:chapter_blp:8}.} Take the structural posterior of the realized effect as $N(9,1.5^2)$. The retrospective estimate adds noise of standard deviation 2, so its posterior predictive distribution is $N(9,1.5^2+2^2)=N(9,2.5^2)$. The observed 4\% is $(4-9)/2.5=-2$ predictive standard deviations from the mean. For a standard normal $Z$, $\Prob(Z\le-2)=0.023$ is the probability of an estimate at least that low and $\Prob(|Z|\ge2)=0.046$ the probability of one at least that far from 9 in either direction. The prediction is not refuted, but it misses by twice its predictive spread and toward a smaller effect.

Demand parameter uncertainty is already inside the 1.5, so the first suspects are the assumptions that can lower the effect with no parameter wrong. Modularity is one: efficiencies that cut marginal cost offset the markup rise. Conduct is the other: if premerger prices already carried coordination, a merger adds little markup (under $H^M$ the effect is zero, against the Bertrand prediction). For the next case, give $H$ prior weight on collusive conduct, add a cost change $\Delta c$ with a prior centered at zero, and treat each retrospective estimate $\hat e$ with standard error $s$ as data, $\hat e\mid\theta_2,H,\Delta c\sim N\{e(\theta_2,H,\Delta c),s^2\}$. The posterior over $(H,\Delta c)$ from this case is then the prior for the next.

\paragraph{Exercise~\ref{ex:chapter_blp:9}.} Let $O_{jk} = 1$ for products with the same owner and zero otherwise. Then $H(\kappa) = O + \kappa(\mathbf1\mathbf1^\top - O)$ gives $H(0) = H^B$ and $H(1) = H^M$. Differentiating $\sum_k H_{jk}(p_k-c_k)s_k$ in $p_j$ gives
\[
s_j + \sum_k H_{jk}\frac{\partial s_k}{\partial p_j}(p_k-c_k) = 0,
\qquad
\bm m = -(H(\kappa)\odot D_p^\top)^{-1}\bm s
\]
where the matrix is invertible. The endpoints internalize the owner's profits and all industry profits, respectively. Perfect competition, $\bm m = 0$, lies outside this family. A merger replaces $O$ by post merger ownership: merged products get weight one and rivals retain weight $\kappa$. Larger $\kappa$ leaves less to internalize; at $\kappa = 1$ the merger changes nothing.

To learn $\kappa$, add a branch to the ownership function of \texttt{blp\_core.py} that returns $H(\kappa)$, and draw $\kappa$ uniformly on $(0, 1)$ with the other parameters from the same prior. The equilibrium solver and the summaries need no change. The markup regressions among those summaries are what separates $H^B$ from $H^M$, since their loadings on $1/s_0 - 1/(1 - s_f)$ differ, and they carry the information about $\kappa$ in between. Because $\kappa$ is now one more coordinate of $\theta$ in Algorithm~\ref{blp:alg:gbc}, the discrete conduct classifier is not needed, and the targets of the noise outsourcing regression become the twelve parameters and $\log\{\kappa/(1 - \kappa)\}$, whose out of fold residuals are added on the logit scale so that draws stay in $(0, 1)$. We drew 30,000 industries from the prior on the minivan design with $T = 40$. Rejection hardly tilts the prior, whose quantiles at 5, 50 and 95 percent are 0.05, 0.51 and 0.95. We trained on 29,600 with three folds and the default hidden layers of 256, 256 and 128 units, held out 400, kept the 1,480 nearest neighbours, and used seed 910,000 plus the draw index.

On the minivan data of Section~\ref{blp:sec:app}, which are simulated under Bertrand Nash so that the truth is $\kappa = 0$, the posterior of $\kappa$ has mean 0.12, median 0.11 and 95\% interval 0.009 to 0.33, with $\Prob(\kappa < 0.1) = 0.44$ and $\Prob(\kappa < 0.25) = 0.93$. The demand parameters remain in place: $\alpha$ has mean 0.87 with interval 0.70 to 1.04 (truth 1), $\sigma_p$ 0.40 with 0.30 to 0.49 (truth 0.4), and $\sigma_x$ 0.71 with 0.44 to 0.93 (truth 0.7). The data rule out weights on rivals' profits much above one third but cannot tell $\kappa = 0$ from $\kappa = 0.05$, and a continuous prior cannot put mass on the point zero, so the mean of 0.12 is partly the boundary at work. A point mass at zero in the prior would answer the question of whether Bertrand Nash holds exactly.

The method is calibrated on this design. On 400 industries held out from the table, 97 percent of the 95\% intervals contain the true $\kappa$, with root mean squared error 0.09 against a prior standard deviation of 0.29, mean interval width 0.46 and correlation 0.95 between posterior mean and truth. Coverage is 94 percent for the 139 cases with true $\kappa$ below 0.33, all 136 cases between 0.33 and 0.67, and 97 percent of the 125 above 0.67. New data at the same parameters (seed 2027) give posterior mean 0.51 with interval 0.18 to 0.75 when $\kappa = 0.5$, and 0.97 with interval 0.93 to 1.00 when $\kappa = 1$. This also addresses the failure in Section~\ref{blp:sec:misspec}, where partial internalization lies outside the training set and the classifier has to choose between the two ends: once $\kappa$ is a coordinate, a value of one half is estimated rather than forced to an end.

\paragraph{Exercise~\ref{ex:chapter_blp:10}.} Let $\bar{\bm\psi}(\theta_2,\beta)=T^{-1}\sum_t\psi(W_t;\theta_2,\beta)=\bar{\bm\psi}_0(\theta_2)-\mathcal C\beta$, where $\bar{\bm\psi}_0(\theta_2)=T^{-1}\sum_tZ_t^\top\delta_t(\theta_2)$ and $\mathcal C=T^{-1}\sum_tZ_t^\top\bm x_t^\top$, and let $\mathcal W$ be the weighting matrix. The criterion $\bar{\bm\psi}^\top\mathcal W\bar{\bm\psi}$ has $\beta$ gradient $-2\mathcal C^\top\mathcal W\bar{\bm\psi}$, which vanishes at the profiled value $\hat\beta(\theta_2)=(\mathcal C^\top\mathcal W\mathcal C)^{-1}\mathcal C^\top\mathcal W\bar{\bm\psi}_0(\theta_2)$ by the definition of $\hat\beta$. The concentrated moment is $\bar{\bm\psi}\{\theta_2,\hat\beta(\theta_2)\}=M\bar{\bm\psi}_0(\theta_2)$ with $M=I-\mathcal C(\mathcal C^\top\mathcal W\mathcal C)^{-1}\mathcal C^\top\mathcal W$. Since $M\mathcal C=0$, $M\bar{\bm\psi}(\theta_2,\beta)=M\bar{\bm\psi}_0(\theta_2)$ for every $\beta$, so its derivative with respect to $\beta$ is zero at the profiled value and everywhere. With $\dot{\bm\psi}_0=\partial\bar{\bm\psi}_0/\partial\theta_2^\top$ and $M^\top\mathcal WM=\mathcal WM$, the first order condition for $\theta_2$ of the concentrated criterion is the estimating equation
\[
\dot{\bm\psi}_0(\theta_2)^\top\mathcal W\,M\,\bar{\bm\psi}_0(\theta_2)=0 .
\]
The unprojected score $\dot{\bm\psi}_0^\top\mathcal W\bar{\bm\psi}$ has $\beta$ derivative $-\dot{\bm\psi}_0^\top\mathcal W\mathcal C\neq0$; the factor $M$ makes the equation orthogonal to $\beta$.

Orthogonality removes the first order effect of errors in $\beta$ on $\hat\theta_2$. Proposition~\ref{blp:prop:pushforward} shows what it leaves: even if $\E\hat\theta_2=\theta_2$, $\E g(\hat\theta_2)=g(\theta_2)+\tfrac12\mathrm{tr}\{\nabla^2g(\theta_2)\mathcal V\}+o(\lVert\mathcal V\rVert)$, and the curvature term comes from the nonlinear map $g$, which is a fixed point, whatever the nuisance. For the logit markup $1/\{\alpha(1-s)\}$, convex in $\alpha$, the bias is $\mathrm{Var}(\hat\alpha)/\{\alpha^3(1-s)\}$, a fraction $\mathrm{Var}(\hat\alpha)/\alpha^2$ of the markup: 4\% when $\hat\alpha$ has a coefficient of variation of 20\%. The pushforward $g_\#\Pi(\cdot\mid\mathcal D)$ averages $g$ over the posterior and has no such term.
